% AFDX 端系统设计文档
% 编译：在本目录执行 tectonic AFDX端系统FPGA设计文档.tex（或 xelatex 两遍）
\documentclass[a4paper,zihao=-4,fontset=none,oneside,openany]{ctexrep}

% ---------- 字体 ----------
\setmainfont{Times New Roman}
\setsansfont{Helvetica}
\setmonofont{Menlo}[Scale=0.82]
% Kaiti SC 是按需下载字体，部分沙箱环境读不到，读不到时退回宋体
\IfFontExistsTF{Kaiti SC}{\def\kaifont{Kaiti SC}}{\def\kaifont{Songti SC}}
\setCJKmainfont{Songti SC}[BoldFont=Heiti SC, ItalicFont=\kaifont]
\setCJKsansfont{Heiti SC}
\setCJKmonofont{Heiti SC}
\setCJKfamilyfont{zhsong}{Songti SC}[BoldFont=Heiti SC]
\setCJKfamilyfont{zhhei}{Heiti SC}
\setCJKfamilyfont{zhkai}{\kaifont}
\providecommand{\songti}{\CJKfamily{zhsong}}
\providecommand{\heiti}{\CJKfamily{zhhei}}
\providecommand{\kaishu}{\CJKfamily{zhkai}}

% ---------- 版式 ----------
\usepackage[a4paper,top=2.4cm,bottom=2.3cm,left=2.6cm,right=2.4cm,headheight=15pt]{geometry}
\usepackage{graphicx}
\usepackage{booktabs}
\usepackage{tabularx}
\usepackage{longtable}
\usepackage{array}
\usepackage{multirow}
\usepackage{etoolbox}
\usepackage{needspace}
\usepackage[section]{placeins}
\usepackage{amsmath}
\usepackage{enumitem}
\usepackage{xcolor}
\usepackage{tikz}
\usepackage{tcolorbox}
\tcbuselibrary{listings}
\usepackage{fancyhdr}
\usepackage[font=small,labelsep=quad,skip=6pt]{caption}
\usepackage[hidelinks,bookmarksnumbered,bookmarksdepth=2]{hyperref}

\linespread{1.2}
\setlength{\parskip}{3pt}
\emergencystretch=3em
\raggedbottom
\clubpenalty=10000
\widowpenalty=10000
\displaywidowpenalty=10000
\setlist{topsep=4pt plus 1pt minus 1pt,itemsep=0pt,parsep=0pt,leftmargin=2em}
\AtBeginDocument{%
  \setlength{\abovedisplayskip}{9pt plus 2pt minus 2pt}%
  \setlength{\belowdisplayskip}{9pt plus 2pt minus 2pt}%
  \setlength{\abovedisplayshortskip}{3pt plus 2pt}%
  \setlength{\belowdisplayshortskip}{6pt plus 2pt minus 2pt}}
% 表格统一字号与行距；正文保留小四字号，插图按原尺寸插入。
\newcommand{\tablestyle}{\small\linespread{1.08}\selectfont
  \renewcommand{\arraystretch}{1.16}}
\AtBeginEnvironment{table}{\tablestyle}
\BeforeBeginEnvironment{longtable}{\FloatBarrier\begingroup\tablestyle}
\AfterEndEnvironment{longtable}{\endgroup}
\setlength{\LTpre}{8pt}
\setlength{\LTpost}{8pt}
\setlength{\LTcapwidth}{\textwidth}
\captionsetup{font={small,stretch=1.08},skip=5pt,format=hang,justification=raggedright,
  singlelinecheck=true}
\captionsetup[table]{position=top}
% 图表贴近引文，在节边界和长表前清空浮动队列。
\renewcommand{\topfraction}{0.9}
\renewcommand{\bottomfraction}{0.8}
\renewcommand{\textfraction}{0.1}
\renewcommand{\floatpagefraction}{0.8}
\setcounter{topnumber}{2}
\setcounter{bottomnumber}{1}
\setcounter{totalnumber}{3}
\setlength{\floatsep}{8pt plus 2pt minus 2pt}
\setlength{\textfloatsep}{10pt plus 2pt minus 2pt}
\setlength{\intextsep}{8pt plus 2pt minus 2pt}
\makeatletter
\def\fps@table{!htbp}
\setlength{\@fptop}{0pt}
\setlength{\@fpsep}{12pt}
\setlength{\@fpbot}{0pt plus 1fil}
\appto\@floatboxreset{\tablestyle}
\makeatother

\ctexset{
  chapter = {format=\zihao{3}\heiti\raggedright, name={第,章}, number=\arabic{chapter},
             beforeskip=0pt, afterskip=12pt},
  section = {format=\zihao{4}\heiti, beforeskip=12pt, afterskip=6pt,
             break=\Needspace{3\baselineskip}},
  subsection = {format=\zihao{-4}\heiti, beforeskip=9pt, afterskip=5pt,
                break=\Needspace{3\baselineskip}},
}
\renewcommand{\thefigure}{\thechapter-\arabic{figure}}
\renewcommand{\thetable}{\thechapter-\arabic{table}}
\renewcommand{\theequation}{\thechapter-\arabic{equation}}

\pagestyle{fancy}
\fancyhf{}
\fancyhead[L]{\small AFDX 端系统 FPGA 设计文档}
\fancyhead[R]{\small\leftmark}
\fancyfoot[C]{\small\thepage}
\renewcommand{\headrulewidth}{0.4pt}
\fancypagestyle{plain}{\fancyhf{}\fancyfoot[C]{\small\thepage}\renewcommand{\headrulewidth}{0pt}}

% ---------- 命令 ----------
\newfontfamily\runcommandfont{Menlo}[Scale=1.0]
\definecolor{commandback}{HTML}{F5F6F7}
\definecolor{commandframe}{HTML}{D9DDE2}
\newtcblisting{runcommand}{
  listing only,
  colback=commandback, colframe=commandframe,
  boxrule=0.4pt, arc=0pt, outer arc=0pt,
  boxsep=0pt, left=7pt, right=7pt, top=6pt, bottom=6pt,
  before skip=7pt, after skip=7pt,
  listing options={
    basicstyle=\linespread{1}\runcommandfont\fontsize{9pt}{12pt}\selectfont,
    columns=fullflexible, keepspaces=true, showstringspaces=false,
    breaklines=true, breakatwhitespace=false, upquote=true,
    literate={-}{{\char45}}1,
    aboveskip=0pt, belowskip=0pt}}
% 代码名不连字，但允许在下划线和斜杠后断行
\ExplSyntaxOn
\NewDocumentCommand{\code}{m}{\texttt{\hyphenchar\font=-1\relax
  \str_map_inline:nn {#1} { ##1 \str_if_in:nnT {_/} {##1} { \penalty 200 \relax } }}}
\ExplSyntaxOff
\newcolumntype{L}{>{\raggedright\arraybackslash}X}
\newcolumntype{C}{>{\centering\arraybackslash}X}
\newcolumntype{R}{>{\raggedleft\arraybackslash}X}
\newcommand{\longtablehead}[2]{%
  \toprule #2 \\ \midrule \endfirsthead
  \caption[]{（续）}\\
  \toprule #2 \\ \midrule \endhead
  \midrule \multicolumn{#1}{r}{接下页}\\ \endfoot
  \bottomrule \endlastfoot}
% 插图：#1 文件名（fig/pdf/ 下），#2 宽度系数，#3 图题，#4 标签
\newcommand{\fig}[4]{%
  \begin{figure}[!htbp]\centering
  \includegraphics[width=#2\linewidth]{fig/pdf/#1.pdf}
  \caption{#3}\label{#4}
  \end{figure}}
% 按印刷尺寸绘制的图：原尺寸插入，不缩放，保证全文字号一致
\NewDocumentCommand{\figN}{o m m m}{%
  \begin{figure}[!htbp]\centering
  \includegraphics{fig/pdf/#2.pdf}
  \IfValueTF{#1}{\caption[#1]{#3}}{\caption{#3}}\label{#4}
  \end{figure}}
% 全文引用的实测数字。最终版本重跑回归和综合后只改这里。
% 出处：docs/FACTS.md、docs/FACTS_DETAIL.md、docs/evidence/。
\newcommand{\CodeVer}{12cb849}
% 性能（tb_sys_perf，77 条 VL，1471 字节，BAG = 1 ms，双向，稳定测量 21 ms）
\newcommand{\Thru}{906.136}
\newcommand{\ThruEightyOne}{953.208}
\newcommand{\ThruTen}{117.680}
\newcommand{\Lat}{24.408}
\newcommand{\Jit}{0}
\newcommand{\Ttx}{0.248}
\newcommand{\Trx}{11.952}
% 突发（tb_gbe_burst）
\newcommand{\BurstSevSeven}{959.272}
\newcommand{\BurstEightOne}{1008.488}
% RTC（tb_sys_rtc）
\newcommand{\RtcIdle}{0}
\newcommand{\RtcLoad}{10788}
% 综合（eLinx 界面 Run Synthesis，EQ6HL130 / CSG484\_H，未布线估算）
\newcommand{\Lut}{32847}
\newcommand{\FF}{13869}
\newcommand{\Clb}{4620}
\newcommand{\Mem}{340}
\newcommand{\Pins}{259}
\newcommand{\LutPct}{24.12}
\newcommand{\FFPct}{10.19}
\newcommand{\ClbPct}{33.93}
\newcommand{\MemPct}{28.33}
\newcommand{\PinPct}{76.40}
\newcommand{\Setup}{2.522}
\newcommand{\Fmax}{182.548}
\newcommand{\Hold}{0.068}
% 命令行综合
\newcommand{\CliCells}{29963}
\newcommand{\CliSetup}{2.749}
\newcommand{\CliFmax}{190.440}
% 验证
\newcommand{\RegCount}{38}
\newcommand{\RegChecks}{1183249131}
\newcommand{\RegIndep}{2305780}
\newcommand{\SvcChecks}{1123}
\newcommand{\CapChecks}{49924}
\newcommand{\SnmpChecks}{699}
\newcommand{\ElinxChecks}{9238}


\begin{document}

% ---------- 封面 ----------
\hypersetup{pageanchor=false}
\begin{titlepage}
\centering
\vspace*{2.2cm}
{\zihao{4}\heiti 全国大学生嵌入式芯片与系统设计竞赛\par}
\vspace{4pt}
{\zihao{5} FPGA 创新设计赛道 \quad 中科亿海微 \quad 选题一：基于以太网 MAC 总线实现 AFDX\par}
\vspace{3cm}
{\zihao{1}\heiti AFDX 端系统\par}
\vspace{10pt}
{\zihao{2}\heiti FPGA 设计文档\par}
\vspace{18pt}
{\zihao{4} 方案与具体设计\par}
\vfill
\begin{tabular}{rl}
\heiti 队伍名称： & \underline{\makebox[7.8cm][l]{[待填写]}} \\[6pt]
\heiti 参赛成员： & \underline{\makebox[7.8cm][l]{[待填写]}} \\[6pt]
\heiti 指导教师： & \underline{\makebox[7.8cm][l]{[待填写]}} \\[6pt]
\heiti 目标器件： & \underline{\makebox[7.8cm][l]{EQ6HL130，CSG484\_H 封装}} \\[6pt]
\heiti 开发环境： & \underline{\makebox[7.8cm][l]{eLinx 3.0.7，ModelSim SE-64 2020.4}} \\[6pt]
\end{tabular}
\vspace{2cm}

{\zihao{4} 2026 年 10 月\par}
\end{titlepage}

\pagenumbering{Roman}
\hypersetup{pageanchor=true}

\chapter*{摘要}
\addcontentsline{toc}{chapter}{摘要}

本设计针对全国大学生嵌入式芯片与系统设计竞赛“基于以太网 MAC 总线实现 AFDX”赛题要求，基于中科亿海微 EQ6HL130 FPGA 研制了高完整性、双网冗余的 \mbox{ARINC} 664 Part 7 航电端系统。端系统逻辑由单一 125~MHz（8~ns 周期）主时钟驱动，以纯 Verilog HDL 实现了 128 条发送虚拟链路（VL）与 4096 条接收虚拟链路的线速硬件处理，并集成 \mbox{ARINC} 615A 软件数据加载（Data Loading）、SNMPv1/v2c 网络管理与高精度 SNTP/RTC 时钟同步服务。

面对大规模虚拟链路状态维护与 125~MHz 物理时序收敛的双重挑战，系统确立了“片上 BRAM 表格化管理与单主控状态机”的核心架构，全芯片划分为 24 个 RTL 模块：接收映射表采用 32-Bank 分块存储与双级读回流水，切断大扇出长连线延迟；管理面协议引擎（ASN.1/BER 与 TFTP）引入 2 分频局部时钟使能（$\text{CE\_DIV}=2$），将复杂状态转移约束在 16~ns 多周期路径中；报文线速收发、RAM 访问与配置总线维持 125~MHz 单周期全速运行，兼顾复杂协议处理与周期级实时响应。

经系统级全功能仿真验证与官方工具链物理签核，设计达成以下核心指标：
\begin{enumerate}[leftmargin=2em,itemsep=1pt]
  \item \textbf{吞吐率与确定性时延：}定长满载（77 VL）有效载荷吞吐率达 \Thru~Mbps（81 VL 达 \ThruEightyOne~Mbps），多 VL 混合变长高负载（\code{var_high}）两方向单向吞吐率均达 \ThruVarHigh~Mbps（双向合计 1809.504~Mbps）；单帧端到端硬件流水时延典型值为 24.424~µs，满载稳态最大时延为 \Lat~µs（其中发送段 $L_\text{tx}$ 典型为 \Ttx~µs、至首前导码 \TtxSof~µs，接收段 $L_\text{rx}$ 首字节为 \Trx~µs、末字节为 \TrxLast~µs），显著优于赛题 1--2~ms 的指标要求。
  \item \textbf{传输抖动控制：}稳态发送间隔抖动实测为 \Jit~µs；严格按官方 1022 答疑 R97 口径，定长满载同 VL 端到端时延极差抖动仅为 \JitSameVl~µs（1 拍，8~ns），多 VL 变长循环工况下为 \JitVarCycle~µs，实测极差源于串行线路上前序在途帧排队深度的动态变化。
  \item \textbf{硬件资源与物理时序收敛：}逻辑综合消耗查找表（LUT）\Lut{} 个（占比 \LutPct\%）、触发器（FF）\FF{} 个（占比 \FFPct\%）、嵌入式存储块 \Mem{} 个 EBR（占比 \MemPct\%）；在 eLinx 官方工具链（\code{S2d + PRES1.50}）下实现 125~MHz 布局布线后静态时序完全收敛，单周期路径最差建立裕量为 $+\RoutedSetup$~ns、保持裕量为 $+\RoutedHold$~ns，多周期路径建立裕量为 $+\MulticycleSetup$~ns，全设计 45,745 个时序端点 0 违例，最高等效工作频率达 \RoutedFmax~MHz。
  \item \textbf{全流程验证覆盖：}构建包含 \RegCount{} 项自动化回归用例的验证平台，累计通过 \RegChecks{} 次硬件断言自检与 \RegIndep{} 项独立参考模型（Golden Model）逐字节比对；FSM 状态覆盖率达 \CovFsmState\%、语句覆盖率达 \CovStmt\%、分支覆盖率达 \CovBranch\%；综合后 VQM 门级网表后仿真完全通过五功能联合测试。
\end{enumerate}

\vspace{0.8em}
\noindent\textbf{关键词：}AFDX 端系统；ARINC 664；FPGA；高完整性通信；静态时序收敛；布局布线

\clearpage

% 印刷目录列章、节；PDF 书签保留小节，便于定位具体机制。
\setcounter{tocdepth}{1}
\begingroup
\small\linespread{1.08}\selectfont
\setlength{\parskip}{0pt}
\tableofcontents
\endgroup
\clearpage
\pagenumbering{arabic}

\chapter{概述}\label{ch:intro}

\section{设计目标}

AFDX 端系统是航电网络中主机与交换网络之间的接口：它要按每条虚拟链路（VL）的带宽分配间隔（BAG）把应用消息装成以太网帧，同时从 A、B 两路网络发出；接收时逐帧检查完整性，把两路冗余副本合并为唯一一份交付给应用。赛题要求端系统支持 ARINC 664 Part 7 通信、高完整性通信、SNMP 网络管理、ARINC 615A 数据加卸载和 RTC 时钟同步五项功能，并给出传输速率 100~Mbps、端到端时延 1--2~ms、抖动几十微秒、128 条发送 VL 和 4096 条接收 VL 的技术指标。

在 FPGA 上实现这些要求，核心矛盾在于逐链路状态的规模与逐字节线速通路之间：4096 条接收 VL 和 128 条发送 VL 的配置、序号与定时状态若放进触发器，资源随链路数线性膨胀；若放进块 RAM，每帧的查表、比较和回写又会落到关键路径上，难以在每拍一字节的数据通路下保持确定的发送时刻。

本设计实现了完整的 AFDX 端系统 \code{AFDX_top}，以单一 125~MHz 时钟驱动两路八位 GMII，全部由 Verilog RTL 自主实现，经 eLinx~3.0.7 综合、ModelSim 仿真验证。设计有三项关键技术：
\begin{enumerate}
  \item \textbf{表加单主状态机。}为在线速下同时支撑 128 条发送 VL 和 4096 条接收 VL，所有逐 VL 配置和状态存放在块 RAM 中，每张表只由一个状态机按固定流水读写，数据通路每拍处理一个字节。这一结构在 EQ6HL130 上只占用 \LutPct\% 的 LUT 和 \MemPct\% 的存储块。
  \item \textbf{提前选、到点发的发送调度。}为使发送时刻严格对齐 BAG，发送引擎在到期前 512 拍选出候选 VL 并预先组好帧头，到点即发。只要没有其他 VL 的帧争用线路，帧的第一个前导字节就恰在允许时刻出现，调度本身不引入抖动。
  \item \textbf{自研双网 MAC。}一个状态机驱动 A、B 两套 CRC 与输出寄存器，两网天然同拍。不经用户侧 FIFO，空闲发送延迟为 0 拍；与 eLinx 生成的三速以太网 IP 配置相比，单口 LUT 用量为 556 对 1400。
\end{enumerate}
在此基础上，端系统每网每方向的线速载荷容量为 956.44~Mbps，是赛题门限的 9.5 倍；线路空闲时应用到应用的时延上界为 26.360~µs，比 1~ms 要求低近两个数量级（推导见第~\ref{sec:budget}~节）。

\section{主要成果}

表~\ref{tab:func} 列出赛题五项功能的实现要点，表~\ref{tab:perf} 给出技术指标的设计值。设计值均由结构参数推导，推导过程见第~\ref{sec:budget}~节；仿真实测结果见测试报告。

\begin{table}\centering\small
\caption{赛题功能与实现要点}\label{tab:func}
\begin{tabularx}{\linewidth}{>{\raggedright\arraybackslash}p{3.4cm}L>{\raggedright\arraybackslash}p{2.0cm}}
\toprule
赛题功能 & 实现要点 & 详见 \\
\midrule
ARINC 664 P7 通信 & 128 条发送 VL 按 BAG 调度，UDP/IP 封装，A、B 双网同时发送；4096 条接收 VL 按 VLID 直接查表 & 第~\ref{ch:tx}、\ref{ch:rx}~章 \\
高完整性通信 & FCS、帧长、IP/UDP 头逐帧检查；每网独立的序号完整性检查；按处理顺序交付首个有效副本的冗余管理 & 第~\ref{ch:rx}~章 \\
SNMP 网络管理 & v1/v2c 的 Get、GetNext、Set 和 v2 Trap，36 个管理对象，BER 逐字节编解码 & 第~\ref{sec:snmp}~节 \\
615A 数据加卸载 & 基于 TFTP 的上传（LUI、LUS、LUR、LUH、数据文件）与介质定义卸载，8~KiB 加载区，CRC-16 校验 & 第~\ref{sec:a615}~节 \\
RTC 时钟同步 & SNTPv4 主从对时，32 位秒加纳秒，8~ns 分辨率，无乘除法器的时间戳换算 & 第~\ref{sec:rtc}~节 \\
\bottomrule
\end{tabularx}
\end{table}

\begin{table}\centering\small
\caption{技术指标与设计值}\label{tab:perf}
\begin{tabularx}{\linewidth}{>{\raggedright\arraybackslash}p{1.9cm}>{\raggedright\arraybackslash}p{2.8cm}>{\raggedright\arraybackslash}p{4.2cm}L}
\toprule
指标 & 选题指南要求 & 设计值 & 依据 \\
\midrule
传输速率 & 100~Mbps & 每网每方向载荷容量 956.44~Mbps & 每拍一字节；1471 字节满载帧连同前导、FCS 和帧间隔占 1538 拍（第~\ref{sec:budget}~节） \\
端到端时延 & $<$1--2~ms & 线路空闲时 $\le$26.360~µs & 发送 $\le$100 拍、线路 1525 拍、接收首字节 $\le$200 拍、逐字节交付 $n-1$ 拍；线路争用另加式~(\ref{eq:jitter}) 的等待 \\
发送抖动 & $\le$ 几十微秒 & 无线路争用时为 0；争用时不超过前面帧的占线、提前授权等待与扫描开销之和，满足配置约束时不超过 50~µs & 提前 512 拍选出、到点起始；BAG 为 $2^k$~ms 时争用只造成一次相位后移（第~\ref{sec:jitter}~节） \\
帧格式 & VLID、SN 1--255，载荷 $\le$1471~B，帧长 $\le$1518~B & 全部按规定生成和检查 & 帧长 $47+\max(n,17)$，64--1518 字节（第~\ref{sec:frame}~节） \\
VL 容量 & 128 发送 / 4096 接收 & 128 / 4096，接收 VLID 任意 & 逐 VL 状态存块 RAM；65536 项 VLID 映射表 \\
\bottomrule
\end{tabularx}
\end{table}

综合后估算方面，占用 LUT \Lut、触发器 \FF、CLB \Clb、存储块 \Mem，8~ns 约束下建立裕量 +\Setup~ns，估算最高频率 \Fmax~MHz（第~\ref{ch:timing}~章）。

\section{文档结构}

第~\ref{ch:req}~章：赛题要求和官方答疑到可验证设计规格的分解；第~\ref{ch:arch}~章：总体架构与关键设计决策；第~\ref{ch:if}~章：顶层接口、线路帧格式与模块间接口；第~\ref{ch:tx}、\ref{ch:rx}~章：收发通路；第~\ref{ch:svc}~章：RTC、SNMP、615A 三项服务；第~\ref{ch:cfg}~章：配置表、异常处理、主机编程流程与配置约束；第~\ref{ch:timing}~章：时序设计规则与资源；第~\ref{ch:verif}~章：验证观察手段；第~\ref{ch:limit}~章：设计边界与后续工作。附录：源文件、参数、服务寄存器、内部接口信号表与术语和参考文献。

\chapter{需求分析与规格分解}\label{ch:req}

本章把选题指南的功能、指标和接口要求，连同官方答疑的口径，分解为编号的设计规格，并给出规格到设计章节和测试用例的追溯关系。下文以 R$n$ 表示官方《FPGA 创新大赛问题汇总》（1010 版）Sheet1 的第 $n$ 行。

\section{赛题需求}

选题指南（正文第 2--4 页）对 AFDX 端系统提出三类要求：
\begin{itemize}
  \item \textbf{功能：}ARINC 664 Part 7 通信、高完整性数据通信、SNMP 网络管理、615A 数据加卸载、RTC 时钟同步。
  \item \textbf{技术指标：}传输速率 100~Mbps；帧格式基于 IEEE 802.3 加 ARINC 664 扩展，含 VLID 和 1--255 的序列号，有效载荷不超过 1471 字节，总帧长不超过 1518 字节；端到端时延小于 1--2~ms，抖动不超过几十微秒；最多 128 条发送 VL、4096 条接收 VL。
  \item \textbf{接口与编码：}顶层模块统一采用 \code{AFDX_top} 模板的 8 个端口；端口、信号和寄存器用小写，参数用大写，重要代码段简洁注释。
\end{itemize}

\section{官方答疑口径}

表~\ref{tab:qa} 汇总与本设计相关的答疑及本设计的处理方式。

\begin{table}\centering\small
\caption{官方答疑与本设计的处理}\label{tab:qa}
\begin{tabularx}{\linewidth}{>{\centering\arraybackslash}p{1.4cm}>{\raggedright\arraybackslash}p{5.2cm}L}
\toprule
答疑 & 要点 & 本设计的处理 \\
\midrule
R3、R61 & MAC 接口为 GMII，即 PHY 接口；MII 可能达不到吞吐 & 顶层 8 个端口承载 A 网八位 GMII 字节流，含前导和 FCS \\
R30 & 不绑定器件型号 & 选用 eLinx 中的 EQ6HL130，CSG484\_H 封装 \\
R46 & MAC 可调用 eLinx IP & 自研双网 MAC（第~\ref{sec:mactx}~节） \\
R47、R48 & 不规定协议版本，参考协议规范 & 按表~\ref{tab:proto} 所列规范实现 \\
R49 & 可酌情增加接口 & 新增 B 网、应用、配置和事件端口（第~\ref{ch:if}~章） \\
R50、R62 & ready 信号可不用；valid 即本拍字节有效 & \code{rx_ready} 不使用；\code{tx_ready} 表示初始化完成 \\
R58 & 吞吐率达到 100~Mbps & 以去重后有效载荷计，设计容量 956.44~Mbps（第~\ref{sec:budget}~节） \\
R59 & 时钟频率与接口速率匹配 & 125~MHz 时钟，GMII 每拍一字节 \\
R60 & GMII 接口时延不超过 2~ms & 以应用口到应用口计时，包含收发两侧全部处理；线路空闲上界 26.360~µs \\
\bottomrule
\end{tabularx}
\end{table}

关于 R60 的计时边界：本设计从应用提交消息末字节量到对端应用交付末字节，只扣除测试设定的线路延迟。这一区间包含发送排队、组帧、MAC、线路串行和接收检查的全部时间，覆盖了 GMII 到 GMII 的区间，因此是更严格的口径。

\section{设计规格}

表~\ref{tab:spec} 把上述需求分解为编号规格。F 为功能规格，P 为性能规格，I 为接口规格，C 为实现约束。

\begin{longtable}{>{\centering\arraybackslash}p{1.0cm}>{\raggedright\arraybackslash}p{2.8cm}>{\raggedright\arraybackslash}p{10.2cm}}
\caption{设计规格}\label{tab:spec}\\
\longtablehead{3}{编号 & 规格 & 内容}
F1 & 用户数据发送 & 每端 128 条发送 VL，每条独立配置 BAG（以时钟周期计，不超过 $2^{30}$）、VLID、源/目的 IP 与端口、最大载荷和队列深度；消息 1--1471 字节；按 VL 配置从 A、B 一网或两网发送 \\
F2 & 用户数据接收 & 4096 条接收 VL，任意 16 位 VLID 可映射；检查前导、FCS、帧长、IP 头校验和、UDP 长度与校验和；每网独立的完整性检查；启用冗余管理且副本到达间隔在 \code{skew_max} 内时，同一消息只交付一次 \\
F3 & SNMP 网络管理 & SNMPv1/v2c 的 Get、GetNext、Set 和 SNMPv2-Trap；36 个管理对象；community 校验 \\
F4 & 615A 数据加卸载 & 基于 TFTP 的上传操作和介质定义卸载；LUH 头文件与数据文件 CRC-16 校验；超时重传 \\
F5 & RTC 时钟同步 & SNTPv4 主从对时；32 位秒加 32 位纳秒；回应须回显本次请求的 Transmit 时间戳 \\
P1 & 吞吐 & 每方向有效载荷不低于 100~Mbps \\
P2 & 端到端时延 & 不超过 2000~µs \\
P3 & 发送抖动 & 不超过 50~µs（“几十微秒”的检查取值） \\
P4 & 帧格式 & 帧长 64--1518 字节（含 FCS）；SN 复位后首帧为 0，之后 1--255 循环 \\
P5 & 容量 & 128 条发送 VL、4096 条接收 VL 全部可配置并可同时启用 \\
I1 & 官方端口 & \code{AFDX_top} 的 8 个端口名称、方向和位宽不变 \\
I2 & 线路接口 & A、B 两路八位 GMII，每拍一字节，帧间隔不少于 12 拍 \\
I3 & 扩展接口 & 应用发送/接收、配置总线、统计事件端口 \\
C1 & 时钟 & 单一 125~MHz 时钟，8~ns 约束下综合建立裕量不小于 2.0~ns \\
C2 & 编码规范 & 信号小写、参数大写；每个模块开头说明作用、接口和关键时序 \\
\end{longtable}

\Needspace{8\baselineskip}
\section{协议依据}

\begin{longtable}{>{\raggedright\arraybackslash}p{3.2cm}>{\raggedright\arraybackslash}p{5.0cm}>{\raggedright\arraybackslash}p{\dimexpr\linewidth-8.2cm-6\tabcolsep\relax}}
\caption{协议依据与版本}\label{tab:proto}\\
\longtablehead{3}{功能 & 依据 & 实现范围}
AFDX 网络 & ARINC 664 Part 7（2005-06-27） & 帧格式、VL 与 BAG、完整性检查、冗余管理 \\
数据加卸载 & ARINC 615A-2（2002-05-10）、ARINC 665-1、RFC 1350 & 上传、介质定义卸载；协议文件版本 A4，兼容 A3 \\
网络管理 & RFC 1157、RFC 3416、RFC 3418 & SNMPv1/v2c，BER 编码 \\
时钟同步 & RFC 4330 & SNTPv4 48 字节报文 \\
\end{longtable}

\section{追溯矩阵}

表~\ref{tab:trace} 给出每条规格对应的设计章节和验证用例，测试用例的结果见测试报告。

\begin{table}\centering\small
\caption{规格追溯矩阵}\label{tab:trace}
\begin{tabularx}{\linewidth}{>{\centering\arraybackslash}p{1.0cm}>{\raggedright\arraybackslash}p{3.0cm}L}
\toprule
规格 & 设计章节 & 验证用例 \\
\midrule
F1 & \ref{sec:txbuf}--\ref{sec:mactx} & \code{tb_tx_basic}、\code{tb_tx_bag}、\code{tb_tx_multi}、\code{tb_tx_drop}、\code{tb_sys_load} \\
F2 & \ref{sec:macrx}--\ref{sec:rxdeliver} & \code{tb_rx_basic}、\code{tb_rx_fault}、\code{tb_rx_icrm}、\code{tb_rm_phase}、\code{tb_sys_fault} \\
F3 & \ref{sec:snmp} & \code{tb_sys_snmp}、\code{tb_sys_snmp_settrap}、\code{tb_judge_snmp} \\
F4 & \ref{sec:a615} & \code{tb_sys_615a}、\code{tb_net_services} \\
F5 & \ref{sec:rtc} & \code{tb_sys_rtc}、\code{tb_rtc_recovery}、\code{tb_rtc_originate} \\
P1--P3 & \ref{sec:budget}、\ref{sec:txeng} & \code{tb_sys_perf}、\code{tb_gbe_burst} \\
P4 & \ref{sec:frame} & \code{tb_mac_tx}、\code{tb_mac_rx}、\code{tb_rx_icrm} \\
P5 & \ref{sec:txtables}、\ref{sec:rxtables} & \code{tb_judge_capacity}、\code{tb_rx_load} \\
I1--I3 & 第~\ref{ch:if}~章 & \code{tb_sys_net}、\code{tb_sys_cfg}、\code{tb_elinx} \\
C1、C2 & 第~\ref{ch:timing}~章 & 综合时序报告、网表仿真 \\
\bottomrule
\end{tabularx}
\end{table}

\chapter{总体架构}\label{ch:arch}

\section{系统框图}

系统采用“顶层应用流式包装器 \code{AFDX_top} + 端系统协议核心 \code{afdx_es}”的两级顶层结构。图~\ref{fig:system} 展示了端系统核心 \code{afdx_es} 的共享存储与双网数据通路，外围由 \code{AFDX_top} 集成发送应用流包装器 \code{afdx_app_tx} 与接收应用流包装器 \code{afdx_app_rx}，对外提供带 2 字节大端 VLID 前缀的流式应用接口（\code{tx_data}/\code{rx_data}）。端系统核心由八个一级模块组成：发送核心和接收核心构成线速数据通路；应用口与 RTC、SNMP、615A 三个服务通过服务通道共用一对收发核心；配置模块提供寄存器、统计计数器和三张配置表的访问接口；时间基准为整个设计提供自由运行的周期计数。各模块的分工见表~\ref{tab:modules}。

\figN[AFDX\_top 系统框图]{f01_system}{端系统核心的共享存储与双网数据通路。示例消息的载荷仅存一份，由描述符授权组帧并发往已启用的两网；接收侧载荷预写在本网环形缓存中，共享引擎只发布获准交付的描述符（分片报文经 \code{afdx_rx_reassembly} 重组后统一交付）。}{fig:system}

\begin{longtable}{>{\raggedright\arraybackslash}p{3.0cm}>{\raggedright\arraybackslash}p{\dimexpr\linewidth-3.0cm-4\tabcolsep\relax}}
\caption{一级模块分工}\label{tab:modules}\\
\longtablehead{2}{模块 & 职责}
\code{afdx_app_tx} / \code{afdx_app_rx} & 顶层应用流包装器：解析/添加 2 字节大端 VLID 前缀，提供双侧 $8\times 4$~KiB（32~KiB）分体应用消息缓冲、发送侧 3 级流水 128 路 VLID 比较映射与接收侧 $4096\times 16$~bit 索引反查 \\
\code{afdx_es} & 端系统协议核心顶层，集成下述八个一级核心子系统 \\
\code{afdx_tx_core} & 发送核心：入口两槽弹性缓冲、消息分页缓存与超长回滚（支持 $1\text{--}\MaxMessageBytes$~B）、128 条 VL 的 BAG 调度与 IPv4 自动分片、组帧与逐片 IPv4 首部校验和计算、双网锁步 MAC（第~\ref{ch:tx}~章） \\
\code{afdx_rx_core} & 接收核心：双网 MAC、帧头与分片解析、完整性检查（IC）、冗余管理（RM）、\ReasmContexts{} 路并发 IPv4 顺序分片重组、交付（第~\ref{ch:rx}~章） \\
\code{afdx_svc} & 服务通道：应用与三个服务按整条消息轮询发送；按 UDP 目的端口将接收消息送至对应服务或应用 \\
\code{afdx_rtc} & SNTPv4 纳秒级时钟同步（第~\ref{sec:rtc}~节） \\
\code{afdx_snmp} & SNMP 代理（第~\ref{sec:snmp}~节） \\
\code{afdx_615a} & 615A 软件数据加载目标端（第~\ref{sec:a615}~节） \\
\code{afdx_cfg} & 地址译码、寄存器、27 个统计计数器、配置表写入握手（第~\ref{ch:cfg}~章） \\*
\code{afdx_timebase} & 32 位周期计数 \code{now} 与粗粒度计时 \code{now_c} \\
\end{longtable}

\FloatBarrier
\section{模块层次与 \RtlFileCount{} 个 RTL 文件划分}

\code{AFDX_top} 顶层下共划分为 \RtlFileCount{} 个 Verilog RTL 源文件（累计 \RtlLineCount{} 行代码，包含 \DesignUnitCount{} 个 \code{module} 设计单元，源码位于 \code{rtl/} 目录），模块归属、层次关系与片上同步 RAM 实例数量见表~\ref{tab:hierarchy}（各源文件的具体职责见附录~\ref{app:files}）。全设计共例化 \RamSdpCount{} 个通用简单双端口 RAM（\code{afdx_ram_sdp}）与 160 个针对物理布线与写保持时序定制的专用同步 RAM 单元（合计 250 个同步 RAM 实例），经综合与装箱映射为 \Mem{} 个 M4K/EBR 物理存储块（占比 \MemPct\%）。

为兼顾新增的应用流封装、\ReasmContexts{} 路并发顺序分片重组、管理面协议状态机复杂度与 125~MHz 物理时序收敛，设计实施了清晰的层次解耦：
\begin{enumerate}[leftmargin=2em,itemsep=1pt]
  \item \textbf{顶层流式包装器（\code{afdx_app_tx.v}、\code{afdx_app_rx.v}、\code{afdx_app_map.v}、\code{afdx_app_payload.v}、\code{afdx_app_read.v}）：}发送包装器从 \code{tx_data} 提取前 2 字节大端 VLID，通过内部 3 级流水 128 路 VLID 比较器完成 $\text{VLID}\rightarrow\text{VL\_idx}$ 映射，并将 $1\text{--}\MaxMessageBytes$~B 载荷写入 $8\times 4$~KiB 分体缓存（\code{afdx_app_payload}，每体拆分为 $4\times 1024\times 9$~bit 子片并配置写数据 LUT 门控隔离）后经 \code{afdx_app_read} 8 信用预取器送入核心；接收包装器通过 \code{afdx_app_map}（$4096\times 16$~bit 同步 RAM）将核心输出的 12 位接收 VL 索引还原为 2 字节大端 VLID 前缀，经 32~KiB 接收应用缓冲无反压连续输出至 \code{rx_data}。
  \item \textbf{发送入口弹性级与接收 \ReasmContexts{} 路分体顺序重组（\code{afdx_tx_buf.v}、\code{afdx_rx_reassembly*.v}）：}\code{afdx_tx_buf.v} 内置两槽位 49~bit 入口弹性缓冲（\code{afdx_tx_ingress_elastic}）与 256~B 超长报文回滚 RAM（\code{afdx_tx_rollback_ram}），核心载荷池由独立封装的 \code{afdx_tx_payload_store}（$131072\times 8$~bit，带 27~bit 本地写捆绑与 LUT 写隔离）实现；接收核心在 \code{afdx_rx_engine} 准入阶段与 \code{afdx_rx_reassembly} 协同查询并预留槽位，经 \code{afdx_rx_deliver} 从双网 4~KiB 环形缓存读出后送入 2 个重组存储体（\code{afdx_rx_reassembly_bank} $\times 2$，每体含 4 张 \code{afdx_rx_reassembly_ram} 元数据表与 8 个槽位），每个槽位由 4 通道交织的 $8192\times 2$~bit 通道 RAM（\code{afdx_rx_reassembly_payload_lane}，全芯片 64 通道共配置 128 个 \code{lut_mask="8888"} 写数据隔离 LUT 驱动全部 256 个 RAM 数据输入端）构成，单拍并行写入 8 位分片字节并消除短路径保持时间违例。
  \item \textbf{\code{afdx_snmp.v}、\code{afdx_615a.v} 与 \code{afdx_mgmt_tx.v}：}管理面核心状态机（\code{afdx_snmp_engine.v} 与 \code{afdx_615a_engine.v}）由 2 分频局部时钟使能（\code{CE_DIV=2}）驱动，运行于 16~ns 多周期使能域；外围报文捕获、RAM 接口及 \code{afdx_mgmt_tx.v} 流式预取发送适配器运行于 125~MHz 单周期主时钟域。
\end{enumerate}

\begingroup\tablestyle
\setlength{\tabcolsep}{5pt}
\setlength{\heavyrulewidth}{0.8pt}
\setlength{\lightrulewidth}{0.35pt}
\begin{longtable}{@{}>{\raggedright\arraybackslash}p{3.4cm}>{\raggedright\arraybackslash}p{\dimexpr\linewidth-5.3cm-4\tabcolsep\relax}>{\raggedleft\arraybackslash}p{1.9cm}@{}}
\caption{\RtlFileCount{} 个 RTL 文件（\DesignUnitCount{} 个设计单元）分工与下级层次}\label{tab:hierarchy}\\
\longtablehead{3}{\textbf{所属域与模块} & \textbf{子模块及下级关系} & \textbf{同步 RAM\newline (SDP/专用)}}
\textbf{\code{AFDX_top}} & 顶层流式包装器：例化 \code{afdx_app_tx}、\code{afdx_app_rx} 与 \code{afdx_es} & --- \\
\addlinespace[3pt]
\textbf{\code{afdx_app_tx}} & 3 级流水 128 路 VLID 比较器、\code{descriptor_ids}（\code{afdx_ram_sdp} $\times 1$）、\code{afdx_app_payload}（$\hookrightarrow$ \code{afdx_app_payload_bank} $\times 8$ $\hookrightarrow$ \code{afdx_app_payload_slice} $\times 32$，$1024\times 9$）、\code{afdx_app_read} & 1 / 33 \\
\addlinespace[3pt]
\textbf{\code{afdx_app_rx}} & \code{afdx_app_map}（$4096\times 16$ VLID 反查 RAM）、\code{descriptor_ids}（\code{afdx_ram_sdp} $\times 1$）、\code{afdx_app_payload}（$\hookrightarrow$ 8 Bank / 32 Slice）、\code{afdx_app_read} & 1 / 34 \\
\addlinespace[3pt]
\textbf{\code{afdx_es}} & 端系统协议核心顶层：集成 \code{afdx_timebase}、\code{afdx_cfg}、\code{afdx_svc}、\code{afdx_tx_core}、\code{afdx_rx_core}、\code{afdx_rtc}、\code{afdx_snmp}、\code{afdx_615a} & --- \\
\addlinespace[3pt]
\textbf{\code{afdx_timebase}} / \textbf{\code{afdx_cfg}} / \textbf{\code{afdx_svc}} & 周期计数 \code{now[31:0]} 与 \code{now_c[15:0]}；寄存器译码与 27 个核心统计计数器；应用与 3 项服务 4 路轮询仲裁与 UDP 端口分发 & 0 / 0 \\
\addlinespace[5pt]
\textbf{\code{afdx_tx_core}} &
\code{afdx_tx_buf}（含 \code{afdx_tx_ingress_elastic} 与 \code{afdx_tx_rollback_ram}）、\code{afdx_tx_payload_store}（$131072\times 8$）\newline
\code{afdx_tx_engine}（含 IPv4 分片表 \code{afdx_ram_sdp} $\times 7$）、\code{afdx_tx_builder}（12 拍帧头与 IPv4 校验和）\newline
\code{afdx_tx_free}（\code{afdx_ram_sdp} $\times 1$）、\code{afdx_mac_tx}（$\hookrightarrow$ \code{afdx_crc32_d8} $\times 2$）\newline
核心表 \code{afdx_ram_sdp} $\times 5$：\code{tx_cfg}、\code{tx_state}、\code{tx_desc}（$256\times 47$）、\code{tx_mnext}、\code{tx_pnext} & 13 / 2 \\
\addlinespace[5pt]
\textbf{\code{afdx_rx_core}} &
\code{afdx_mac_rx} $\times 2$（各含 1 个 \code{afdx_crc32_d8}）、\code{afdx_rx_parse} $\times 2$、\code{afdx_rx_engine}、\code{afdx_rx_deliver}\newline
\code{rx_map}（32 Bank $\times 2$ 写分片，\code{afdx_ram_sdp} $\times 64$）、\code{rx_dq}（$32\times 73$~bit，\code{afdx_ram_sdp} $\times 1$）、\code{rx_ring_a/b}（\code{afdx_ram_sdp} $\times 2$）、\code{rx_cfg}（\code{afdx_rx_table_slice} $\times 8$）、\code{rx_state}（\code{afdx_rx_table_slice} $\times 11$）\newline
\code{afdx_rx_reassembly} $\hookrightarrow$ \code{afdx_rx_reassembly_bank} $\times 2$（每体含 \code{afdx_rx_reassembly_ram} $\times 4$ 元数据表 + \code{afdx_rx_reassembly_payload_slot} $\times 8$ $\hookrightarrow$ \code{afdx_rx_reassembly_payload_lane} $\times 32$，$8192\times 2$） & 67 / 91 \\
\addlinespace[5pt]
\textbf{\code{afdx_rtc}} & \code{afdx_ntp_fraction} $\times 3$（NTP 小数换算单元） & 0 / 0 \\
\addlinespace[5pt]
\textbf{\code{afdx_snmp}} &
\code{afdx_snmp_engine}（\code{CE_DIV=2} 多周期域）、\code{afdx_mgmt_tx}、\code{afdx_ram_sdp} $\times 2$：\code{request_ram}、\code{response_ram} & 2 / 0 \\
\addlinespace[5pt]
\textbf{\code{afdx_615a}} &
\code{afdx_615a_engine}（\code{CE_DIV=2} 域）、\code{afdx_mgmt_tx}、\code{afdx_ram_sdp} $\times 6$：\code{packet_rx}、\code{packet_tx}、\code{lur_ram}、\code{header_ram}、\code{load_ram}、\code{denial_ram} & 6 / 0 \\
\midrule
\multicolumn{2}{l}{\textbf{全设计合计：\RtlFileCount{} 个 RTL 文件（\DesignUnitCount{} 个设计单元），物理映射为 \Mem{} 个 M4K/EBR}} & \textbf{\RamSdpCount{} / 160} \\
\end{longtable}
\endgroup

\section{关键设计决策}

\begin{longtable}{>{\raggedright\arraybackslash}p{2.6cm}>{\raggedright\arraybackslash}p{3.4cm}>{\raggedright\arraybackslash}p{2.4cm}>{\raggedright\arraybackslash}p{\dimexpr\linewidth-8.4cm-8\tabcolsep\relax}}
\caption{关键设计决策}\label{tab:decisions}\\
\longtablehead{4}{设计问题 & 候选方案 & 选择 & 理由}
时钟结构 & 双网各自时钟域；统一 125~MHz 单时钟 & 单时钟 & GMII 接口线速为 1~Gbit/s，载荷容量为 956.44~Mbps，满足赛题 100~Mbps 的要求；核心内部收发与控制逻辑在同一时钟域协同运行 \\
顶层应用接口 & 直接暴露多信号核心口；2B VLID 前缀流式包装 & 2B VLID 流式包装 & 严格符合官方 1026/1030 答疑 R120-3/R124/R126 要求，保持 \code{AFDX_top} 模板端口定义，同时内部 \code{afdx_es} 保留完整显式控制接口 \\
长包分片与重组 & 整包拆为多条独立消息入队；单描述符按片调度 + \ReasmContexts{} 路双 Bank 顺序重组 & 单描述符按片调度 + \ReasmContexts{} 路重组 & $8192$~B 长包仅占 1 个消息描述符与 16 个缓存页，由 \code{tx_engine} 记录剩余字节数与片偏移并按 BAG 逐片授权；接收端在逐帧 IC/RM 后经 2-Bank $\times$ 8 槽位支持 \ReasmContexts{} 条 VL 并发顺序偏移重组 \\
MAC 实现 & eLinx 官方 MAC IP；自研 & 自研 & 直接控制前导码、FCS 和帧间隔；输出寄存器直接驱动 GMII，空闲发送延迟为 0 拍 \\
双网 MAC & 双 MAC 独立发送；单 MAC 同时驱动两网 & 单 MAC 驱动两网 & 两网副本同拍起始，节省一套 MAC；单网 VL 也按整帧占用双网发送时隙（第~\ref{sec:mactx}~节） \\
VL 状态存储 & 寄存器阵列；块 RAM 状态表 & 块 RAM 状态表 & 寄存器用量随 VL 数量线性增长；采用块 RAM 后，由对应引擎按固定流水读写，配置写入须等待访问许可 \\
VLID 映射 & 全地址直接查表；哈希加 CAM；限制 VLID 范围 & 直接查表 & 接收核心采用 32-Bank $\times 2$ 分片 $65536\times 13$~bit 直接查表（3 级或树归并流水，无哈希冲突）；应用发送采用 3 级流水 128 路比较器完成映射 \\
发送调度 & 轮询扫描；最早到期优先；预取窗口扫描 & 预取窗口扫描 & 每拍至多一级比较；512 拍窗口为扫描回访与组帧留出时间，无争用时可在允许时刻启动发送 \\
接收写入 & 整帧检查后写入；边接收边预写入 & 先预写入后提交 & 载荷边接收边写入环形缓存，检查失败时回退预写入指针，释放本帧占用的空间 \\
服务发送 & 服务独立发送通路；按消息轮询共享发送核心 & 共享发送核心 & 服务报文与用户报文遵守同一 BAG 调度；无服务报文时应用直通，零附加延迟 \\
服务实现 & 软核处理器；硬件状态机 & 硬件状态机 & 协议解析与报文组装均由硬件状态机完成；RTC 时间戳在报文握手当拍采样 \\
服务状态序列 & 各分支独立展开；公共状态序列加返回状态寄存器 & 公共状态序列 & SNMP 和 615A 复用取字节、整数拼接与报文输出过程，统一同步 RAM 的访问时序 \\
\end{longtable}

\section{时钟与复位}\label{sec:clkrst}

端系统逻辑采用单一主时钟 \code{clk}，工作频率 125~MHz，周期 8~ns，所有内部触发器均在 \code{clk} 上升沿同步更新。\code{reset} 为高有效的同步复位，外部须保证其释放与 \code{clk} 上升沿严格同步；板级实现时由物理复位同步器提供。

为满足航电高层管理协议（SNMP 与 ARINC 615A）在 125~MHz 主频下的静态时序收敛要求，系统在统一主时钟树下设置了管理面局部时钟使能结构：
\begin{enumerate}[leftmargin=2em,itemsep=1pt]
  \item \textbf{多周期引擎结构（$\text{CE\_DIV}=2$）：}在 \code{afdx_snmp} 与 \code{afdx_615a} 顶层生成每两拍有效一次的时钟使能信号 \code{ce}，并通过多级低扇出寄存器复制树（\code{CE_GROUPS}）驱动核心引擎内部逻辑。全工程共有 \MulticycleRegs{} 个触发器（由 \MulticycleRules{} 条多周期约束规则覆盖）约束为 2 周期建立路径（16~ns 建立时间预算），大幅缓解深层协议分支的组合时序压力。
  \item \textbf{单周期外围与接口隔离：}应用流包装器、以太网报文接收捕获、分片重组引擎、双端口协议 RAM 写口、32 位配置总线、RTC 纳秒时钟基准以及用于流式预取的 \code{afdx_mgmt_tx} 适配模块均保持在 125~MHz 单周期全速运行。这种结构既避免了异步跨时钟域（CDC）带来的亚稳态风险与同步 FIFO 开销，又使关键控制与存储接口维持周期级确定性。
\end{enumerate}

核心顶层采用与 \code{clk} 同步的 GMII 字节流接口。连接外部物理 PHY 时，接收数据随 PHY 恢复的 RX\_CLK 变化，发送侧需向外输出 GTX\_CLK；因此板级外围需在接收引脚与核心之间设置帧级弹性缓冲，将 RX\_CLK 域字节流同步至 \code{clk} 域，发送侧则由 \code{clk} 经输出寄存器驱动 GTX\_CLK。

复位结束后，发送核心清零 \code{tx_state}、\code{tx_cfg} 及分片进度表，并将 256 个页号写入空闲页池；接收核心清零 \code{rx_map}、\code{rx_cfg} 和 \code{rx_state}，分片重组引擎复位 \ReasmContexts{} 个重组槽位状态，应用包装器清零映射表与描述符槽。深度最大的 \code{rx_map} 物理划分为 32 个子库并行清零，配合初始化定时器计数需 65536 拍。两个核心的 \code{init_done} 经流水打拍并在核心顶层相与锁存输出（从解除复位到 \code{afdx_es} 就绪为 65539 拍，即 524.312~µs；顶层 \code{AFDX_top} 的应用入口 \code{tx_ready} 在包装器同步完成清零后于约第 65,550 拍拉高）。表初始化完成之前，\code{cfg_ready} 保持 0、配置写不被接受，凡在此期间开始接收的以太网帧均整帧忽略。

\section{时间基准}\label{sec:timebase}

\code{afdx_timebase} 为全设计产生三个信号：
\begin{itemize}
  \item \code{now[31:0]}：自由运行的周期计数，复位为 0，每拍加 1，约 34.36~s 回绕。帧起点以 1 拍（8~ns）为粒度，所有时间戳的先后比较均通过作差并检查符号位实现，因此 BAG 可配置为不超过 $2^{30}$（约 8.59~s）的任意整数拍（例如 100003 拍）。
  \item \code{now_c[15:0]}：粗粒度计时，等于 \code{now[27:12]}，每 4096 拍（32.768~µs）加 1，仅用于接收侧的冗余管理。
  \item \code{byte_en}：字节使能，复位期间为 0，此后恒为 1，即每拍都是字节边界。
\end{itemize}

\section{端到端数据通路}\label{sec:path}

图~\ref{fig:path} 给出一帧用户消息从发送端应用到接收端应用的完整路径。在顶层 \code{AFDX_top} 中，应用发送流先经 \code{afdx_app_tx} 剥离 2 字节大端 VLID 前缀并映射为内部 VL 索引，随后进入 \code{afdx_es} 发送侧，依次经过入口弹性级、分页缓存、调度与分片引擎、组帧器和双网 MAC；两份以太网帧经 A、B 网到达接收端后，各自经 MAC 和解析器检查，由接收引擎做完整性检查和冗余管理，送入交付模块；单帧消息直通输出，IPv4 分片消息则进入 \code{afdx_rx_reassembly} 完成分片拼接与全报文 UDP 校验，最后经 \code{afdx_app_rx} 加上 2 字节大端 VLID 前缀流式交付给上层应用。

帧内数据通路每拍处理一个字节，对应 1~Gbit/s。发送侧先把应用消息写入分页缓存，满足调度条件后，一边读取载荷，一边生成以太网帧；接收侧将到达的载荷预写入环形缓存，帧检查通过后再提交给交付与重组模块。各段的拍数和上界在第~\ref{sec:budget}~节推导。

\section{设计预算}\label{sec:budget}

本节由各模块的状态机拍数推导吞吐、时延、抖动、缓存和握手等待的设计值，作为第~\ref{ch:tx}--\ref{ch:cfg}~章模块设计的约束。拍数均以 125~MHz 时钟计，1 拍 = 8~ns；仿真实测结果见测试报告。

本节及第~\ref{ch:tx}--\ref{ch:cfg}~章以相邻两级输出有效时刻之间的周期差作为阶段延迟。状态机从 IDLE 态响应请求并启动操作需要 1 个周期，因此阶段延迟比操作本身的执行时间多 1 拍。例如，调度引擎的入队阶段共耗时 6 拍，其中响应请求需 1 拍，执行 ENQ 操作需 5 拍。

\subsection{吞吐}

单帧载荷为 $n$ 字节（$1\le n\le 1471$）的以太网帧连同 8 字节前导码、4 字节 FCS 和 12 拍帧间隔，在物理线路上占用的时长为
\begin{equation}
C(n) = 8 + 47 + \max(n,17) + 12 = 67 + \max(n,17)\ \text{拍}.
\end{equation}
满载帧 $n=1471$ 占 $C=1538$ 拍，最短帧 $n\le 17$ 占 84 拍。每网每方向的单帧最大有效载荷容量为
\begin{equation}
R_\text{max} = \frac{1471 \times 8~\text{bit}}{1538 \times 8~\text{ns}} = 956.44~\text{Mbps}.
\end{equation}
设 BAG 为 $B$ 拍，$N$ 条 VL 各按该间隔发送满载帧时，单向有效载荷率为 $N \times 1471 \times 8~\text{bit}/(B \times 8~\text{ns})$，可调度条件为 $N \times C \le B$。当 $\text{BAG}=1$~ms（$B=125000$ 拍）时：
\begin{itemize}[leftmargin=2em,itemsep=1pt]
  \item \textbf{声明 63~VL 满载规划基准（R119/R138-3/R139-3）：}单端激活 $N=63$ 条发送 VL 与 $63$ 条接收 VL，每毫秒占用物理线路 $63\times 1538 = 96,894$ 拍，线路占用率为 $96894/125000 = 77.5152\%$（约 $77.5\%$），单向去重有效载荷吞吐率为 $63\times 1471\times 8~\text{bit}/1~\text{ms} = \ThruSixtyThree$~Mbps（若计入顶层 2 字节 VLID 前缀则输入流速率为 $\ThruSixtyThreeRaw$~Mbps），双向全双工合计有效载荷吞吐率为 $\ThruSixtyThreeBi$~Mbps。
  \item \textbf{77~VL 与 81~VL 超规格极限压力基准：}$N=77$ 条 VL 时单向有效载荷率为 $\Thru$~Mbps（线路占用率 $94.7\%$）；物理极限 $N = \lfloor 125000/1538 \rfloor = 81$ 条 VL 时，单向有效载荷率达 $\ThruEightyOne$~Mbps（线路占用率 $99.66\%$）。
\end{itemize}

线速通路上各环节的吞吐预算如下：
\begin{itemize}
  \item \textbf{应用消息接收：}顶层 \code{afdx_app_tx} 以每拍 1 字节无间断接收带 2 字节 VLID 前缀的消息并写入 $8\times 4$~KiB 分体缓冲；核心 \code{afdx_tx_buf} 正常每拍接收一字节，每次申请新缓存页暂停 6 拍；消息接收结束后，写入描述符及等待入队应答共暂停约 10 拍。一条占 3 页的 1471 字节消息在核心口共耗时 $1471 + 3\times 6 + 10 = 1499$ 拍，小于该帧占用线路的 1538 拍；一条占 16 页的 $8192$ 字节最大分片消息写入核心耗时 $8192 + 16\times 6 + 10 = 8298$ 拍（66.384~µs），而其拆分为 6 个以太网帧跨 5 个 BAG 周期发送，入口写入速率远高于线路消耗速率。
  \item \textbf{授权与准备：}引擎完成一帧或一片授权约需 9--12 拍，组帧器 \code{afdx_tx_builder} 采用 12 拍（步骤 \code{0..11}）准备序列在帧头阶段硬件计算 20 字节 IPv4 首部校验和（IPv4 首分片 UDP 校验和按 RFC 768 固定置 \code{0x0000}，无需预读整包载荷），并在计划起始时刻前 8 拍（\code{START_LEAD_REGISTERED = -8}）完成准备，均在当前帧发送期间重叠完成，远小于最短帧的 84 拍时隙。
  \item \textbf{接收重组与交付：}核心交付口与 \ReasmContexts{} 路重组引擎均支持每拍 1 字节线速吞吐；最短以太网帧占 84 拍，接收引擎 30 拍（相位 \code{0..30}）帧处理日程可在连续两帧最短帧到达期间处理双网副本；单网到达的载荷率至多 $1471/1538$ 字节每拍，低于交付与重组速率。
\end{itemize}

\subsection{端到端时延}

端到端时延的起止点定义见第~\ref{ch:req}~章。下文按收发通路无其他消息争用、当前 VL 已满足 BAG 发送间隔的基准条件，给出核心分段口径下各区段的典型拍数和设计上界（表~\ref{tab:latbudget}），发送侧和接收侧的逐拍时序见图~\ref{fig:pipe}。若从顶层 \code{AFDX_top} 流式接口计时（R119/R139 全消息入口首字节至接收末字节），由于 \code{afdx_app_tx} 与 \code{afdx_app_rx} 均采用“整包收齐并完成 VLID 映射后转发”的存储转发边界，单条 $n$ 字节消息在无排队争用时额外叠加顶层发送写入（$n+2$ 拍加跨层预取）与顶层接收缓存转发（$n$ 拍加 2 字节 VLID 前缀输出）；例如在 63~VL 周期稳态激励下，$1471$~B 消息顶层全消息最大端到端时延实测为 \LatSixtyThreePMax~µs，旋转突发（B2）最大时延为 \LatSixtyThreeBTwo~µs；对于拆分为 $F$ 片的 $8192$~B 长消息（$F=6$，$\text{BAG}=1$~ms），端到端完成时延由 $(F-1)\times\text{BAG}$ 主导，实测为 $5212.304$~µs。

图~\ref{fig:pipe}(a) 中，消息入队时直接成为候选。若末字节在 RECV 状态下正常接收，且无需在消息尾部申请新页，由于载荷写口增加了 1 拍流水寄存、核心输出端插入了 1 级流缓冲（双缓冲结构），从接收末字节到输出首个前导码字节固定为 32 拍（0.256~µs）；按 1022 答疑 R95 规则计至以太网帧 SFD 结束（8 字节前导码尾），发送处理时延 $L_\text{tx}$ 典型为 39 拍（0.312~µs）。单字节消息或末字节跨页时多约 6 拍。图~\ref{fig:pipe}(b) 给出单网单帧的接收过程，由于 \code{rx_map} 采用 32 子库多 Bank 分块存储增加 2 拍读回流水，从 GMII 最后一个 FCS 字节到首个交付字节由 24 拍更新为 26 拍（0.208~µs）。

\figN[收发两侧的逐拍时序（无其他消息争用、VL 已到期）]{f20_pipe}{无其他消息争用、VL 已到期时的流水时序：(a) 发送准备与起始；(b) 接收检查与交付}{fig:pipe}

\begin{table}\centering\small
\caption{端到端时延预算（拍）}\label{tab:latbudget}
\begin{tabularx}{\linewidth}{>{\raggedright\arraybackslash}p{2.5cm}>{\raggedleft\arraybackslash}p{1.2cm}>{\raggedleft\arraybackslash}p{1.2cm}L}
\toprule
区段 & 典型 & 上界 & 构成 \\
\midrule
发送侧至前导码 & 32 & 100 & 分页缓存写载荷寄存 1 拍、写描述符并请求入队 3 拍，引擎入队 6 拍，授权 9 拍，组帧准备 11 拍，输出流缓冲与 MAC 起始 2 拍；计至 SFD 结束（$L_\text{tx}$）为 39 拍 \\
线路 & 1525 & 1525 & 前导码首字节至 FCS 末字节，共 $8+1518$ 字节（$n=1471$） \\
接收至首字节 & 26 & 200 & MAC 检测帧尾 3 拍，解析器分类 3 拍，多 Bank 分块映射与引擎判定入队 14 拍，交付模块读块 RAM 队头并读取环形缓冲 6 拍；上界另计另一网描述与过期清理让出 \\
逐字节交付 & $n-1$ & $n-1$ & 首字节到末字节，每拍一字节，$n=1471$ 时为 1470 拍 \\
\midrule
合计 & 3053 & 3295 & 24.424~µs（遇 AGE 扫描 1 拍为 3054 拍 / 24.432~µs） / 26.360~µs \\
\bottomrule
\end{tabularx}
\end{table}

存在其他消息时，还须计入发送排队、提前授权占槽、线路占用与接收排队交付的等待；端到端时延的上界由各 VL 的 BAG、帧长和流量配置共同决定。

\FloatBarrier
\subsection{发送抖动}\label{sec:jitter}

记 VL $v$ 的帧允许时刻为 $T_v$（上一帧实际起始时刻加 BAG），实际起始时刻为 $S_v$，发送抖动为 $J_v = S_v - T_v$，即同一 VL 相邻两帧起始间隔与 BAG 之差。

与第 3.7.2 节考察消息到达时 VL 已到期的直通延迟不同，本节考察消息在允许时刻之前已提前排队时的准点发送精度。当线路和待发槽位均未被其他帧占用时，发送抖动为 0。调度器能否在 $T$ 时刻准点启动发送，取决于候选能否在 $T$ 之前完成扫描、授权和准备。一条 VL 在 $T-512$ 进入提前窗口，此后：
\begin{itemize}
  \item \textbf{扫描回访：}无候选时，引擎每 2 拍检查一条 VL，轮询 128 条共需 256 拍。每扫描 16 条 VL 至多插入一次入队操作（6 拍），每发送一帧还需插入一次发送起始确认（TX Start Confirm）（4 拍）；每次插入操作都会清空 5 级扫描流水并重新发出（约 6 拍，见第~\ref{sec:txeng}~节），折合每次入队约 12 拍、每次发送起始确认（TX Start Confirm）约 10 拍。计入这些开销后，任一 VL 相邻两次被检查的间隔约为 400 拍。
  \item \textbf{授权与准备：}授权 9 拍、组帧器准备 11 拍、起始判断 2 拍，共 22 拍。
\end{itemize}
$400 + 22 = 422 < 512$，留有 90 拍裕量。因此在上述条件下，MAC 在 $T$ 时刻准点输出前导码首字节。

线路争用时，物理链路同一时刻只能发送一帧，组帧器也仅设一个待发槽位。$v$ 到期且消息已排队时，可能因以下三种原因推迟发送：
\begin{itemize}
  \item \textbf{线路占用：}排在 $v$ 前面的帧逐个发完，每帧线路占用 $C(n_u)=67+\max(n_u,17)$ 拍。
  \item \textbf{提前授权占槽等待：}扫描按 VL 索引轮转，选中首条已到期或将在 512 拍内到期的候选 VL。若扫描先遇到一条尚未到期的 VL $u$ 并发出授权，组帧器的待发槽位将被 $u$ 占用，直到 $T_u$ 时刻才启动发送，期间 $v$ 无法获得授权。每次此类提前授权引入的占槽等待为 $w_u < 512$ 拍。
  \item \textbf{轮询扫描延迟：}扫描从上次选中位置继续向后轮询，在遍历到 $v$ 之前所选中的 VL 均先于 $v$ 发送；扫描与这些帧的发送重叠进行，扫描范围不超过一圈，连同 $v$ 自身的授权和准备不超过约 422 拍。
\end{itemize}
记 $Q_v$ 为在 $v$ 之前发送的其他 VL 帧的集合，则
\begin{equation}
J_v \le \Big(\sum_{u \in Q_v} \big(C(n_u) + w_u\big) + 422\Big) \times 8~\text{ns},
\label{eq:jitter}
\end{equation}
其中 $w_u$ 只对在自身允许时刻之前被授权的帧非零。服务报文各走自己的 VL，同样按 BAG 调度，已包含在 $Q_v$ 中。

一次争用只造成一次相位后移。下一次允许时刻按实际起始时刻加 BAG 计算（第~\ref{sec:txeng}~节）。多条 VL 发生发送争用后，后发 VL 的相位后移，此后各自按 BAG 更新允许时刻。要持续保持错开的相位，需满足两个条件：
\begin{itemize}
  \item 各 VL 的 BAG 互为整数倍。ARINC 664 P7 规定 BAG 取 $2^k$~ms，满足这一点；任取整数拍的 BAG 会使相位相互漂移，周期性地再次发生争用。
  \item VL 在允许时刻之前已有消息排队。计时到期时若队列为空，\code{timer_active} 被清除，下一条消息到达即发，相位改由应用的提交时刻决定。
\end{itemize}
两个条件都满足时，争用只在启动时发生，稳定后同一 VL 相邻两帧的起始间隔恰为 BAG。

以 P3 的 50~µs 作为发送抖动的检查取值，按式~(\ref{eq:jitter}) 对应 $\sum_{u\in Q_v}(C(n_u)+w_u) \le 6250-422 = 5828$ 拍。例如：
\begin{itemize}
  \item 同时到期的其他 VL 至多 3 条最长帧（$w_u=0$）：$3\times 1538 = 4614$ 拍，满足；
  \item 前面有 2 条最长帧且都被提前授权：$2\times(1538+512) = 4100$ 拍，满足；
  \item 77 条最短帧同时到期：$76 \times 84 = 6384$ 拍，不满足。
\end{itemize}
同时参与争用的 VL 数量超出上述示例时，主机须错开各 VL 首条消息的提交时刻，使任一 VL 之前排队发送的帧满足该条件；在上述两个条件下，错开只需在启动时做一次（第~\ref{sec:cfgrule}~节）。

与 ARINC 664 P7 的抖动公式 $40~\text{µs} + \sum(20+L_\text{max})\times 8~\text{bit}/1~\text{Gbit/s}$ 相比，$C(n)\times 8~\text{ns}$ 恰等于公式中对应帧的线路占用时间；本设计的扫描开销 422 拍（3.4~µs）与至多 8 次提前授权等待（共 4096 拍）之和小于公式中的 40~µs 项。因此对满足 P7 公式、且排在前面的提前授权帧不超过 8 个的配置，式~(\ref{eq:jitter}) 给出的发送抖动上界不超过按 P7 抖动公式计算的值，也不超过 P7 规定的 500~µs 抖动上限。

\subsection{存储预算与缓存规划}\label{sec:bufbudget}

在存储架构规划中，设计严格解耦了“全量静态链路配置”与“动态在途流量缓冲深度”：
\begin{itemize}
  \item \textbf{静态配置容量：}硬件配置表与映射表完全支持 128 条发送 VL 与 4096 条接收 VL（加上接收映射表 65536 个 VLID 全地址空间）的并发静态配置；接收包装器 \code{afdx_app_map} 同步维护 $4096\times 16$~bit 的 $\text{VL\_idx}\rightarrow\text{VLID}$ 反查映射。
  \item \textbf{顶层应用流包装缓冲：}\code{afdx_app_tx} 与 \code{afdx_app_rx} 各自配备 $8\times 4$~KiB（共 32~KiB）分体消息缓存与 32 项描述符队列，单侧即可容纳 4 条 $8192$~B 最长消息或 22 条 $1471$~B 单帧消息，吸收顶层连续突发与核心分页申请之间的短时速率差；\code{afdx_tx_buf} 另设 256~B 回滚 RAM（\code{afdx_tx_rollback_ram}），对超过 $8192$~B 的异常超长消息安全回收已分配缓存页并上报 \code{ev_app[1]}。
  \item \textbf{动态发送载荷缓存：}发送核心配置 128~KiB 共享载荷存储池（\code{afdx_tx_payload_store}，划分为 256 页 $\times$ 512 字节），采用基于单向链表的动态页分配机制。单帧最长报文（1471 字节）占 $\lceil 1471/512 \rceil = 3$ 页，$8192$~B 最大分片消息占 16 页。在团队声明的 63~VL 满载规格下，63 条单帧大包同时在途占 $63\times 3 = 189$ 页（73.8\%）；在 60 条 $1471$~B 短包与 $1472/3000/8192$~B 三条分片长包并发的 63~VL 混合满载工况下，峰值占 $180 + 3 + 6 + 16 = 205$ 页（80.1\%），均留有充裕空闲页裕量（77/81~VL 超规格单帧压力下分别占 231/243 页）；当超规格突发导致空闲页耗尽时，硬件执行安全丢弃并置位统计脉冲，确保系统免于死锁。
  \item \textbf{接收端环形缓冲、描述队列与 \ReasmContexts{} 路重组缓存：}A、B 网各设 4096 字节接收环形缓冲（\code{rx_ring_a/b}），可安全容纳连续两个最长以太网帧（$2\times 1471$ 字节）。交付描述队列 \code{rx_dq} 设为 $32\times 73$~bit；分片帧经 IC/RM 去重通过后立即由 \code{afdx_rx_deliver} 搬移至 \code{afdx_rx_reassembly} 的专用重组槽位（2-Bank $\times$ 8 槽位 = \ReasmContexts{} 路并发上下文，每槽位 8~KiB，总计 128~KiB 重组载荷 RAM），使跨多个 BAG 周期到达的长报文分片完全不驻留占用 4~KiB 线速环形缓冲；超时未收齐的残缺上下文在 $250,000,000$ 拍（2~s）后由硬件自动老化回收。
\end{itemize}

\subsection{握手等待}

\begin{longtable}{>{\raggedright\arraybackslash}p{2.4cm}>{\raggedright\arraybackslash}p{\dimexpr\linewidth-2.4cm-4\tabcolsep\relax}}
\caption{握手等待与初始化预算}\label{tab:handshakebudget}\\
\longtablehead{2}{接口或操作 & 等待上界与计算依据}
入队应答 & 发送引擎一次只执行一种操作（含分片元数据读写的授权操作约 9--12 拍）。每扫描满 16 次或已有候选时可插入入队操作，入队应答等待不超过 64 拍。 \\
\code{m_ready} 反压 & 核心入口经两槽弹性缓冲隔离后，取页暂停 6 拍；消息末尾写描述符、入队应答、异常时页面回滚释放和事件脉冲存在短时停顿；顶层 \code{AFDX_top} 则由 32~KiB \code{afdx_app_tx} 缓存吸收核心跨页停顿，在缓存未满时保持 \code{tx_ready=1} 连续接收整条应用消息。 \\
应用口等待 & \code{app_tx_ready} 受服务通道轮询仲裁控制（第~\ref{sec:svcch}~节）。三个服务均有报文待发时，应用至多等待其各发送一条：RTC 报文 48 字节、SNMP 响应不超过 256 字节、615A 报文不超过 516 字节；经管理发送适配器流式预取后，字节发送以线速每拍 1 字节推进，连同预取启动与入队等待，握手等待上界约为 1100 拍（约 8.8~µs）。 \\
配置写 & 向配置表写入时，需等待对应收发核心通过 \code{cfg_ok} 给出许可。发送引擎最坏连续执行发送起始确认、入队和分片授权约 28 拍；接收引擎单次帧处理日程（含分片重组槽位查询与预留）为 30 拍（相位 \code{0..30}），双网连续处理两份描述加 AGE 中断等待上界约 64 拍（第~\ref{sec:cfgok}~节）。 \\
初始化 & 最大的 \code{rx_map} 逻辑容量 65536 项，硬件划分为 32 个子库并行清零（2048 拍），配合全局清表定时器计数及跨级输出寄存，从解除复位至核心 \code{afdx_es} 就绪耗时 65539 拍（524.312~µs），顶层 \code{AFDX_top} \code{tx_ready} 于约第 65,550 拍有效。 \\
\end{longtable}

服务通道按整条消息轮询，限制了服务连续占用消息通路的时间，表中的约 8.8~µs 是提交口的握手等待预算。入队后，用户与服务报文按各自 VL 的 BAG 参与调度，共享组帧器和发送 MAC；服务 VL 的线路占用与提前授权等待统一计入式~(\ref{eq:jitter}) 及第~\ref{sec:cfgrule}~节的配置预算。


\chapter{接口设计}\label{ch:if}

\section{顶层端口}

按官方 1026/1030 答疑澄清（R110、R120-3、R124、R126），顶层模块 \code{AFDX_top} 完整保留选题模板的 8 个核心端口（名称、方向和位宽不变，其中 \code{tx_data}/\code{rx_data} 直接承载带 2 字节大端 VLID 前缀的应用层流式消息），并显式引出 A、B 双网千兆 GMII 物理端口、32 位配置总线及硬件统计事件脉冲。顶层共声明 26 个端口信号，总位宽为 \Pins{} 位（占 EQ6HL130-CSG484\_H 可用 339 个用户 I/O 的 \PinPct\%）。顶层端口定义见表~\ref{tab:ports}；内部协议核心 \code{afdx_es} 同时保留显式的核心应用端口 \code{app_tx_*}/\code{app_rx_*}（第~\ref{sec:intif}~节）。

\begin{longtable}{>{\raggedright\arraybackslash}p{5.2cm}>{\centering\arraybackslash}p{1.0cm}>{\raggedright\arraybackslash}p{7.6cm}}
\caption{\code{AFDX_top} 顶层端口定义（26 组信号，共 \Pins{} 引脚）}\label{tab:ports}\\
\longtablehead{3}{端口 & 方向 & 含义}
\multicolumn{3}{l}{\textit{选题模板端口（2 字节 VLID 前缀应用流接口）}} \\*
\code{clk} & 入 & 125~MHz 全局主时钟 \\
\code{reset} & 入 & 同步复位，高有效 \\
\code{tx_data[7:0]}、\code{tx_valid} & 入 & 应用层发送字节流：\code{tx_valid} 连续为 1 区间构成一条消息（前 2 字节大端 VLID + $1\text{--}\MaxMessageBytes$~B 载荷），相邻消息间至少空闲 1 拍 \\
\code{tx_ready} & 出 & 应用发送就绪：内部表初始化完成且发送包装缓存可接纳新消息时为 1 \\
\code{rx_data[7:0]}、\code{rx_valid} & 出 & 应用层接收字节流：\code{rx_valid} 连续为 1 区间输出一条去重/重组后的完整消息（前 2 字节大端 VLID + 载荷），相邻消息间至少空闲 1 拍 \\
\code{rx_ready} & 入 & 按官方答疑 R62 口径保留端口兼容，硬件内部不使用（无反压连续交付） \\
\midrule
\multicolumn{3}{l}{\textit{A/B 双冗余网络 GMII 物理接口}} \\*
\code{gmii_txd_a[7:0]}、\code{gmii_tx_en_a} & 出 & A 网 GMII 发送字节流与发送使能（含前导码、SFD 与 FCS） \\
\code{gmii_txd_b[7:0]}、\code{gmii_tx_en_b} & 出 & B 网 GMII 发送字节流与发送使能，与 A 网锁步同拍起始 \\
\code{gmii_rxd_a[7:0]}、\code{gmii_rx_dv_a} & 入 & A 网 GMII 接收字节流与数据有效 \\
\code{gmii_rxd_b[7:0]}、\code{gmii_rx_dv_b} & 入 & B 网 GMII 接收字节流与数据有效 \\
\midrule
\multicolumn{3}{l}{\textit{配置与统计事件接口}} \\*
\code{cfg_wr}、\code{cfg_rd}、\code{cfg_addr[17:0]}、\code{cfg_wdata[31:0]} & 入 & 32 位字寻址配置总线写/读使能、地址与写数据 \\
\code{cfg_rdata[31:0]}、\code{cfg_rvalid}、\code{cfg_ready} & 出 & 配置读返回数据、读数据有效（固定 2 拍）与写就绪 \\
\code{ev_tx[6:0]} & 出 & 核心发送统计脉冲：\code{accept}、\code{drop_dis}、\code{drop_full}、\code{drop_long}、\code{drop_nopage}、\code{tx_a}、\code{tx_b} \\
\code{ev_a[7:0]}、\code{ev_b[7:0]} & 出 & A/B 网接收分类与异常脉冲（含 \code{ev_a/b[7]} 分片重组拒绝），见表~\ref{tab:rxev} \\
\code{ev_rm_drop}、\code{ev_deliver} & 出 & 冗余管理去重丢弃脉冲、有效消息交付脉冲 \\
\code{ev_app[2:0]} & 出 & 顶层应用流包装器丢包脉冲：\code{[0]} 发送未映射 VLID 丢弃，\code{[1]} 发送超 $8192$~B 或缓存满回滚丢弃，\code{[2]} 接收包装缓存溢出丢弃 \\
\bottomrule
\end{longtable}

顶层参数 \code{CE_DIV}、\code{TX_PAGE_AW}、\code{TX_OFF_AW}、\code{LOOKAHEAD}、\code{RX_RING_AW}、\code{RX_DQ_AW}、\code{RX_CONFIG_BANKS} 的默认值分别为 2、8、9、512、12、5、16，透传至核心 \code{afdx_es}。

\subsection{顶层端口映射与网络分层定义}

按照 ARINC 664 Part 7 标准分层与官方 1022/1026/1030 答疑澄清（R62、R110、R111、R112、R120-3、R124、R126）：
\begin{enumerate}[leftmargin=2em,itemsep=1pt]
  \item \textbf{顶层应用流接口（模板 8 端口）：}\code{clk}、\code{reset}、\code{tx_data}、\code{tx_valid}、\code{tx_ready}、\code{rx_data}、\code{rx_valid}、\code{rx_ready} 构成面向机载分区应用的流式收发边界；每条消息以 2 字节大端 VLID 为首，以 \code{valid} 连续高电平定界、消息间至少空 1 拍定帧。
  \item \textbf{双网 GMII 物理接口：}独立引出 A 网（\code{gmii_txd_a}/\code{gmii_tx_en_a}、\code{gmii_rxd_a}/\code{gmii_rx_dv_a}）与 B 网（\code{gmii_txd_b}/\code{gmii_tx_en_b}、\code{gmii_rxd_b}/\code{gmii_rx_dv_b}）八位千兆 GMII 端口，传输含 7 字节前导码、1 字节 SFD 与 4 字节 FCS 的完整以太网字节流。
  \item \textbf{配置同步窥探与核心显式接口：}主机通过 \code{cfg_*} 配置发送 VL 表与接收映射表时，\code{AFDX_top} 硬件自动窥探 \code{tx_cfg} 字 2（VLID）与 \code{rx_map} 写入事务，同步更新 \code{afdx_app_tx} 与 \code{afdx_app_rx} 的 VLID 映射，无需额外配置指令；内部 \code{afdx_es} 保留显式 \code{app_tx_*}/\code{app_rx_*} 信号供核心级与服务联合测试直接复用。
  \item \textbf{MAC 实现与评测依据（R111、R112）：}MAC 模块采用纯 RTL 自研实现；功能仿真波形、日志、覆盖率记录、门级网表与布局布线后 STA 报告共同构成签核依据。
\end{enumerate}

\section{GMII 接口}\label{sec:gmii}

A、B 两网各有一组八位 GMII，在 \code{valid} 为 1 的每拍传送一个字节。

\begin{itemize}
  \item \textbf{发送：}\code{gmii_tx_en_a/b}（核心内为 \code{tx_valid}）从前导码首字节至 FCS 末字节连续为 1，帧间至少空闲 12 拍，见图~\ref{fig:wgmii}。GMII 输出直接来自 MAC 的输出寄存器，发送适配延迟为 0。若某条 VL 未使能其中一路网络，则该帧发送期间对应网络的使能与数据全程保持为 0。
  \item \textbf{接收：}\code{gmii_rx_dv_a/b}（核心内为 \code{rx_valid}）的上升沿表示一帧开始，下降沿表示该帧结束。相邻两帧的低电平间隔最短为 1 拍。MAC 先寄存输入，再逐拍查找帧起始定界符（SFD）、累加 CRC 并剥离 FCS。
  \item \textbf{复位：}复位期间输出为 0；对于表初始化完成前开始接收的帧，接收核心将整帧忽略（第~\ref{sec:clkrst}~节）。
  \item \textbf{时钟与错误信号：}接口采用全双工模式，收发字节均与 \code{clk} 同步，不含 TX\_CLK、RX\_CLK、TX\_ER、RX\_ER、CRS 和 COL。接收核心根据字节流检查前导码、帧长和 FCS；连接外部 PHY 时，外围适配逻辑须处理 PHY 错误标记并丢弃带错帧。收发时钟适配见第~\ref{sec:clkrst}~节。
\end{itemize}

\figN[GMII 发送一帧的接口时序（1 字节消息，最短帧）]{f31_wgmii}{GMII 发送一帧的接口时序（1 字节消息，最短帧）。第 0 拍为前导码首字节，图中按标注的周期段压缩展示；示例载荷为 B0、SN 为 17，两网同拍起始，仅第 11 字节（接口号，A 为 \texttt{0x20}、B 为 \texttt{0x40}）和 FCS 不同；某网未启用时该网整帧保持为 0}{fig:wgmii}

\section{以太网帧与 IPv4 分片格式}\label{sec:frame}

图~\ref{fig:frame} 给出了本设计发送的 AFDX 单帧与首分片格式。目的 MAC 和源 MAC 按 ARINC 664 Part 7 的固定格式编码：目的 MAC 为 \code{03 00 00 00} 加 16 位 VLID，源 MAC 为 \code{02 00 00} 加 16 位端系统用户 ID 和 1 字节接口号（A 网 \code{0x20}、B 网 \code{0x40}）。当应用层消息总长 $M \le L_\text{max}$（$L_\text{max}\le 1471$~B）时，按不分片单帧发送；当 $L_\text{max} < M \le \MaxMessageBytes$~B 时，硬件按 IPv4 分片规则自动拆分为多个以太网帧（表~\ref{tab:frame}）：
\begin{itemize}[leftmargin=2em,itemsep=1pt]
  \item \textbf{不分片帧与首分片（42 字节协议头）：}包含 14 字节以太网头、20 字节 IPv4 头与 8 字节 UDP 头。IPv4 片偏移为 0，首分片置 \code{MF=1}（不分片帧 \code{MF=0}）；UDP 长度字段填写整条应用消息总长加 8（$M+8$），UDP 校验和字段按 RFC 768 填 \code{0x0000}。非末分片的 IPv4 数据段长度（含首片 8 字节 UDP 头）向下取整为 $\min(L_\text{max}+8, 1479)$ 的最大 8 字节整数倍（当 $L_\text{max}=1471$~B 时，首片承载 8 字节 UDP 头加 $1464$~B 应用数据，共 $1472$~B IPv4 载荷）；应用数据不足 17 字节时补零至 17 字节。
  \item \textbf{后续分片（34 字节协议头）：}仅含 14 字节以太网头与 20 字节 IPv4 头，不含 UDP 头。同一应用消息的所有分片共享同一个 16 位 IPv4 标识（Identification），片偏移字段以 8 字节为单位递增（例如 $L_\text{max}=1471$~B 时第 2 片偏移为 $1472/8=184$，承载最多 $1472$~B 应用数据；$8192$~B 消息共拆分为 6 个以太网帧，末片 \code{MF=0} 承载 $840$~B 应用数据）。后续分片若末片应用数据不足 25 字节，则硬件补零至 25 字节，使加 34 字节头、1 字节 SN 与 4 字节 FCS 后的以太网总帧长仍不小于 64 字节。
\end{itemize}

\figN[AFDX 帧格式]{f06_frame}{AFDX 单帧与首分片格式及具体字节位置。单帧应用数据长度 $1 \le n \le 1471$ 字节（后续分片无 8 字节 UDP 头，应用数据从第 34 字节起始，至多 1472 字节）；示例 $n=4$ 时，p0--p3 位于第 42--45 字节，补零 13 字节，SN 位于第 59 字节，含 FCS 的帧长为 64 字节。地址编码采用 VL 7 与 USER\_ID=\texttt{0x1234}}{fig:frame}

\begin{table}\centering\small
\caption{AFDX 以太网帧与 IPv4 分片字段（十六进制）}\label{tab:frame}
\begin{tabularx}{\linewidth}{>{\centering\arraybackslash}p{1.9cm}>{\raggedleft\arraybackslash}p{1.2cm}L}
\toprule
字节位置 & 长度 & 内容 \\
\midrule
0--3 & 4 & \code{03 00 00 00}，Part 7 规定的目的 MAC 固定部分 \\
4--5 & 2 & VLID，高字节在前 \\
6--8 & 3 & \code{02 00 00}，源 MAC 固定部分 \\
9--10 & 2 & 端系统用户 ID（寄存器 \code{USER_ID}） \\
11 & 1 & 接口号：A 网 \code{0x20}，B 网 \code{0x40} \\
12--13 & 2 & \code{08 00}（IPv4） \\
14--33 & 20 & IPv4 头：\code{45 00}；总长度（首片/单帧 $n+28$，后续分片 $n+20$）；标识（每条应用消息递增 1，同消息各分片共享）；标志与片偏移 \texttt{\{2'b00, MF, offset[12:0]\}}；TTL \code{0x01}；协议号 \code{0x11}（UDP）；逐片硬件计算的 IPv4 头校验和；源 IP；目的 IP \\
34--41 & 8/0 & 仅单帧与首分片存在 UDP 头：源端口、目的端口、总长度 $M+8$、校验和（发送填 \code{0x0000}）；后续分片无此 8 字节 \\
42（或 34）起 & $n$ & 本帧承载的应用数据（单帧 $1\le n\le 1471$，首片 $\le 1464$，后续片 $\le 1472$） \\
其后 & 补零 & 单帧/首片 $n<17$ 时补 0 至 17~B；后续分片 $n<25$ 时补 0 至 25~B \\
再后 & 1 & AFDX 序列号 SN（每发出一个以太网帧/分片独立递增） \\
最后 & 4 & FCS（CRC32，由 MAC 添加） \\
\bottomrule
\end{tabularx}
\end{table}

由此，单帧/首片含 FCS 的帧长为 $47+\max(n,17)$，后续分片含 FCS 的帧长为 $39+\max(n,25)$，物理帧长均严格落在 64--1518 字节范围内；A、B 两份帧只有第 11 字节和 FCS 不同。复位后，各 VL 首帧的 SN 均为 0，之后按 1、2、…、255、1、… 循环（分片报文的每一个分片均独立消耗一个递增 SN 并遵守所属 VL 的 BAG 间隔）。组播目的 IP 按 Part 7 图 3-38 配置为 \code{224.224.VLID[15:8].VLID[7:0]}，发送核心直接使用配置表中的目的 IP，因此单播服务（如 615A）也可共用该通路。

端系统以 VL 为独立调度与查表粒度：发送侧支持 $1\text{--}\MaxMessageBytes$~B 消息的单帧直发或 IPv4 自动分片；接收侧先对每一到达帧执行 FCS/IPv4 校验、基于 VLID 的逐帧 IC 序号检查与双网 RM 去重，不分片帧直接交付上层应用或内部服务，分片帧则进入 \ReasmContexts{} 路并发重组引擎按片偏移顺序拼装为完整 UDP 报文（若发送方提供了非零 UDP 校验和则同时验证全报文校验和）后统一交付。

CRC32 采用反射多项式 \code{0xEDB88320}，初值为全 1，计算范围覆盖从目的 MAC 到 SN 的全部字节。结果取反后按低字节在前的顺序发送。收发通路均采用组合模块 \code{afdx_crc32_d8}，以 32 位当前 CRC 值和 8 位数据为输入，输出更新后的 CRC 值。

\section{应用接口}

\subsection{顶层 2B-VLID 流式边界与核心应用接口}

应用层消息交互包含\textbf{顶层公开流式接口}（\code{AFDX_top} 的 \code{tx_data}/\code{rx_data}）与\textbf{内部核心应用接口}（\code{afdx_es} 的 \code{app_tx_*}/\code{app_rx_*}）两层：
\begin{itemize}[leftmargin=2em,itemsep=1pt]
  \item \textbf{顶层流式接口（R120-3/R124/R126）：}发送端在 \code{tx_ready=1} 时拉高 \code{tx_valid}，连续送入 2 字节大端 VLID 前缀与 $1\text{--}\MaxMessageBytes$~B 应用载荷（期间 \code{tx_valid} 不得间断），消息结束后至少拉低 \code{tx_valid} 1 拍以标志帧尾。包装器 \code{afdx_app_tx} 将消息写入 32~KiB 分体缓冲并通过 3 级流水比较器查得 7 位发送 VL 索引：若 VLID 未映射则整包丢弃并脉冲上报 \code{ev_app[0]}；若消息超过 $8192$~B 或包装缓冲/描述符队列满则回滚丢弃并上报 \code{ev_app[1]}。接收端由 \code{afdx_app_rx} 通过 $4096\times 16$~bit 映射 RAM 反查 16 位 VLID，在 \code{rx_valid} 连续高电平区间依次输出 2 字节大端 VLID 与完整载荷（相邻消息间至少空 1 拍；若 32~KiB 接收包装缓冲溢出则上报 \code{ev_app[2]}）。
  \item \textbf{核心显式消息接口（\code{afdx_es}）：}\code{app_tx_valid} 与 \code{app_tx_ready} 同拍为 1 时传送一个载荷字节，\code{app_tx_last} 标记末字节；7 位 VL 索引 \code{app_tx_vl} 和可选 UDP 端口在首字节锁存。没有服务报文争用时，\code{app_tx_ready} 经两槽入口弹性缓冲连接 \code{afdx_tx_buf}，每拍接收一字节，仅在跨 512~B 页分配（暂停 6 拍）或消息末尾入队确认（暂停约 10 拍）时拉低；顶层 32~KiB 包装缓冲将该跨页停顿对外部流接口屏蔽。三个管理服务同时有待发报文时，核心应用口轮询等待上界约 1100 拍（约 8.8~µs，第~\ref{sec:budget}~节）。
\end{itemize}

核心缓存页用尽、消息超长（$>8192$~B）或 VL 队列已满时，发送核心丢弃该消息并计数，自动回收已分配的缓存页，随后继续接收并丢弃该消息的剩余字节。

图~\ref{fig:wapp}(a) 以核心口无服务争用、长度为 $2\le n\le512$ 字节且不跨页的消息为例：接收首字节后，分配缓存页需 6 拍；接收末字节后，写入消息描述并完成入队需 10 拍。\code{app_tx_ready} 为 0 时，上游预取器保持当前字节及其有效标记，直至完成握手；图中的 C0 是下一条消息的首字节。

\subsection{接收消息的连续交付}

在核心接收接口上，一条消息（单帧或重组完成的长消息）的 $n$ 个字节在连续 $n$ 拍内送出，期间 \code{app_rx_valid} 保持为 1，\code{app_rx_last} 与末字节同拍；12 位 VL 索引、端口和来源网络在整条消息期间保持固定，见图~\ref{fig:wapp}(b)。顶层 \code{AFDX_top} 在此基础上前置 2 字节大端 VLID，以连续 $n+2$ 拍 \code{rx_valid=1} 无反压输出。

\figN[核心应用口时序]{f32_wapp}{核心应用接口（\code{afdx_es}）时序：(a) 无争用、空闲引擎及时接纳时的消息提交与握手示例；(b) 接收消息连续交付。顶层 \code{AFDX_top} 由 \code{afdx_app_tx}/\code{afdx_app_rx} 封装为带 2 字节大端 VLID 前缀的流式接口；示例中的 6 拍取页和 10 拍提交间隔不代表所有争用条件下的等待上界}{fig:wapp}

\section{配置总线}

配置总线支持寄存器和三张配置表的写入，以及寄存器和计数器的读取，时序如下（$t$ 为周期号，图~\ref{fig:wcfg}）：
\begin{itemize}
  \item \textbf{读：}$t$ 拍 \code{cfg_rd} 为 1，则 $t+2$ 拍 \code{cfg_rvalid} 为 1、\code{cfg_rdata} 有效。复位结束后可以每拍连续读。同拍读写同一寄存器时，读操作返回更新前的值。
  \item \textbf{写寄存器：}\code{cfg_wr} 与 \code{cfg_ready} 同拍为 1 时，完成写操作。
  \item \textbf{写配置表：}$t+1$ 拍起 \code{cfg_ready} 变 0，待写寄存器保持地址和数据，直到核心在某拍 $u$ 给出 \code{cfg_ok} 时写入，$u+1$ 拍 \code{cfg_ready} 恢复；不受阻时 $u=t+1$，最长等待 64 拍（第~\ref{sec:cfgok}~节）。
\end{itemize}

图~\ref{fig:wcfg} 中，A、R、T 分别对应连续寄存器读、寄存器写和配置表写；寄存器写期间 \code{cfg_ready} 保持为 1。

\figN[配置总线时序]{f33_wcfg}{配置总线的连续读、寄存器写与接收配置表写时序}{fig:wcfg}

\section{事件接口}

\code{ev_*} 端口在统计事件发生时输出单周期脉冲，与第~\ref{ch:cfg}~章的计数器一一对应，供外部逻辑监视各类事件。对于每条发送消息，发送核心均产生且仅产生 \code{accept} 或四类丢弃脉冲中的一种。对于表初始化完成后开始接收的每一帧，接收核心在对应网络上产生一个帧分类脉冲；若判定为正确帧，还会输出一个处理结果脉冲（第~\ref{sec:rxclass}~节）。

\section{内部接口}\label{sec:intif}

一级模块之间及核心内部子模块之间的接口以寄存器输出为主。ready 握手与 \code{cfg_ok} 采用组合输出；\code{afdx_svc} 的数据选择和接收分发也采用组合直通，仅包含选择与门控逻辑。内部接口主要采用以下四种约定：
\begin{itemize}
  \item \textbf{字节流握手}（\code{valid}/\code{ready}/\code{last}）：\code{valid} 与 \code{ready} 同拍为 1 时传送一个字节；附带字段在首字节锁存并保持到末字节。
  \item \textbf{请求应答}（\code{req}/\code{ack}）：请求持续有效，直到接收到应答脉冲；应答与结果在同一拍输出。
  \item \textbf{单深缓冲握手（Single-Slot Handshake）}（\code{valid}/\code{take} 或 \code{valid}/\code{ready}）：生产者置 \code{valid} 并保持字段不变，消费者接收后，生产者才能送入下一项。
  \item \textbf{写允许授权}（\code{cfg_ok}）：表的读写控制模块通过组合信号给出写允许授权。许可有效时，核心寄存外部写请求，并在下一拍写入 RAM；该拍不再发出冲突的表读请求（第~\ref{sec:cfgok}~节）。
\end{itemize}

表~\ref{tab:ifsum} 汇总各内部接口，各信号的定义见附录~\ref{app:intif}。

\begin{table}\centering\small
\caption{内部接口汇总}\label{tab:ifsum}
\begin{tabularx}{\linewidth}{>{\raggedright\arraybackslash}p{2.6cm}>{\raggedright\arraybackslash}p{3.5cm}>{\raggedright\arraybackslash}p{1.6cm}L}
\toprule
接口 & 两端 & 约定 & 关键时序 \\
\midrule
包装发送 \code{m_*} & \code{afdx_app_tx} $\rightarrow$ \code{afdx_es} & 字节流 & 经 8 信用预取器连续供给核心应用口 \\
消息 \code{m_*} & \code{afdx_svc} $\rightarrow$ \code{tx_buf} & 字节流 & 经两槽弹性缓冲解耦；跨页取页低 6 拍，消息末尾低约 10 拍 \\
入队 \code{enq_*} & \code{tx_buf} $\rightarrow$ \code{tx_engine} & 请求应答 & 携带 14 位消息长度，应答不晚于请求后 64 拍 \\
授权 \code{grant_*} & \code{tx_engine} $\rightarrow$ \code{tx_builder} & 单深缓冲握手 & 含分片偏移、标识与首/后续片标志；当前帧 \code{tx_sof} 起可接受下一授权 \\
发送起始/分片完成 & \code{tx_builder} $\rightarrow$ \code{tx_engine} & 单拍脉冲 & 回报首前导字节时刻 $t_\text{start}$ 及当前分片读末指针 \code{fragment_done_*} \\
帧字节 \code{s_*} & \code{tx_builder} $\leftrightarrow$ \code{mac_tx} & 字节流 & 经一级双缓冲流寄存器解耦；MAC 在前导码末拍接收第 0 字节 \\
帧描述 \code{desc_*} & \code{rx_parse} $\rightarrow$ \code{rx_engine} & 单深缓冲握手 & \code{f_end} 后第 3 拍给出；含分片偏移与 \code{MF}，单帧判定日程 30 拍 \\
重组预约 \code{query/reserve} & \code{rx_engine} $\leftrightarrow$ \code{rx_reassembly} & 查询应答 & 帧准入第 5/6 拍发起查询，19 拍返回槽位可接纳结果，第 28/29 拍原子提交 \\
交付队列 \code{rx_dq} & \code{rx_engine} $\rightarrow$ \code{rx_deliver} & $32\times 73$ 队列 & 队列非空后的第 6 拍输出首个 \code{frag_valid} \\
重组与交付 \code{d_*} & \code{rx_deliver} $\rightarrow$ \code{rx_reassembly} $\rightarrow$ \code{afdx_svc} & 字节流 & 单帧直通，分片帧按顺序写入槽位并在末片校验通过后连续输出 \\
服务发送 \code{svc_tx_*} & 三个服务 $\rightarrow$ \code{afdx_svc} & 字节流 & 按整条消息轮询 \\
配置表写 & \code{afdx_cfg} $\rightarrow$ 收发核心 & 写允许授权 & 等待不超过 64 拍（第~\ref{sec:cfgok}~节） \\
\bottomrule
\end{tabularx}
\end{table}

\chapter{发送通路}\label{ch:tx}

本章按数据流向说明发送核心 \code{afdx_tx_core} 的五个部分：分页缓存 \code{afdx_tx_buf} 与空闲页池 \code{afdx_tx_free}、调度引擎 \code{afdx_tx_engine}、组帧器 \code{afdx_tx_builder} 和双网 MAC \code{afdx_mac_tx}。

\section{发送核心与配置表}\label{sec:txtables}

图~\ref{fig:tx} 是发送核心的框图。应用和服务的消息经 \code{afdx_tx_buf} 分页写入载荷 RAM 并登记入队；\code{afdx_tx_engine} 按 BAG 选出下一帧并授权给 \code{afdx_tx_builder}；组帧器预读载荷、拼接帧头，逐字节交给 \code{afdx_mac_tx} 发往两网。

\figN[发送核心框图]{f07_tx}{发送核心及其配置、状态、载荷存储之间的连接}{fig:tx}

发送侧共 7 张表，全部由 \code{afdx_ram_sdp}（一写一读、同步读）例化，每张表的写口和读口各只归一个模块，见表~\ref{tab:txtables}。

\begin{table}\centering\small
\caption{发送侧配置与状态表}\label{tab:txtables}
\begin{tabularx}{\linewidth}{>{\raggedright\arraybackslash}p{2.1cm}>{\raggedleft\arraybackslash}p{1.9cm}>{\raggedright\arraybackslash}p{1.9cm}>{\raggedright\arraybackslash}p{1.9cm}L}
\toprule
表 & 深度$\times$位宽 & 写口 & 读口 & 内容 \\
\midrule
\code{tx_cfg} & 1024$\times$32 & 配置总线 & tx\_engine & 每条 VL 8 个字，地址 \{VL[6:0], 字号[2:0]\} \\
\code{tx_state} & 128$\times$63 & tx\_engine & tx\_engine & \code{timer_active}、\code{next_time}、\code{sn}、\code{q_count}、\code{q_head}、\code{q_tail} \\
\code{tx_mnext} & 256$\times$8 & tx\_engine & tx\_engine & 同一 VL 队列中的下一条消息（首页号） \\
\code{tx_desc} & 256$\times$44 & tx\_buf & tx\_engine & 消息描述：长度、\code{port_en}、源/目的端口 \\
\code{tx_pnext} & 256$\times$8 & tx\_buf & tx\_builder & 消息内的下一页 \\
\code{tx_payload} & 131072$\times$8 & tx\_buf & tx\_builder & 128~KiB 载荷，地址 \{页号[7:0], 偏移[8:0]\} \\
\code{tx_free} & 256$\times$8 & tx\_free & tx\_free & 空闲页号 FIFO \\
\bottomrule
\end{tabularx}
\end{table}

\code{tx_cfg} 中每条 VL 占 8 个字，其中 6 个配置字有效，字段见表~\ref{tab:cfgfmt}。配置字合并在同一张 1024$\times$32 的表中，相比拆分为六张浅表减少了块 RAM 占用；引擎在授权操作中逐拍读取六个配置字，步骤见表~\ref{tab:txops}。

\code{afdx_tx_core} 例化 7 张表和 5 个子模块，并持有 \code{tx_cfg} 的写口：复位后先用 1024 拍把 \code{tx_cfg} 清零，此后只在 \code{cfg_we} 与引擎的 \code{cfg_ok} 同拍为 1 时寄存写地址与数据，并在下一拍写入配置表。

\code{tx_cfg} 清零、空闲页装满和 \code{tx_state} 清零三者都完成后 \code{init_done} 置 1，引擎和分页缓存才开始工作。

\section{空闲页池}\label{sec:txfree}

\code{afdx_tx_free} 管理载荷缓存的 256 个页号；每页从分配到归还期间由一条消息独占。它是一块 256$\times$8 的 RAM 加读写指针和计数组成的 FIFO，复位后用 257 拍依次写入页号 0--255，装完后开放取页。

页池提供一个取页口和两个还页口，均带握手：
\begin{itemize}
  \item 取页口 \code{pop_req}/\code{pop_ready}/\code{pop_ack}/\code{pop_ok}/\code{pop_page} 接 \code{tx_buf}，请求应答式；
  \item 还页口 1 \code{ret1_*} 接 \code{tx_builder}，读完一页还一页；
  \item 还页口 2 \code{ret2_*} 接 \code{tx_buf}，放弃消息时退页。
\end{itemize}

取页由四步状态机完成，见表~\ref{tab:freefsm}。从请求被接受到 \code{pop_ack} 共 4 拍，其中 2 拍用于同步 RAM 读出和寄存。

\begin{table}\centering\small
\caption{空闲页池取页状态}\label{tab:freefsm}
\begin{tabularx}{\linewidth}{>{\raggedright\arraybackslash}p{2.4cm}L>{\raggedright\arraybackslash}p{2.6cm}}
\toprule
状态 & 动作 & 转移 \\
\midrule
\code{POP_IDLE} & \code{pop_ready}$=1$；接受请求时若 \code{count}$=0$ 立即应答 \code{pop_ok}$=0$，否则发出读 \code{rd_ptr} & 有页 $\rightarrow$ \code{POP_READ} \\
\code{POP_READ} & 等待同步 RAM 读出 & $\rightarrow$ \code{POP_CAPTURE} \\
\code{POP_CAPTURE} & 寄存读出的页号 & $\rightarrow$ \code{POP_REPLY} \\
\code{POP_REPLY} & 给出 \code{pop_ack}、\code{pop_ok}$=1$ 和页号 & $\rightarrow$ \code{POP_IDLE} \\
\bottomrule
\end{tabularx}
\end{table}

每个还页口先把页号锁存至各自的保持寄存器（\code{ret*_ready} 为该寄存器空），再经共用写口逐个写入 FIFO 尾部；两个寄存器均满时还页口 1 优先。取页与还页同拍发生时 \code{count} 不变。

FIFO 深度等于页数，任一页号在池中至多出现一次。写口每拍可写入一个页号，两个还页口各自最快每 2 拍还一页。

\section{分页缓存}\label{sec:txbuf}

\code{afdx_tx_buf} 从消息接口逐字节收下消息，按页写入 \code{tx_payload}，在 \code{tx_pnext} 中把同一消息的页串成链，最后写 \code{tx_desc} 并向引擎入队。载荷缓存共 256 页、每页 512 字节，地址为 \{页号[7:0], 偏移[8:0]\}。

分页缓存的接口包括消息接口 \code{m_*}（表~\ref{tab:ifm}）；取页口和还页口 2；\code{tx_payload}、\code{tx_pnext}、\code{tx_desc} 的写口；入队请求应答口（表~\ref{tab:ifgrant}）；五个统计脉冲 \code{ev_accept}、\code{ev_drop_dis}、\code{ev_drop_full}、\code{ev_drop_long}、\code{ev_drop_nopage}。

\subsection{消息接收与页分配}

图~\ref{fig:txbuf} 是其状态机，各状态的动作见表~\ref{tab:txbuffsm}。\code{m_ready} 由状态组合产生，初始化完成后只在 IDLE、RECV、DRAIN 三个状态为 1；字节传输以 \code{m_valid} 与 \code{m_ready} 同时为 1 为条件。

\figN[分页缓存状态机]{f19_txbuf}{分页缓存的消息接收、页分配、入队与异常退页流程}{fig:txbuf}

\begin{table}\centering\small
\caption{分页缓存各状态}\label{tab:txbuffsm}
\begin{tabularx}{\linewidth}{>{\raggedright\arraybackslash}p{2.1cm}L>{\raggedright\arraybackslash}p{4.0cm}}
\toprule
状态 & 动作 & 转移 \\
\midrule
IDLE & 收下首字节存入 \code{held_data}，取样 VL 与端口，长度记 1 & $\rightarrow$ ALLOC \\
ALLOC & 置 \code{pop_req} & $\rightarrow$ WAIT\_PAGE \\
WAIT\_PAGE & 等 \code{pop_ack}；得页后记入 6 个页寄存器之一，首页记为 \code{enq_head}，否则写 \code{tx_pnext}[上一页] $=$ 新页；把暂存字节写到新页偏移 0 & 无页 $\rightarrow$ RETURN；暂存字节是末字节 $\rightarrow$ DESC；否则 $\rightarrow$ RECV \\
RECV & 每拍收一字节写入当前页，偏移加 1；偏移到 511 时置 \code{boundary} & 已满 1471 字节仍有数据 $\rightarrow$ RETURN；跨页字节暂存 $\rightarrow$ ALLOC；末字节 $\rightarrow$ DESC \\
DESC & 写 \code{tx_desc}[\code{enq_head}] $=$ \{长度, \code{port_en}, 端口\} & $\rightarrow$ ENQUEUE \\
ENQUEUE & 置 \code{enq_valid} & $\rightarrow$ REPLY \\
REPLY & 等 \code{enq_ack}，记下 \code{enq_result} & 收下 $\rightarrow$ EVENT；拒绝 $\rightarrow$ RETURN \\
RETURN & 经还页口 2 逐页退回已用页 & 消息已结束 $\rightarrow$ EVENT；否则 $\rightarrow$ DRAIN \\
DRAIN & 继续收字节但不保存 & 末字节 $\rightarrow$ EVENT \\
EVENT & 按原因给出一个统计脉冲 & $\rightarrow$ IDLE \\
\bottomrule
\end{tabularx}
\end{table}

写满一页后，将下一字节暂存于 \code{held_data}，取到新页后写入其偏移 0。取页由跨页字节到达触发。最长消息占 3 页，6 个页寄存器为更小的页长参数留有余量。

\FloatBarrier
\subsection{时序与异常处理}

以首字节被收下为第 0 拍：第 1 拍发取页请求，第 2 拍被空闲页池接受，第 6 拍 \code{pop_ack} 到达并写入首字节，第 7 拍回到 RECV。每取一页 \code{m_ready} 为 0 共 6 拍。消息末字节之后，DESC、ENQUEUE 各 1 拍，引擎入队 6 拍后应答，EVENT 1 拍，\code{m_ready} 共为 0 约 10 拍。

超长（第 1472 个字节到来）、取页失败与入队被拒三种情况均转入 RETURN 态：先退回本消息已用的全部页，再继续接收并丢弃剩余字节。被丢弃消息占用的页在其统计脉冲之前全部退回。\code{reason} 记录首次出现的原因，每条消息在 EVENT 状态产生一个统计脉冲。

每条消息的统计脉冲出现在末字节被收下后 96 拍以内、下一条消息末字节之前，因此第 $k$ 个脉冲对应第 $k$ 条消息。

\code{m_ready} 连续为 0 的上界由消息末尾决定：$2 + 64 + 9 + 1 = 76$ 拍，取设计约束 96 拍（第~\ref{sec:budget}~节）。一条 1471 字节消息在应用口共占 1499 拍，短于它在线路上的 1538 拍。

\section{调度引擎}\label{sec:txeng}

\code{afdx_tx_engine} 是 \code{tx_state} 和 \code{tx_mnext} 的唯一读写控制模块，负责入队判决、按 BAG 选帧、授权组帧器和根据实际起始时刻更新下一次允许时刻。

\subsection{队列状态与操作仲裁}

\code{tx_state} 每条 VL 一项 63 位，见表~\ref{tab:txstate}。每条 VL 的消息队列是一条单向链表：\code{q_head}、\code{q_tail} 存首尾消息的首页号，\code{tx_mnext}[首页号] 指向下一条消息，因此 128 条 VL 共用 256 项链表。

\begin{table}\centering\small
\caption{\texttt{tx\_state} 字段}\label{tab:txstate}
\begin{tabularx}{\linewidth}{>{\centering\arraybackslash}p{1.6cm}>{\raggedright\arraybackslash}p{2.4cm}L}
\toprule
位 & 字段 & 含义 \\
\midrule
62 & \code{timer_active} & \code{next_time} 有效；为 0 表示可立即发送 \\
61:30 & \code{next_time} & 下一帧最早起始的 \code{now} 值 \\
29:22 & \code{sn} & 下一帧的序号，复位为 0，之后 1--255 循环 \\
21:16 & \code{q_count} & 排队消息数，不超过配置的 \code{qmax}（至多 63） \\
15:8 & \code{q_head} & 队首消息首页号 \\
7:0 & \code{q_tail} & 队尾消息首页号 \\
\bottomrule
\end{tabularx}
\end{table}

引擎一次只做一种操作，见图~\ref{fig:txfsm}。IDLE 下按“起始确认 $>$ 入队 $>$ 授权 $>$ 扫描”的优先级选择，各操作按 \code{phase} 逐拍推进，做完即回 IDLE，见表~\ref{tab:txops}。

\code{want_enq} 表示入队请求有效且未处理，并满足已有候选或扫描配额已满；\code{want_grant} 表示已有候选、授权端就绪且无待起始帧。无上述操作时，IDLE 每 2 拍检查一条 VL。

\figN[发送调度引擎状态机]{f08_txfsm}{发送调度引擎的起始确认、消息入队、帧授权与扫描流程}{fig:txfsm}

\begin{table}\centering\small
\caption{调度引擎的操作}\label{tab:txops}
\begin{tabularx}{\linewidth}{>{\raggedright\arraybackslash}p{1.5cm}>{\raggedright\arraybackslash}p{2.4cm}L>{\raggedleft\arraybackslash}p{1.1cm}}
\toprule
操作 & 触发 & 逐拍动作 & 拍数 \\
\midrule
START & 组帧器报告 $t_\text{start}$（先锁存到 \code{ack_pending}） & 0：读 \code{tx_state}[\code{pend_vl}]；1：保存；2：写 \code{timer_active}$=1$、\code{next_time}$=t_\text{start}+$\code{bag}，清 \code{pend_valid} & 3 \\
ENQ & \code{enq_valid}，且已有候选或扫描配额已满 & 0：读 \code{tx_cfg} 字 0 和 \code{tx_state}；1：保存；2：$d=$\code{now}$-$\code{next_time}；3：$d+512$；4：判决，收下则写回 \code{q_count}$+1$、新队尾，非空队列另写 \code{tx_mnext}[旧队尾]，并在条件满足时直接定为候选 & 5 \\
GRANT & 有候选且 \code{grant_ready} & 0--4：依次读配置字 1--5；1：保存字 0 和 \code{tx_state}；2：读 \code{tx_desc}、\code{tx_mnext}[队首]；4：保存描述；7：给出授权，写回 SN 前进、\code{q_count}$-1$、队首前移，记 \code{pend_vl} & 8 \\
扫描 & 无上述请求且无候选 & 每 2 拍发出一条 VL 的读，经 5 级流水判定（图~\ref{fig:scan}） & 2/VL \\
\bottomrule
\end{tabularx}
\end{table}

\FloatBarrier
\subsection{入队与扫描}

入队判决按顺序检查：VL 未启用或两网都未启用时返回结果码 1；$n$ 大于 \code{max_payload} 时返回 3；\code{q_count} 不小于 \code{qmax} 时返回 2；其余情况接受入队并返回 0。收下时若当前没有候选、该 VL 不是 \code{pend_vl} 且已在 512 拍窗口内，直接定为候选，省去一次扫描。

扫描每 2 拍发出一条 VL 的 \code{tx_state} 读，结果经图~\ref{fig:scan} 的 5 级流水判定。记 $d=$\code{now}$-$\code{next_time}，判据为：
\begin{itemize}
  \item \code{q_count}$\ne 0$、（\code{timer_active}$=0$ 或 $d+$\code{LOOKAHEAD}$\ge 0$）且不是 \code{pend_vl}：定为候选，扫描停在下一条；
  \item \code{q_count}$=0$、\code{timer_active}$=1$ 且 $d\ge 0$：写回 \code{timer_active}$=0$，防止 \code{now} 回绕约 34~s 后把旧时刻误判为将来；
  \item 判定完成时 \code{scan_commit} 指针前移至下一条 VL；任一操作插入时流水线作废，自 \code{scan_commit} 重新发出，并按更新后的状态重新判定。
\end{itemize}

\figN[扫描流水]{f21_scan}{扫描流水。每 2 拍发出一条 VL，第 $c+5$ 拍给出判定；START、ENQ 或 GRANT 插入时清空流水，从 \texttt{scan\_commit} 重新扫描}{fig:scan}

扫描从上次选中的下一条 VL 继续，128 条 VL 轮流被检查。入队操作仅在已有候选或扫描满 16 条（\code{SCAN_QUOTA}）时插入，使持续入队期间仍能推进扫描；无候选时，任一 VL 相邻两次被检查的间隔约 400 拍（第~\ref{sec:budget}~节）。

\subsection{授权状态与发送时刻}

一帧从授权到真正起始之间，该 VL 的 \code{next_time} 仍是旧值。\code{pend_vl} 保存该帧的 VL，扫描和入队将其排除在候选之外；起始确认写回新的 \code{next_time} 时清除该标记。组帧器的一个“下一帧”位置对应一个 \code{pend_vl} 寄存器。

128 条 VL 扫描一遍 256 拍，长于最短帧的 84 拍时隙；若到期后才开始找，发送时刻会随扫描位置浮动。引擎在到期前 \code{LOOKAHEAD}$=512$ 拍内选出候选并完成授权，组帧器备好长度和校验和，\code{now} 到达允许时刻 $T$ 时 MAC 恰好发出第一个前导字节。回访约 400 拍加授权和准备 22 拍小于 512，见第~\ref{sec:budget}~节和图~\ref{fig:pipe}。

\code{next_time} 以实际起始时刻 $t_\text{start}$ 加 \code{bag} 计算。线路争用使某帧晚发时，下一帧的允许时刻随之顺延，同一 VL 相邻两帧的起始间隔始终不小于 BAG。

授权时把当前 \code{sn} 交给组帧器并写回下一值：复位后每条 VL 第一帧为 0，其后 1、2、…、255、1。

\section{组帧器}\label{sec:txbld}

\code{afdx_tx_builder} 把授权的字段组成完整帧内容（目的 MAC 到 SN），在允许时刻交给 MAC，并把实际起始时刻回报给引擎。

它内部最多同时有两帧：一帧在 \code{stream_active} 状态向 MAC 送字节，另一帧在 \code{pending} 位置准备并等待。\code{grant_ready}$=\neg$\code{pending}，\code{pending} 在本帧的 \code{tx_sof} 那拍清除，因此下一帧的授权可在当前帧起始后立即进入，准备在当前帧发送期间完成。

\subsection{帧头准备与载荷发送}

组帧器的结构见图~\ref{fig:builder}：帧头走 336 位移位寄存器，载荷走预读流水和 8 项字节 FIFO，补零和 SN 由计数器产生，三者经一个选择器送出。

\figN[组帧器数据通路]{f22_builder}{组帧器数据通路。帧头与载荷分两路，在选择器汇合后交给 MAC；页归还由独立的读元数据流水触发，与载荷进入字节 FIFO 同时发生}{fig:builder}

收下授权的同拍锁存全部字段和 \code{USER_ID}，IP 标识取 16 位计数器并加 1。之后 10 拍算 IP 头校验和，每拍只做一次加法，见表~\ref{tab:prep}。

\begin{table}\centering\small
\caption{组帧器准备步骤}\label{tab:prep}
\begin{tabularx}{\linewidth}{>{\centering\arraybackslash}p{1.4cm}L}
\toprule
\code{prep} & 动作 \\
\midrule
0 & IP 总长 $=n+28$，UDP 长 $=n+8$；累加器置初值 \code{0x4500}$+$\code{0x0111}$=$\code{0x04611} \\
1--6 & 依次累加 IP 总长、标识、源 IP 高/低 16 位、目的 IP 高/低 16 位 \\
7--8 & 两次把高 4 位进位加回低 16 位 \\
9 & 取反得校验和，置 \code{prepared} \\
\bottomrule
\end{tabularx}
\end{table}

起始条件为 \code{prepared}、上一帧已送完，且 \code{grant_timed}$=0$、允许时刻已过或距允许时刻至多 4 拍。4 拍提前量补偿起始路径上的四级寄存器：\code{now} 寄存 1 拍、差值 \code{now_q}$-T$ 寄存 1 拍、\code{s_valid} 寄存 1 拍、MAC 判断起始并寄存输出 1 拍。于是线路空闲时第一个前导字节恰在 \code{now}$=T$ 那一拍出现。

起始拍把 42 字节帧头并行装入移位寄存器，A、B 两路共用，只在第 11 字节（接口号）B 网换为 \code{40}。

MAC 每取一字节，按 \code{part} 依次送出帧头 42 字节、载荷 $n$ 字节、$n<17$ 时补零到 17 字节、带 \code{s_last} 的 SN。\code{tx_sof} 那拍的 \code{now} 作为 $t_\text{start}$ 回报引擎。

起始拍起，只要 FIFO 中不足 5 字节就每拍发出一个 \code{tx_payload} 读；读请求寄存、同步 RAM 和输出寄存共 3 拍后写入 FIFO。

当前页与下一页号分开保存：起始时读 \code{tx_pnext}[首页]，每次跨页再读下一页的后继，跨页时请求连续不断。读到页末字节或消息末字节时，该页号随数据一起进入 4 项归还 FIFO，经还页口 1 退回空闲页池。

\subsection{时序与异常处理}
\begin{itemize}
  \item \textbf{FIFO 容量：}发出读时 FIFO 中至多 4 字节，另有至多 3 个读在途，最多 $4+3+1=8$ 字节，等于 FIFO 深度。
  \item \textbf{供数时序：}MAC 在起始后第 8 拍才取第 0 个字节，帧头又占 42 拍，第一个载荷字节在起始后第 50 拍才被取走，而 FIFO 在起始后第 4 拍就开始填充；此后每拍补入一字节、取出一字节。下一页号在跨页前 500 拍以上已读出。
  \item \textbf{准备时间：}下一帧准备需 11 拍，可与当前帧的发送重叠；最短帧时隙为 84 拍。
\end{itemize}
从消息末字节被收下到 \code{tx_sof} 的名义值为 31 拍，上界 100 拍，逐拍分解见第~\ref{sec:budget}~节图~\ref{fig:pipe}。

MAC 报 \code{tx_underrun} 时本帧记为坏帧，不产生 \code{ev_tx_a}/\code{ev_tx_b}；其余帧在 \code{tx_eof} 时按本帧启用的网络各给一个脉冲。

\Needspace{14\baselineskip}
\section{双网发送 MAC}\label{sec:mactx}

\code{afdx_mac_tx} 用一个状态机驱动 A、B 两套 CRC 和输出寄存器，两网同拍输出。发送节拍由 MAC 控制，上游按其握手信号提供字节。图~\ref{fig:macfsm}(a) 是其状态机，(b) 是第~\ref{sec:macrx}~节的接收 MAC。

eLinx 生成的三速以太网 IP 带用户侧 FIFO、保留 FCS，且每网一个实例；相同器件下的资源对照见测试报告。

发送状态机仅在 \code{byte_en} 有效时推进。接收状态机使用寄存后的有效标记和数据；PREAMBLE 或 \code{BAD_PRE} 因有效标记撤销而退出时，报告前导错误。

各状态的动作见表~\ref{tab:macfsm}，发送一帧的字节时间见表~\ref{tab:macbyte}，$L$ 为取走的字节数。起始那一拍取样 \code{s_en_a}、\code{s_en_b}；未启用的网络在整帧期间 \code{tx_valid} 与 \code{tx_data} 保持为 0；两网均未启用时状态机仍按相同节拍推进，使各帧时序与网络使能无关。上游一直有帧等候时，帧间隔恰为 12 拍。

数据阶段若上游未备好字节（欠载），MAC 给出 \code{tx_underrun} 脉冲，发出 4 字节不取反的错误 FCS，然后继续接收并丢弃上游这一帧的剩余字节。

复位期间线路输出全为 0，复位结束后先等满 12 拍帧间隔再允许首帧。

\figN[双网 MAC 状态机]{f10_macfsm}{双网 MAC 状态机：(a) 发送及欠载处理；(b) 接收及前导检查}{fig:macfsm}

\begin{table}\centering\small
\caption{发送 MAC 各状态}\label{tab:macfsm}
\begin{tabularx}{\linewidth}{>{\raggedright\arraybackslash}p{2.1cm}L>{\raggedright\arraybackslash}p{3.4cm}}
\toprule
状态 & 动作 & 转移 \\
\midrule
GAP & 输出 0；\code{gap_left} 递减 & \code{gap_left}$=0$ 且 \code{s_valid} $\rightarrow$ PREAMBLE，同拍取样两网使能、给 \code{tx_sof} \\
PREAMBLE & 输出 7 个 \code{55} 和 \code{D5}；末拍 \code{s_ready}$=1$ 取走第 0 字节 & 取到字节 $\rightarrow$ DATA；未备好 $\rightarrow$ FCS（欠载） \\
DATA & 每拍输出一字节并更新两网 CRC & 末字节已出 $\rightarrow$ FCS；未备好 $\rightarrow$ FCS（欠载） \\
FCS & 输出 4 字节 FCS，低字节先发；第 4 拍后撤 \code{valid}、给 \code{tx_eof}，\code{gap_left}$=11$ & 正常 $\rightarrow$ GAP；欠载 $\rightarrow$ DRAIN \\
DRAIN & \code{s_ready}$=1$，继续接收并丢弃上游剩余字节 & 收到 \code{s_last} $\rightarrow$ GAP \\
\bottomrule
\end{tabularx}
\end{table}

\begin{table}\centering\small
\caption{MAC 发送的字节时间}\label{tab:macbyte}
\begin{tabularx}{\linewidth}{>{\raggedright\arraybackslash}p{3.0cm}>{\raggedright\arraybackslash}p{3.6cm}L}
\toprule
字节时间 & 线上内容 & 本拍动作 \\
\midrule
起始前一拍 & 空闲 & 判断起始，取样两网使能 \\
0--6 & \code{55} & \\
7 & \code{D5} & 取走第 0 个字节 \\
8 至 $8+L-1$ & 第 0 至 $L-1$ 个字节 & 取走下一字节，取走带 \code{s_last} 的字节后停止 \\
其后 4 拍 & FCS & A、B 两网各自更新 CRC32 \\
再后至少 12 拍 & 空闲 & 第 12 拍可判断下一帧起始 \\
\bottomrule
\end{tabularx}
\end{table}

\chapter{接收通路}\label{ch:rx}

接收核心 \code{afdx_rx_core} 为 A、B 两路网络各配置一套接收 MAC 和帧头解析器，两网共用同一接收引擎与交付模块。MAC 在接收过程中逐字节校验 CRC，解析器同步检查帧头并将载荷预写入本网的环形缓冲。对于通过解析器检查的帧，接收引擎根据帧描述依次执行 VL 查表、完整性检查和冗余管理，再通知解析器提交数据或回退写指针；已提交的消息最终由交付模块逐字节送往应用口。

\section{接收核心与配置表}\label{sec:rxtables}

\figN[接收核心框图]{f11_rx}{接收核心及配置、状态和载荷存储之间的连接}{fig:rx}

接收侧的片上表见表~\ref{tab:rxtables}。\code{rx_map} 以 16 位 VLID 直接寻址，任意 4096 个 VLID 均可映射到接收 VL 索引 \code{idx}；映射表定义为 65536$\times$13 的存储阵列，读出数据直接锁存至一级寄存器，查表延迟固定为 2 拍。

\begin{table}\centering\small
\caption{接收侧配置与状态表}\label{tab:rxtables}
\begin{tabularx}{\linewidth}{>{\raggedright\arraybackslash}p{2.0cm}>{\raggedleft\arraybackslash}p{1.9cm}>{\raggedright\arraybackslash}p{1.9cm}>{\raggedright\arraybackslash}p{1.9cm}L}
\toprule
表 & 深度$\times$位宽 & 写口 & 读口 & 内容 \\
\midrule
\code{rx_map} & 65536$\times$13 & 配置总线 & \code{rx_engine} & \code{valid}[12]、\code{idx}[11:0] \\
\code{rx_cfg} & 4096$\times$32 & 配置总线 & \code{rx_engine} & \code{en}、\code{net_a}、\code{net_b}、\code{ic_en}、\code{rm_en}、\code{skew_max}[16:5]、\code{max_payload}[27:17] \\
\code{rx_state} & 4096$\times$44 & \code{rx_engine} & \code{rx_engine} & 每网 PSN 及有效位；\code{rm_last}、\code{rm_time} 及有效位 \\
\code{rx_ring_a/b} & 4096$\times$8 & 本网 \code{rx_parse} & \code{rx_deliver} & 每网一个环形缓冲，按到达顺序存载荷 \\
\code{rx_dq} & 32 项寄存器 & \code{rx_engine} & \code{rx_deliver} & 待交付消息的来源网络、起始地址、长度、\code{idx}、UDP 端口 \\
\bottomrule
\end{tabularx}
\end{table}

\code{rx_cfg} 中的 \code{skew_max} 以 \code{now_c} 的计时单位（32.768~µs）计。\code{afdx_rx_core} 例化 A、B 两路接收 MAC 与解析器、共用接收引擎、交付模块及 5 张表，其中 \code{rx_map}、\code{rx_cfg} 的写口由核心统一管控。复位后，核心用 65536 拍逐项清零 \code{rx_map}（前 4096 拍同步清零 \code{rx_cfg}），引擎另行清零 \code{rx_state}；全部完成后置位 \code{init_done}。运行期间仅在 \code{cfg_ok} 为 1 时锁存配置写请求并于下一拍写入配置表。核心将各模块的统计脉冲按位相或，寄存一拍后输出。

\section{接收 MAC}\label{sec:macrx}

\code{afdx_mac_rx} 把 GMII 输入寄存一拍后，在 \code{rx_valid} 为 1 的每一拍处理一个字节，状态机见图~\ref{fig:macfsm}(b)。
\begin{itemize}
  \item \textbf{帧起始定界（SFD 检测）：}前导字节 \code{0x55} 的数量不限，检测到 \code{0xD5} 后以紧随的下一字节作为帧首字节。若收到其他值，或在检测到 \code{0xD5} 之前 \code{rx_valid} 已回到 0，则判为前导码错误。复位后若 \code{rx_valid} 已为 1，MAC 先等待其回到 0，再从下一帧开始接收。
  \item \textbf{剥离 FCS：}每接收一字节均更新 CRC 和帧长；输出比输入晚 4 字节，收到第 $k+4$ 字节时输出第 $k$ 字节，最后 4 字节（FCS）不送后级。
  \item \textbf{帧结束：}寄存后的 \code{rx_valid} 首次为 0 时，MAC 判定帧接收结束，输出 \code{f_end} 和 \code{f_err}。错误判定依次检查前导码错误、帧长不在 64--1518 范围内、CRC 余式不等于 \code{0xDEBB20E3}，均无异常则判为正确帧。MAC 向后级最多输出 1514 字节；若帧长超限，则按帧长错误结束本帧。
  \item \textbf{帧间：}\code{f_end} 之后立即回到空闲，可分辨间隔为 1 个低电平周期的两帧。
\end{itemize}
以 GMII 上 FCS 末字节到达时刻为第 0 拍：第 1 拍该字节锁存至输入寄存器，第 2 拍寄存后的 \code{rx_valid} 拉低，第 3 拍输出 \code{f_end}、\code{f_err} 和 \code{f_len}。剥离 FCS 后的最后一个有效字节在第 2 拍从 \code{m_data} 输出，比第 3 拍的 \code{f_end} 早 1 拍，解析器据此完成校验和累加。

\section{帧头解析与预写入缓存}

\code{afdx_rx_parse} 由字节位置计数器 \code{position} 驱动，随输入字节流式检查帧头；另设两个辅助状态机分别负责长度计算与帧尾分类，结构见图~\ref{fig:parse}。

\subsection{流式检查与长度计算}

\figN[帧头解析]{f23_parse}{帧头解析：(a) 按字节位置检查并预写入载荷；(b) 收到 UDP 长度后的两拍长度计算；(c) 帧尾的两拍分类}{fig:parse}

帧头必须满足以下条件，任一条件不满足均判为“帧头错”：
\begin{itemize}
  \item 第 0--3 字节为 \code{03 00 00 00}；第 12--13 字节为 \code{08 00}；第 14 字节为 \code{0x45}；第 23 字节为 \code{0x11}；
  \item IP 不分片（MF 为 0、片偏移为 0）；IP 头 10 个 16 位字的反码和为 \code{0xFFFF}；
  \item 记 $n$ 为 UDP 长度减 8，要求 $1\le n\le 1471$，IP 总长等于 $n+28$，且含 FCS 的帧长等于 $47+\max(n,17)$；
  \item UDP 校验和非 0 时，伪头、UDP 头与 $n$ 字节数据的反码和为 \code{0xFFFF}。
\end{itemize}
补零字节的值不检查。解析器同时记录 VLID、源/目的端口、$n$ 和 SN。

在第 39 字节收到 UDP 长度低字节后，\code{length_phase} 控制长度计算分两拍完成：先在 LEN\_SIZES 步计算 $n$、检查长度是否合法，并求出期望 IP 总长；再在 LEN\_LIMITS 步计算载荷末位置、SN 位置、期望帧长和环形缓冲所需空间。第 42 字节到来时，上述结果均已就绪。

IP 校验和与 UDP 校验和分别使用 20 位和 26 位累加器，每两字节累加一次：前者累加 IP 头的 10 个 16 位字，后者累加 UDP 伪头、UDP 头及载荷。进位折叠采用两级寄存器流水，与累加过程并行；到帧尾后第 1 拍（FIN\_SETTLE），UDP 校验和结果最终稳定。

\subsection{载荷预写入与提交}

从第 42 字节开始，解析器将载荷边接收边写入本网环形缓冲，待最终判定后再决定是否提交。环形缓冲采用三个指针（图~\ref{fig:ring}）：读指针 \code{rd} 由交付模块维护；已提交写指针 \code{wr_commit} 和预写入指针 \code{wr_trial} 由本网解析器维护。

计算剩余空间时以 \code{wr_commit} 为准，不计入正在接收的帧。第 42 字节到来时 $n$ 已知，若剩余空间小于 $n$，解析器便不写入本帧载荷，并将其记为溢出。

\figN[接收环形缓冲的三个指针]{f12_ring}{接收环形缓冲的三个指针。帧通过全部检查后，解析器将 \texttt{wr\_commit} 更新为 \texttt{wr\_trial}；否则将 \texttt{wr\_trial} 回退至 \texttt{wr\_commit}}{fig:ring}

\subsection{帧尾分类与描述交接}

每次 \code{f_end} 有效时，解析器均对该帧进行一次分类。\code{finish_phase} 在 \code{f_end} 后第 2 拍进入 FIN\_REPORT 并作出判定，类别脉冲和帧描述均经寄存后输出，在第 3 拍有效。分类时，若 \code{f_err} 表示 FCS 错误或帧长错误（含前导码错误），则直接采用该结果；否则依次判定帧头错误、溢出或正确。

解析器在帧尾完成错误帧的分类与计数，并将预写入指针 \code{wr_trial} 回退至已提交写指针 \code{wr_commit}；正确帧则生成帧描述字交由接收引擎处理（表~\ref{tab:ifdesc}）。两网在同一拍结束的正确帧，其帧描述也在同一拍给出。

同一网相邻两帧的 \code{f_end} 至少相隔 73 拍（最短帧 64 字节加前导码 8 字节和 1 拍间隔）。引擎处理一份描述至多需要 13 拍，即使两网各有一份描述，也只需 26 拍即可处理完毕，短于同一网相邻描述的到达间隔。若数据帧的首个前导字节在清表完成前到达，解析器将忽略整帧，不生成帧描述，也不计入统计。

\section{完整性检查与冗余管理}

\code{afdx_rx_engine} 是 \code{rx_state} 的唯一读写控制模块，一次处理一份帧描述或执行一步冗余记录过期清理，两网描述同拍到达时先处理 A 网。

\subsection{帧描述处理流程}

帧描述的处理流程见图~\ref{fig:decide}。引擎先查询 \code{rx_map}，再读取 \code{rx_cfg} 和 \code{rx_state}，随后依次执行配置检查、完整性检查、冗余管理和交付入队。任一步未通过，均按对应原因计数，并通知解析器回退预写入指针。

\figN[接收引擎对一份帧描述的判定流程]{f13_decide}{接收引擎对一份帧描述的配置检查、完整性检查、冗余管理与交付判定}{fig:decide}

引擎有 INIT、IDLE、FRAME、AGE 四个操作态，FRAME 按 \code{phase} 0--11 逐拍推进，每拍至多一级比较，见图~\ref{fig:rxeng} 和表~\ref{tab:rxphase}。对于未通过检查的帧，引擎提前跳转至写回或报告步骤，因此一份描述的处理时间为 5--13 拍（含 IDLE 取描述的 1 拍）。图中的 \code{desc} 为 \code{desc_a_valid} 或 \code{desc_b_valid}；AGE 的任一步收到新描述时均返回 IDLE。

\figN[接收引擎状态机]{f24_rxeng}{接收引擎状态机：FRAME 处理帧描述，AGE 在空闲时检查冗余记录}{fig:rxeng}

\begin{table}\centering\small
\caption{接收引擎 FRAME 各步}\label{tab:rxphase}
\begin{tabularx}{\linewidth}{>{\centering\arraybackslash}p{1.2cm}L}
\toprule
\code{phase} & 动作 \\
\midrule
（IDLE） & 取走帧描述（\code{take}），以 VLID 读 \code{rx_map} \\
0--1 & 等待同步读出，保存映射 \\
2 & 映射无效 $\rightarrow$ 11（未知 VL）；否则以 \code{idx} 读 \code{rx_cfg} 和 \code{rx_state} \\
3--4 & 等待读出，保存配置和状态 \\
5 & VL 或本网未启用、$n>$\code{max_payload} $\rightarrow$ 11（配置不允许）；同时计算 PSN 的一步后继序号、到达时间差、SN 与 \code{rm_last} 的原始距离 \\
6 & 时间差取绝对值，计算 PSN 的两步后继序号，原始距离折算到 1--255 环上 \\
7 & 完整性检查，更新本网 PSN；不通过 $\rightarrow$ 10 \\
8 & 冗余管理判定 \code{deliver_allowed} \\
9 & 执行冗余去重判定；若 \code{rx_dq} 已满则记为溢出，否则写入交付队列并更新 \code{rm_last}、\code{rm_time} \\
10 & 状态有变化则写回 \code{rx_state} \\
11 & 给出 \code{resolve} 与 \code{commit}，按处理结果产生一个统计脉冲，返回 IDLE \\
\bottomrule
\end{tabularx}
\end{table}

\FloatBarrier
\subsection{配置与完整性判定}

\code{rx_map} 映射无效时，引擎将该帧计为“未知 VL”；VL 未启用、本网未启用或 $n$ 大于 \code{max_payload} 时，计为“配置不允许”。

完整性检查（\code{ic_en}$=1$ 时）每网独立进行。记本网上一帧的序号为 PSN，满足以下任一条件即可通过检查：清表后本网首次收到该 VL；SN 为 0（对端复位）；SN 为 PSN 在 1--255 环上递增 1 步或 2 步的后继序号，后者允许丢一帧。PSN 为 0 时，后继序号分别为 1 和 2。

启用完整性检查时，通过映射与配置检查的帧均以当前 SN 更新本网 PSN，包括序号检查失败的帧；后续序号连续时可恢复同步。未通过序号检查的帧计为“完整性检查丢弃”。

引擎先对每份副本独立进行完整性检查，通过后再进行冗余管理，符合 ARINC 664 P7 的处理顺序。

\subsection{冗余窗口与过期清理}

冗余管理（\code{rm_en}$=1$ 时）按处理顺序交付首个有效副本；两网描述同拍待处理时先处理 A 网。副本到达间隔在 \code{skew_max} 内时，同一消息只交付一次，判据如下：
\begin{itemize}
  \item 尚无记录，或到达时刻与 \code{rm_time} 之差超过 \code{skew_max}：交付；当两份副本的处理顺序与到达顺序相反时（如引擎忙期间 B 网先完成解析、A 网后完成解析，引擎转入空闲时按 A 网优先仲裁先处理了 A 网），对 16 位时间差取绝对值。
  \item 否则 SN 等于 \code{rm_last}：另一网的重复帧，丢弃并计“冗余管理丢弃”。
  \item 否则 SN 为 0 或 \code{rm_last} 为 0：交付。
  \item 否则在 1--255 环上，SN 领先 \code{rm_last} 1--127 步则交付，其余丢弃。
\end{itemize}
对于判定可交付的帧，若 \code{rx_dq} 已满，则计为“溢出”并回退预写入指针；否则提交环形缓冲中的载荷，将帧描述写入 \code{rx_dq}，并更新 \code{rm_last} 和 \code{rm_time}。两网描述同拍到达时，从描述有效到引擎完成判定，后处理的一份也不超过 27 拍（设计约束取 64 拍）。

\code{now_c} 每变化一次，引擎在空闲时通过 AGE 状态的 4 个步骤检查下一条 VL；有帧描述到来时立即返回 IDLE 处理帧，之后重试该项过期清理检查。判据为：若记录已超过 \code{0x4000} 个计时单位（约 0.537~s），清除 \code{rm_time_valid}。

帧描述处理优先于过期清理，清理被打断后会重新检查同一项。4096 项约需 0.134~s 检查一遍，记录在约 0.671~s 内失效，早于 16 位时间差的半圈时间 1.074~s。

\section{交付}\label{sec:rxdeliver}

\code{afdx_rx_deliver} 按入队顺序从 \code{rx_dq} 取出帧描述，再从对应环形缓冲每拍读出并输出一个字节，同时释放已读空间，结构见图~\ref{fig:deliver}。

\figN[交付模块]{f25_deliver}{交付模块的队头选择、环形缓冲读取与字节输出流水}{fig:deliver}

图~\ref{fig:deliver} 中 C1--C6 依次为组内选择寄存、队头寄存、读请求寄存、同步 RAM 读出、两网数据寄存和输出寄存六个采样边界。队头选择状态机按 \code{HEAD_IDLE}（组内 8 选 1，C1）、\code{HEAD_SELECT}（组间 4 选 1 锁存至 \code{dq_data} 并给出 \code{dq_pop}，C2）和 \code{HEAD_RELEASE}（等待出队握手完成）推进；队头进入 \code{dq_data} 后，读控制在下一拍发起首次读取，随后每拍将地址加 1，直到 \code{left} 减为 0 时停止。

读请求、同步 RAM、数据寄存三级流水由 \code{valid_pipe}、\code{last_pipe} 对齐，第一个 \code{d_valid} 在 \code{dq_count} 变为非零后的第 6 拍有效。\code{d_last} 有效的当拍，模块重新允许取队头。

一条消息的 $n$ 个字节在连续 $n$ 拍内送出，\code{d_last} 与末字节同拍；VL 索引、端口和来源网络在整条消息期间不变。交付口不带反压。

线路空闲时，从帧尾到首个 \code{d_valid} 的典型延迟为 24 拍，设计上界为 200 拍（图~\ref{fig:pipe}）。若前一消息尚未交付完，当前消息需等待；取队头与读流水共约 6 拍，可在前一消息末字节之后的 16 拍内开始交付。32 项描述队列用于暂存长载荷交付期间到达的短帧描述（第~\ref{sec:budget}~节）。

\FloatBarrier
\Needspace{22\baselineskip}
\section{统计分类}\label{sec:rxclass}

对于清表完成后开始接收的每一帧（一段 \code{rx_valid} 连续为 1 的字节流），解析器均在帧尾离开线路后的 32 拍内，在对应网络上输出且仅输出一个类别脉冲。对于分类为正确的帧，接收引擎还输出且仅输出一个处理结果脉冲；该脉冲在引擎完成判定后的 8 拍内有效，且不晚于帧尾后的 150 拍。每网的 \code{ev} 位定义见表~\ref{tab:rxev}。

\begin{table}\centering\small
\caption{接收统计脉冲 \texttt{ev\_a}、\texttt{ev\_b}}\label{tab:rxev}
\begin{tabularx}{\linewidth}{>{\centering\arraybackslash}p{1.0cm}>{\raggedright\arraybackslash}p{3.0cm}>{\raggedright\arraybackslash}p{2.2cm}L}
\toprule
位 & 名称 & 产生模块 & 含义 \\
\midrule
0 & 正确帧 & \code{rx_parse} & 通过 MAC 和帧头检查 \\
1 & FCS 错误 & \code{rx_parse} & CRC 校验不符 \\
2 & 帧长错误 & \code{rx_parse} & 帧长越界、超长或前导码错误 \\
3 & 帧头错误 & \code{rx_parse} & 固定字段、IP/UDP 长度或校验和错误 \\
4 & 未知 VL & \code{rx_engine} & \code{rx_map} 无效 \\
5 & 配置不允许 & \code{rx_engine} & VL 或本网未启用，或超过 \code{max_payload} \\
6 & 完整性检查丢弃 & \code{rx_engine} & SN 不连续 \\
7 & 溢出 & 两处相或 & 环形缓冲剩余空间不足，或 \code{rx_dq} 已满 \\
\bottomrule
\end{tabularx}
\end{table}

位 7 由两处溢出脉冲相或产生：解析器报告环形缓冲剩余空间不足，引擎报告 \code{rx_dq} 已满。同一网上，相邻帧尾至少相隔 73 拍，引擎判定至多 27 拍，解析器的第 $k$ 帧类别脉冲与引擎的第 $k-1$ 帧结果脉冲分时到达，可相或合并。

两网合计另有 \code{ev_rm_drop}（冗余管理丢弃）和 \code{ev_deliver}（交付）。这些脉冲与第~\ref{ch:cfg}~章的计数器一一对应，记录每类接收结果。

\addtocontents{toc}{\protect\Needspace{6\baselineskip}}
\chapter{服务功能}\label{ch:svc}

RTC、SNMP 和 615A 均以 UDP 报文工作，经服务通道 \code{afdx_svc} 与应用共享收发核心。服务报文按配置的 VL 参与 BAG 调度、完整性检查和冗余管理。


\section{服务通道}\label{sec:svcch}

\code{afdx_svc} 由组合选择和五个状态寄存器组成：发送侧为 \code{busy}、\code{owner}、\code{turn}，接收侧为 \code{rx_active}、\code{rx_claimed}，结构见图~\ref{fig:svc}。

\figN[服务通道]{f26_svc}{服务通道：(a) 应用与服务的整消息发送仲裁；(b) 按 UDP 目的端口进行接收分发}{fig:svc}

应用口与三个服务发送口按整条消息轮询，次序为应用、RTC、SNMP、615A。选中的来源独占消息通路，直至末字节握手，随后从下一来源开始选择。服务发送口均空闲时，应用的数据和 ready 组合直通发送核心。发送仲裁器由三个状态寄存器构成：\code{busy} 指示当前消息传输忙，\code{owner} 锁存当前来源，\code{turn} 指示下一轮轮询起点；末字节握手时 \code{busy} 清零、\code{turn} 指向当前来源的下一个。由此任一来源在提出请求后，至多等其余三个来源各传完一条消息。

接收核心的交付口同时广播给应用和三个服务。各服务按 UDP 目的端口给出认领信号，首字节的认领结果保持到末字节：被认领的消息只送服务，应用口的 valid 和 last 为 0；其余消息原样送给应用。三个服务使用不同的 UDP 端口，且都不对交付口施加反压。认领信号由各服务按 \code{d_dport} 组合产生，首字节当拍直接决定应用口的 \code{valid}，之后由寄存的 \code{rx_claimed} 保持。\code{rx_active} 选择当前认领或寄存的认领结果，其反相信号与 \code{d_valid} 相与形成应用口有效信号。

\section{RTC 时钟同步}\label{sec:rtc}

\code{afdx_rtc} 维护 32 位秒加 32 位纳秒的实时时钟，每拍前进 8~ns，纳秒满 $10^9$ 向秒进位；秒的起点为 2000-01-01 00:00:00 UTC，线上 NTP 秒为其加 3155673600。同步采用 RFC 4330 的 SNTPv4 主从方式，交互见图~\ref{fig:rtc}。

\figN{f14_rtc}{SNTP 主从对时交互及四个时间戳的采样位置}{fig:rtc}

\subsection{时间戳采样与收发时序}

\code{afdx_rtc} 包含本地时钟、一个 48 字节接收移位寄存器、发送字节状态机、校时计算状态机和三个 \code{afdx_ntp_fraction} 换算器（发送时间戳编码、接收或参考时间戳编码、回应时间戳解码），两个状态机见图~\ref{fig:rtcfsm}。

从端每隔 \code{PERIOD} 发一个请求，记下发送时刻 $t_1$；主端在收到请求时记 $t_2$，回应发送时记 $t_3$，并把请求的 Transmit 时间戳逐位复制到回应的 Originate 字段；从端收到回应时记 $t_4$。从端只接受 Originate 与本次请求逐位相等的回应，随后计算
\begin{equation}
\text{offset} = \frac{(t_2-t_1)+(t_3-t_4)}{2}
\end{equation}
并加到当前时钟。一个周期内未收到有效回应时，下一周期以新的 Transmit 时间戳重发，直至校时成功。

发送时间戳取消息首字节与发送核心的握手时刻，接收时间戳取核心交付首字节的时刻。时间戳在握手后开始换算，并在对应字段发出前完成。

\figN[RTC 的两个状态机]{f27_rtc}{RTC 状态机：(a) 报文发送；(b) 回应时间戳解码与时钟校正}{fig:rtcfsm}

主端有待回应请求或从端倒计时归零时启动发送。无反压时每字节用 4 拍：TX\_GROUPS 先在两组字段中各选一个 32 位字，TX\_WORD 选字，TX\_BYTE 选字节并置 \code{tx_valid}，TX\_HANDSHAKE 等 \code{tx_ready}。一个 48 字节报文至少 192 拍。Transmit 字段位于第 40--47 字节，距首字节握手至少 160 拍，换算耗时 65 拍。主端的 Receive 字段 $t_2$ 在请求首字节到达时开始换算，回应在整条请求收完后开始发送。

\subsection{时间换算与偏差校正}

\code{afdx_ntp_fraction} 负责 NTP 的 32 位二进制小数与本地纳秒值之间的换算。纳秒编码为 NTP 小数时，用 32 轮、每轮 2 拍的逐位比较减法计算 $\lfloor \text{ns}\times 2^{32}/10^9 \rfloor$；反方向用逐位加与右移并在末轮四舍五入。每拍最多一级 32 位加减，整纳秒时间戳往返换算得到原值。

秒差和纳秒差分开处理，各拍至多一级 32 位运算，步骤见图~\ref{fig:rtcfsm}(b)：先依次解码回应中的 $t_2$、$t_3$ 小数部分（各 65 拍），再用 9 拍计算出两段差、求和、除以 2，并把奇数秒的半秒折入纳秒、处理借位进位。校正量在 ADD\_OFFSET 步加至当前时钟，至 APPLY 步写回共经过 4 拍，故在 LATENCY 步补偿 32~ns；一次校时约 140 拍，期间时钟正常推进。

本端正在等待回应且发送状态机空闲时，以版本 VN=4、模式 mode=4 及 Originate 与本次请求 Transmit 逐位相等作为校时计算的启动条件，服务器 stratum 与全零时间戳不参与判定。迟到的旧回应因 Originate 不符而被丢弃。计算期间或 \code{rx_done} 有效的周期不发起新请求。

\subsection{报文字段与寄存器}

请求与回应均为 48 字节，所有多字节字段高字节先发，见表~\ref{tab:sntp}。双方 UDP 源、目的端口均为 \code{UDP_PORT}（默认 123）。

\begin{table}\centering\small
\caption{SNTPv4 报文字段}\label{tab:sntp}
\begin{tabularx}{\linewidth}{>{\centering\arraybackslash}p{1.4cm}LL}
\toprule
字节 & 请求（从端） & 回应（主端） \\
\midrule
0 & \code{0x23}：LI=0、VN=4、mode=3 & \code{0x24}：LI=0、VN=4、mode=4 \\
1 & stratum=0 & stratum=1 \\
2--3 & Poll=4、Precision=$-26$ & 同左（约 14.9~ns） \\
4--11 & Root Delay、Root Dispersion 均为 0 & 同左 \\
12--15 & 0 & Reference ID=\code{LOCL} \\
16--23 & 0 & 最近一次配置装入的时刻 \\
24--31 & Originate=0 & 请求 Transmit 的逐位复制 \\
32--39 & Receive=0 & $t_2$ \\
40--47 & Transmit=$t_1$ & $t_3$ \\
\bottomrule
\end{tabularx}
\end{table}

RTC 寄存器位于 \code{0x80}--\code{0x88}，见附录表~\ref{tab:regs}。写时间先写 \code{SEC}、\code{NS} 两个暂存字，再以 \code{CTRL} 的装入位一次生效；读 \code{SEC} 时同时锁存纳秒值，随后读 \code{NS} 得到同一时刻的值。偏差以 \code{OFFSET_SEC}$\times 10^9 +$\code{OFFSET_NS} 纳秒表示。


\section{SNMP 网络管理}\label{sec:snmp}

\code{afdx_snmp} 是服务通道上的 SNMP 代理，认领 UDP 目的端口 161，回应的目的端口为请求的源端口，经配置的发送 VL 送出。图~\ref{fig:snmp} 是其处理流程。

\figN{f15_snmp}{SNMP 代理处理流程}{fig:snmp}

\subsection{协议对象与操作}

MIB 包含 35 个只读对象与 1 个读写对象（\code{USER_ID}），共 36 个。对象分为 sysDescr、sysUpTime 两个系统对象，以及演示用私有子树 \code{1.3.6.1.4.1.424242.3} 下的 34 个对象，覆盖版本号、\code{USER_ID}、RTC 秒和同步次数、全部 27 个收发统计计数器、SNMP 请求数、community 错误数和 Trap 发送数。对象值直接取自寄存器。

\Needspace{4\baselineskip}
代理支持以下协议版本和操作：
\begin{itemize}
  \item \textbf{版本：}SNMPv1（version=0）和 v2c（version=1）。
  \item \textbf{Get / GetNext：}Get 精确匹配标量 OID；GetNext 按子标识符逐个比较，返回严格大于请求 OID 的第一个对象，可从 \code{1.3} 遍历全部对象。对象或实例未命中时，v1 回应 noSuchName，v2c Get 统一回应 noSuchObject；v2c GetNext 遍历至末尾时回应 endOfMibView。
  \item \textbf{Set：}\code{USER_ID} 为读写配置对象，类型为 Gauge32，取值 0--65535。代理先检查 OID、类型和取值范围，再经配置总线写入并读回，以读回值组装回应。只读对象、类型不符和越界分别按 v1/v2c 的错误码回应，配置不变。
  \item \textbf{Trap：}发送 SNMPv2-Trap（PDU \code{0xA7}）到 UDP 端口 162，含 sysUpTime.0 和 snmpTrapOID.0 两个变量绑定。\code{TRAP_ENABLE} 已置位且 \code{ENABLE} 从 0 变为 1 时发 coldStart；community 不符且 \code{TRAP_ENABLE} 已置位时发 authenticationFailure。
\end{itemize}
每个请求按单变量绑定处理；community 不符时不回应并计数 \code{BAD_COMMUNITY}。

\subsection{请求解析与回应组装}

代理的主状态机有 44 个状态，按功能分为 7 组，见图~\ref{fig:snmpfsm}。请求先完整写入 256$\times$8 的请求 RAM，解析与查找通过公共读取状态序列访问 RAM：原状态分支将返回位置写入 \code{return_state}，读取序列用 4 拍完成越界检查、发出读请求、等待同步读出和锁存数据，随后转回原状态分支。\code{Fetch} 状态组由解析、\code{Lookup}、\code{Value} 和 \code{Encode} 共用；每个解析状态只处理一个已就绪的字节，单拍逻辑只有一级比较。Set 未命中时经 Value 组的 \code{SET_ERROR} 编码错误回应，异常处理完成后返回接收状态。

\figN[SNMP 代理状态机（按功能分组）]{f28_snmpfsm}{SNMP 代理的请求解析、对象访问与回应组装状态机}{fig:snmpfsm}

代理按 BER 格式依次解析 SEQUENCE、version、community 和 PDU；PDU 内依次是 request-id、error-status、error-index 和变量绑定表。长度字段支持短形式和 \code{0x81}、\code{0x82} 两种定长形式。回应在 256$\times$8 的回应 RAM 中逐字节组装，原样保留 request-id 和 community，再经服务发送口输出。

对象查找按对象表顺序逐个尝试 36 个候选，每个候选按子标识符逐段比较：多字节子标识符按 BER 的 7 位分组累加还原成整数后比较数值，使 GetNext 按 OID 的数值顺序选择对象。

LEN1--LEN4 四拍自内向外计算出值、变量绑定、绑定表、PDU 和整条消息的长度，BUILD 再按“类型、长度、内容”逐字段写入回应 RAM；community、OID 和回显的值用 COPY 状态从请求 RAM 逐字节复制。回应写完后由发送状态逐字节读出，每字节 4 拍送入服务发送口。

代理一次处理一个请求，只在 IDLE 接收新请求；处理期间到达的请求仍被认领并计入请求数，但不回应，由管理端按超时重发。超过 256 字节的请求整条丢弃。

\subsection{配置访问与寄存器}

代理经与主机共享的配置读通路访问寄存器：主机访问其他地址时主机优先，同一地址的主机读和代理读在同一拍合并；代理写只在主机无读写且 \code{cfg_ready} 为 1 时授予。

SNMP 寄存器位于 \code{0x90}--\code{0xA3}，包括启用、回应 VL、community 长度和内容、请求计数、Trap 开关与计数及 sysDescr，见附录表~\ref{tab:regs}。

\section{615A 数据加卸载}\label{sec:a615}
\begingroup
\setlength{\parskip}{0pt}

\code{afdx_615a} 实现 ARINC 615A 的目标端：加载端经 TFTP 送来上传清单，目标主动读取头文件和数据文件，校验后把数据留在 8~KiB 加载区，并通过状态文件报告进度与结果；目标也能以介质定义卸载把加载区作为文件 \code{MEMORY.BIN} 读出。协议文件版本默认 \code{A4}，兼容 \code{A3}；目标标识为 \code{DEMC3_1}。

设计范围为单文件上传和介质定义卸载，信息操作与操作员定义卸载在初始化阶段拒绝；FIND、WAIT、ABORT、整文件重试和多文件加载不在实现范围内。TFTP 采用固定 512 字节块长，请求中的 RFC 2347 选项按忽略处理。ERROR 报文按目的 TID 关联会话，来源端口不参与校验。请求按操作码和零结尾字段解析，名称与模式字段的读取不以实际接收长度为终止条件。

\subsection{报文分派与异常处理}

\code{afdx_615a} 的主状态机约 70 个状态，按图~\ref{fig:a615fsm} 分组。收到的报文先整条存入 1~KiB 的接收缓冲，由分派状态 DISPATCH 按操作码、端口（TID）和块号分派给六类处理过程。图中的 Transfer/status 包含 List、Header、Data、Download 和 Status 处理；End 完成后返回等待状态。

各处理过程共用字节读取、整数拼接与报文输出状态序列。读字节以 \code{return_state} 保存返回位置，读整数另用 \code{num_return} 保存字段处理的返回位置；报文输出在发送 RAM 中逐字节组包，发送后按返回状态寄存器 \code{after_tx} 转移，通常返回等待报文状态 \code{WAIT_PACKET}。这些公共序列使各处理分支共用相同的同步读与发送时序。

数据或文件检查失败时，锁存错误码 \code{reason} 并转入 \code{FAILURE} 状态，依次发送 TFTP ERROR 和失败状态文件；状态文件传输失败（含重试耗尽）时结束。请求名称或模式错误只回 ERROR；LCI/LNO 在初始化阶段拒绝。

\figN[615A 目标状态机（按功能分组）]{f29_615a}{615A 目标端的请求分派、文件传输与状态报告状态机}{fig:a615fsm}

接收缓冲一次保存一个报文。当前报文未取走时，新到报文仍被认领，但不写入缓冲；对端未收到回应时，按超时机制重发。

\subsection{TFTP 传输与重传}

TFTP 使用 RRQ、WRQ、DATA、ACK、ERROR 五种报文，octet 模式，块长 512 字节，块号从 1 开始；WRQ 以 ACK0 建立传输。请求目的端口为 \code{REQUEST_PORT}（默认 59），建立传输后双方使用各自的传输标识符（TID），目标 TID 从 \code{0x8100} 递增。主文件与状态文件分别保存 TID 和块号。

报文在 RAM 中保留到回应到达；超时（默认 2~s）后原样重发上一报文，重传上限默认 1 次，耗尽以原因 5 失败。重复 DATA 只重发 ACK，不重复写 RAM。

\subsection{上传与文件校验}

上传流程见图~\ref{fig:a615}(a)：
\begin{enumerate}
  \item 加载端读取 \code{DEMC3_1.LUI}，目标返回接受状态 0001。
  \item 目标写 \code{DEMC3_1.LUS} 报告初始状态。
  \item 加载端写 \code{DEMC3_1.LUR} 上传清单，含头文件名和件号；版本不支持则以原因 1 拒绝。
  \item 目标发 LUS 0002 后读取清单中的 LUH 头文件；头文件校验通过后读取其列出的数据文件。
  \item 数据块按序写入加载区并累计 CRC。首个满块写入后暂缓 ACK，先经独立传输发送一次进度 LUS，再确认数据块。
  \item 末块（短块或空块）到达后，长度和 CRC 均正确才确认，并发送 LUS 0003、进度 100；否则发送 TFTP ERROR 和 LUS 1003。
\end{enumerate}

\figN[615A 交互]{f16_615a}{615A 单文件上传与介质定义卸载交互}{fig:a615}

两种传输的末块均少于 512 字节；文件长度为 512 字节的整倍数时，补发零长度 DATA 块。上传在 LUH 校验完成后先报告进度，再请求数据文件。

LUH 支持 ARINC 665-1 的 8002 版本：读取头文件字数、数据文件列表指针、数据文件数、文件名、数据文件字数和 CRC-16，检查头文件实际长度、版本、列表数量、头文件 CRC 及数据文件的长度和 CRC。件号和硬件 ID 不参与匹配；支持文件、用户数据、扩展点、Load CRC-32、解压和解密不参与当前加载处理。

CRC-16 的多项式为 \code{0x1021}，最高位优先、初值 \code{0xFFFF}、结果不取反，以 8 位并行的组合函数每拍处理一字节。

新 LUS 计数加 1，重传不加。进度取已写字节占预期总字节的整数百分比，成功前上限为 99。计算时先由移位相加完成乘 100，再以预期总字节数为减数逐拍迭代求商，最后通过减 10 迭代分离十位与个位，转换为十进制 ASCII 字符。

\subsection{介质定义卸载}

卸载流程见图~\ref{fig:a615}(b)。加载端读 \code{DEMC3_1.LND}，目标回应后发初始 LNS；加载端写 LNR 列表，只接受 \code{MEMORY.BIN}；目标发 LNS 0002 后主动写该文件，按 512 字节块停等发送最近一次成功上传的内容，首个满块确认后报告进度，末块确认后发 LNS 0003。卸载只读加载区，不改变已加载字节数。

上传和卸载互斥，操作在最终状态文件获 ACK 后才记为成功；新上传开始时撤销旧的有效长度。加载区采用 8~KiB 片上 RAM，供主机经配置总线按字节读回或用于介质定义卸载。


\subsection{存储与寄存器}

五块 RAM 分别存放接收报文（1~KiB）、发送与重传报文（1~KiB）、LUR（512~B）、LUH（2~KiB）和加载数据（8~KiB），均以 \code{afdx_ram_sdp} 例化。

寄存器 \code{0xB0}--\code{0xBB} 见表~\ref{tab:regs}。\code{STATUS} 报告空闲、进行中、成功、失败及原因码：0 无错误；1 元数据不支持或无效（版本、文件名、数量）；2 容量或长度超限；3 CRC 错；4 对端 TFTP ERROR 或模式错误；5 重传耗尽。加载区由 \code{MEMORY_ADDRESS} 选字节、\code{MEMORY_DATA} 读回。
\endgroup

\chapter{配置、复位与统计}\label{ch:cfg}

\code{afdx_cfg} 负责配置总线的地址译码、配置寄存器和 27 个统计计数器的维护、读数据多路选择，并在配置表写入时与收发核心执行许可握手。模块输出全部来自寄存器，结构见图~\ref{fig:cfg}。

\section{地址映射}

配置地址为 18 位字地址，映射见表~\ref{tab:addrmap}。各条 VL 通过自己的 \code{en} 位控制启用。对表外地址的写请求予以忽略，读请求返回 0。

\begin{longtable}{>{\raggedright\arraybackslash}p{4.2cm}>{\centering\arraybackslash}p{1.2cm}>{\raggedright\arraybackslash}p{\dimexpr\linewidth-5.4cm-6\tabcolsep\relax}}
\caption{配置地址映射}\label{tab:addrmap}\\
\longtablehead{3}{地址 & 读写 & 内容}
\code{0x00000} & 只读 & \code{VERSION}，\code{0xC3110000} \\
\code{0x00001} & 只读 & \code{STATUS}，位 0 为 \code{tx_ready} \\
\code{0x00002} & 读写 & \code{USER_ID[15:0]}，填入源 MAC 第 9--10 字节 \\
\code{0x00040}--\code{0x00046} & 只读 & 发送计数器：\code{accept}、\code{drop_dis}、\code{drop_full}、\code{drop_long}、\code{drop_nopage}、\code{tx_a}、\code{tx_b} \\
\code{0x00048}--\code{0x0004F} & 只读 & A 网接收计数器，对应 \code{ev_a} 位 0--7 \\
\code{0x00050}--\code{0x00057} & 只读 & B 网接收计数器，对应 \code{ev_b} 位 0--7 \\
\code{0x00058}、\code{0x00059} & 只读 & \code{rm_drop}、\code{deliver} \\
\code{0x0005A}、\code{0x0005B} & 只读 & \code{a_gdrop}、\code{b_gdrop}（保留，恒为 0） \\
\code{0x00080}--\code{0x000BB} & --- & RTC、SNMP、615A 寄存器（附录表~\ref{tab:regs}） \\
\code{0x01000} $+$ VL$\times$8 $+$ 字号 & 只写 & \code{tx_cfg} \\
\code{0x10000} $+$ VLID & 只写 & \code{rx_map}，数据取位 12:0 \\
\code{0x20000} $+$ idx & 只写 & \code{rx_cfg}，idx $<$ 4096 \\
\end{longtable}

三张表按地址高位译码：地址位 17:10 为 \code{0000_0100} 时选中 \code{tx_cfg}，位 17:16 为 \code{01} 时选中 \code{rx_map}，位 17:12 为 \code{10_0000} 时选中 \code{rx_cfg}。

\section{配置表格式}

三张配置表的字段见表~\ref{tab:cfgfmt}。复位后三张表全部清零，即全部 VL 未启用、全部 VLID 未映射。

\begin{longtable}{>{\raggedright\arraybackslash}p{2.6cm}>{\centering\arraybackslash}p{1.4cm}>{\raggedright\arraybackslash}p{2.6cm}>{\raggedright\arraybackslash}p{\dimexpr\linewidth-6.6cm-8\tabcolsep\relax}}
\caption{配置表字段}\label{tab:cfgfmt}\\
\longtablehead{4}{表与字 & 位 & 字段 & 含义}
\code{tx_cfg} 字 0 & 0 & \code{en} & VL 启用 \\
 & 1、2 & \code{net_a}、\code{net_b} & 从 A、B 网发送 \\
 & 8:3 & \code{qmax} & 队列最多排队的消息数，0--63 \\
 & 19:9 & \code{max_payload} & 单帧最大有效载荷 $L_\text{max}$，1--1471 字节（应用消息 $M>L_\text{max}$ 且 $L_\text{max}\ge 8$ 时触发 IPv4 分片，消息总长上限 $\MaxMessageBytes$~B） \\
\code{tx_cfg} 字 1 & 31:0 & \code{bag} & BAG，单位为拍，$1~\text{ms}=125000$，不超过 $2^{30}$ \\
\code{tx_cfg} 字 2 & 15:0 & VLID & 目的 MAC 低 16 位（顶层 \code{AFDX_top} 自动窥探同步至 \code{afdx_app_tx}） \\
\code{tx_cfg} 字 3、4 & 31:0 & 源 IP、目的 IP & IPv4 地址 \\
\code{tx_cfg} 字 5 & 31:16、15:0 & 源端口、目的端口 & 消息不带端口时使用的默认 UDP 端口 \\
\code{tx_cfg} 字 6、7 & --- & 保留 & 写 0 \\
\midrule
\code{rx_map}[VLID] & 12 & \code{valid} & 该 VLID 已映射（顶层 \code{AFDX_top} 自动窥探同步至 \code{afdx_app_rx}） \\
 & 11:0 & \code{idx} & 接收 VL 索引，0--4095 \\
\midrule
\code{rx_cfg}[idx] & 0 & \code{en} & VL 启用 \\
 & 1、2 & \code{net_a}、\code{net_b} & 接收 A、B 网的帧 \\
 & 3、4 & \code{ic_en}、\code{rm_en} & 启用完整性检查、冗余管理 \\
 & 16:5 & \code{skew_max} & 两网最大到达时间差，单位 32.768~µs \\
 & 27:17 & \code{max_payload} & 单帧最大载荷 $L_\text{max}$，1--1471 字节（后续 IPv4 分片允许 $\le L_\text{max}+8$） \\
\end{longtable}

\section{配置读写与仲裁}

读数据的多路选择分两拍完成，读延迟固定为 2 拍：第 1 拍按地址低 3 位从四组数据源（全局寄存器、发送计数器、A 网计数器、B 网计数器）及三项服务寄存器中并行预选出候选字并锁存，同时按地址位 17:3 锁存组号；第 2 拍按组号从各路寄存值中选出目标字送至 \code{cfg_rdata}，并拉高 \code{cfg_rvalid}，每拍至多经过一级 8 选 1 逻辑。

\figN[配置模块]{f30_cfg}{配置事务示例：(a) 暂存对 \code{rx_cfg[7]} 的写入，等待许可后写 RAM，并对比 RX/TX 的就绪重武装；(b) 读回 \code{USER_ID}=42 的两拍候选锁存与选择}{fig:cfg}

\code{cfg_ready} 为 1 时，总线在当拍接受 \code{cfg_wr} 写请求：普通寄存器和服务寄存器直接写入；配置表写请求则先装入暂存寄存器并拉低 \code{cfg_ready}，等待对应核心给出 \code{cfg_ok}（第~\ref{sec:intif}~节）。\code{tx_ready} 为 0 时，配置表写暂存寄存器保持清零，\code{cfg_ready} 为 0，不接受表初始化期间的写请求。

SNMP 代理读写管理对象使用同一条总线。顶层的仲裁规则是：主机在本拍写入，或主机与代理读取不同地址时，优先处理主机请求；主机读与代理读同一地址时合并为一次读，两者各自在 2 拍后取得数据；代理写只在主机无读写且 \code{cfg_ready} 为 1 时授予总线访问权。主机与代理的读数据分别由各自的 2 拍有效标记对齐。

\section{配置写与查表的冲突仲裁}\label{sec:cfgok}

配置表的读口归收发引擎，写口归配置总线。配置写通过 \code{cfg_ok} 与引擎读请求互斥，避免同步 RAM 同拍同址读写冲突。各引擎的许可条件为：
\begin{align}
\text{发送侧：}\ & \text{cfg\_ok} = \text{enable} \wedge (\text{op}=\text{IDLE}) \wedge \neg\text{want\_enq} \wedge \neg\text{want\_grant} \\
\text{接收侧：}\ & \begin{aligned}[t]
  \text{cfg\_ok} = \text{enable} &\wedge (\text{op}\in\{\text{IDLE},\text{AGE}\}) \\
  &\wedge \neg\text{desc\_a\_valid} \wedge \neg\text{desc\_b\_valid}
\end{aligned}
\end{align}
\code{cfg_ok} 由引擎自身的状态寄存器与同级模块的寄存输出（入队请求、帧描述有效、\code{grant_ready}）组合生成。\code{afdx_cfg} 接收一次写表请求后，保持写使能（\code{tx_cfg_we}、\code{rx_map_we} 或 \code{rx_cfg_we}）及地址、数据；当对应核心的 \code{cfg_ok} 为 1 时，引擎本拍不置位对应配置表的读使能寄存器，核心模块在该拍把写请求锁存到写口寄存器并于下一拍写入 RAM，\code{afdx_cfg} 撤销写请求。接收侧在许可后的下一拍恢复 \code{cfg_ready}；发送侧通过 \code{tx_cfg_ok_q} 与 \code{tx_cfg_rearm} 再延后一拍恢复就绪，保证本地应答与真实写请求完成对齐。时序为：
\begin{equation}
\begin{aligned}
&\text{cfg\_ready}(t+1)=0,\qquad
u=\min\{\,\tau \ge t+1 : \text{cfg\_ok}(\tau)=1\,\},\\
&\text{cfg\_ready}_{\mathrm{RX}}(u+1)=1,\qquad
\text{cfg\_ready}_{\mathrm{TX}}(u+2)=1.
\end{aligned}
\end{equation}
其中 $t$ 为总线接受写请求的时刻，$u-t$ 为写表等待时间，上界为 64 拍（第~\ref{sec:budget}~节）。


\section{复位与初始化}

复位与表初始化过程见第~\ref{sec:clkrst}~节。配置模块复位后，各寄存器恢复默认值，27 个计数器清零。\code{tx_ready} 为 0 期间，配置表写暂存寄存器保持清零，\code{cfg_ready} 保持为 0；此时不接受配置写请求，读访问照常进行。

\section{统计计数器}

27 个计数器均为 32 位，计满后回绕至 0，仅在复位时清零。为隔离发送核心与接收核心长距离跨模块引入的时序关键路径，\code{afdx_cfg} 对来自 TX/RX 通路的 5 路事件统计脉冲（\code{ev_tx}、\code{ev_a}、\code{ev_b}、\code{ev_rm_drop}、\code{ev_deliver}）插入了\textbf{两级流水寄存}（\code{_p} 与 \code{_q}）后再驱动 32 位加法器，降低了跨芯片远距离脉冲直接驱动大位宽累加器的布线延迟。

计入事件输入两级流水打拍（\code{_p} 与 \code{_q}）及读数据通路的两级流水结构，在 $t$ 拍发出配置总线读请求，$t+2$ 拍返回对应事件脉冲截至 $t-3$ 拍在核心源头产生的累计次数（对应配置模块内部计数值在 $t-1$ 拍及以前的累加状态）。计数器分为：
\begin{itemize}
  \item \textbf{发送（7 个）：}接纳的消息数；因 VL 未启用、队列满、超长、无空闲页而丢弃的消息数；A、B 网发出帧数。
  \item \textbf{接收，每网 8 个：}正确帧、FCS 错误、帧长错误、帧头错误、未知 VL、配置不允许、完整性检查丢弃、溢出。
  \item \textbf{接收，两网合计 2 个：}冗余管理丢弃、交付。
  \item \textbf{保留 2 个：}\code{a_gdrop}、\code{b_gdrop}，地址保留，恒为 0。
\end{itemize}

\section{异常处理汇总}\label{sec:errsum}

表~\ref{tab:errsum} 汇总全设计对各类异常的检测位置、处理方式和统计。发送侧拒收消息后，先回收已分配的页，再继续接收并丢弃剩余字节。

\begin{longtable}{>{\raggedright\arraybackslash}p{3.6cm}>{\raggedright\arraybackslash}p{2.0cm}>{\raggedright\arraybackslash}p{5.0cm}>{\raggedright\arraybackslash}p{2.8cm}}
\caption{异常处理汇总}\label{tab:errsum}\\
\longtablehead{4}{异常 & 检测位置 & 处理 & 统计}
\multicolumn{4}{l}{\textit{发送与顶层发送包装}} \\*
顶层发送 VLID 未映射 & \code{afdx_app_tx} & 整条消息丢弃，不进入核心 & \code{ev_app[0]} \\
顶层发送超 $\MaxMessageBytes$~B 或包装缓存满 & \code{afdx_app_tx} & 回滚包装写指针，丢弃本条消息 & \code{ev_app[1]} \\
VL 或两网均未启用 & \code{tx_engine} & 拒绝入队，回滚回收已用缓存页 & \code{drop_dis} \\
队列已满 & \code{tx_engine} & 同上 & \code{drop_full} \\
核心消息超过 $\MaxMessageBytes$~B，或分片时 $L_\text{max}<8$，或单帧超过 $L_\text{max}$（当配置不分片） & \code{tx_engine}、\code{tx_buf} & 从回滚 RAM 回收全部已分配页，继续接收并丢弃剩余字节 & \code{drop_long} \\
无空闲页 & \code{tx_buf} & 同上 & \code{drop_nopage} \\
MAC 发送下溢（TX Underrun） & \code{mac_tx} & 发送未取反的 FCS，使对端判为错误帧，继续接收并丢弃本帧剩余字节 & 不计 \code{tx_a}/\code{tx_b} \\
\midrule
\multicolumn{4}{l}{\textit{接收、分片重组与顶层接收包装}} \\*
前导码错误、帧长越界 & \code{mac_rx} & 整帧丢弃 & 帧长错误 \\
FCS 错误 & \code{mac_rx} & 整帧丢弃 & FCS 错误 \\
固定字段、长度、非 8B 对齐分片或校验和错误 & \code{rx_parse} & 回退预写入指针 & 帧头错误 \\
环形缓冲剩余空间不足 & \code{rx_parse} & 不写入本帧 & 溢出（\code{ev_a/b[7]}） \\
VLID 未映射 & \code{rx_engine} & 回退预写入指针 & 未知 VL \\
VL 或本网未启用、单帧超 $L_\text{max}$（后续片超 $L_\text{max}+8$） & \code{rx_engine} & 回退预写入指针 & 配置不允许 \\
SN 不连续 & \code{rx_engine} & 回退预写入指针，PSN 仍更新为本次 SN & 完整性检查丢弃 \\
另一网已交付的副本 & \code{rx_engine} & 回退预写入指针 & \code{rm_drop} \\
交付队列已满，或分片缺首片/乱序/重叠/16 槽全满 & \code{rx_engine}、\code{rx_reasm} & 回退预写入指针，拒绝进入重组槽位 & 溢出/拒绝（\code{ev_a/b[7]}） \\
末片总长不符、非零全包 UDP 校验错或 2~s 老化超时 & \code{rx_reassembly} & 废弃重组槽位并释放上下文 & 溢出/拒绝（\code{ev_a/b[7]}） \\
顶层接收包装缓存溢出 & \code{afdx_app_rx} & 整包丢弃，保护已入队消息边界 & \code{ev_app[2]} \\
表初始化完成前开始的帧 & \code{rx_parse} & 整帧忽略 & 不计 \\
\midrule
\multicolumn{4}{l}{\textit{服务与配置}} \\*
SNMP community 不符 & \code{afdx_snmp} & 不响应；Trap 启用时发 \code{authenticationFailure} & \code{BAD_COMMUNITY} \\
SNMP 请求超过 256 字节或代理忙 & \code{afdx_snmp} & 端口匹配后丢弃，由管理端超时重发 & 忙时计入 \code{REQUESTS} \\
615A 数据/文件检查错、状态传输失败、请求错误 & \code{afdx_615a} & 数据/文件错：ERROR 及失败状态文件；状态传输失败（含重试耗尽）：直接结束；名称/模式错：仅 ERROR；LCI/LNO：初始化拒绝 & \code{STATUS} 原因码 1--5（第~\ref{sec:a615}~节） \\
SNTP 响应版本、模式或 Originate 不符 & \code{afdx_rtc} & 丢弃，下一周期重发请求 & --- \\
表初始化期间写配置 & \code{afdx_cfg} & \code{cfg_ready}$=0$，不接受 & --- \\
未映射的地址 & \code{afdx_cfg} & 忽略写请求，读出 0 & --- \\
\end{longtable}

\section{主机编程流程}\label{sec:prog}

\subsection{初始化与启用}

主机经配置总线按下列顺序使端系统进入工作状态。每次写表都要等 \code{cfg_ready} 为 1（第~\ref{ch:if}~章）。
\begin{enumerate}
  \item \textbf{初始化：}释放复位后轮询 \code{STATUS}，等待位 0（\code{tx_ready}）变为 1，约 524~µs。
  \item \textbf{端系统标识：}写 \code{USER_ID}。
  \item \textbf{发送 VL：}对每条发送 VL 先写字 1--5（BAG、VLID、IP、默认端口），最后写字 0，同时设置 \code{en}、网络选择、\code{qmax} 和 \code{max_payload}。
  \item \textbf{接收 VL：}对每条接收 VL 先写 \code{rx_cfg}[idx]，再写 \code{rx_map}[VLID]$=\{1, \text{idx}\}$，使映射生效时配置已经就绪。
  \item \textbf{服务：}为每项服务配置一条发送 VL（目的 IP、端口与服务端口一致），写服务的 \code{TX_VL}、端口和 community 等寄存器，最后写 \code{ENABLE}。
  \item \textbf{开始收发：}应用开始提交消息。VL 较多时按第~\ref{sec:jitter}~节错开各 VL 首条消息的提交时刻。
\end{enumerate}

\subsection{运行中更新与服务配置}

\code{tx_cfg} 的 \code{en} 仅在入队时检查；已排队的消息继续参与授权，授权阶段逐拍读取字 0--5。修改一条发送 VL 时，主机应先停止向它提交消息，再把字 0 的 \code{en} 写 0，并保持该字其余字段不变；待该 VL 已提交的消息全部发出后，修改其余字，最后重新启用。

通过 \code{cfg_ok} 保证配置写与表读互斥，多字配置按字依次生效。接收侧的 \code{rx_cfg} 和 \code{rx_map} 各自通过单拍整字写入更新；两表的写入分别生效，历史状态在配置更新后保留。

应用与服务的配置还须满足以下条件：
\begin{itemize}
  \item 应用须等待表示收发核心均已就绪的 \code{tx_ready} 有效后，再提交消息。\code{app_tx_ready} 由应用仲裁选中条件与发送核心的 \code{m_ready} 相与生成。
  \item 各服务的 UDP 端口须互不冲突；RTC 的模式与时间，以及 SNMP 的 \code{ENABLE}、\mbox{community} 和 \code{TX_VL}，仅在对应服务没有在途消息时修改。
  \item 615A 忙时，忽略对启用、发送 VL、请求端口、超时、重传上限和文件版本的控制写请求；写 \code{0xB4} 后至少等待 4 拍，且只在非忙状态下读 \code{0xB5}。
\end{itemize}

\section{配置约束与网络规划规格}\label{sec:cfgrule}

主机须检查各项配置的组合是否满足带宽、抖动、缓存和分片重组并发度要求。结合官方 1030 答疑 R138-3/R139-3，本设计声明的\textbf{端系统满载网络规划规格}为单端 63 条发送 VL / 63 条接收 VL（$\text{BAG}=1$~ms，$L_\text{max}=1471$~B，双网冗余全双工，物理线路占用率 $77.5\%$，留有约 $22.5\%$ 突发排空与管理服务裕量）。具体约束及超限处理如下：
\begin{itemize}
  \item \textbf{带宽：}A、B 两网由同一个 MAC 状态机同拍驱动，只发 A 网或只发 B 网的帧同样占用一个发送时隙，所以约束针对全部发送 VL，不分网络：$\sum_v C(n_v)/\text{BAG}_v \le 1$，其中 $n_v$ 为单帧最大载荷 $L_\text{max}$，$C(n_v)$ 为第~\ref{sec:budget}~节式~(3-1) 的线路占用拍数，服务 VL 计入求和。长包分片消息在每个 BAG 周期仅调度 1 个分片帧，单周期带宽占用与单帧 VL 一致。
  \item \textbf{抖动：}BAG 取 $2^k$~ms；周期业务宜错开各 VL 首条消息相位，在 63~VL 满载规格下即使发生旋转三元组背靠背突发（B2），同 VL 最大端到端极差抖动仍保持在 \JitSixtyThreeBTwo~µs（$<100$~µs），而超规格 81~VL 对齐突发将因排队逼近 1~ms 周期而超出百微秒门限。
  \item \textbf{发送与重组缓存：}发送核心共 256 页（128~KiB），单帧 $1471$~B 占 3 页，$8192$~B 最大长包占 16 页；主机规划在途驻留页数须满足 $\sum_v \text{qmax}_v \times \lceil M_v / 512 \rceil + 9 \le 256$（例如 60 条短包加 $1472/3000/8192$~B 三条长包共占 205 页，安全满足；多条并发 $8192$~B 长包则须受 256 页总预算与顶层 32~KiB 包装缓冲约束）。接收侧配备 \ReasmContexts{} 个并发重组槽位（奇偶 VL 索引各 8 个槽位），同一时刻处于跨 BAG 分片重组中的活跃长包 VL 数在奇偶各 Bank 内均不得超过 8 条，且同 VL 新长包须在前一长包末分片发出后再提交。
  \item \textbf{交付：}交付口每拍一字节。若某条 VL 关闭冗余管理，且同时接收两网数据，则每条消息交付两次。所有接收 VL 的载荷交付总速率须低于 1 字节每拍；两网独立满速且都要交付时，所需交付速率超过该上限，超出部分按溢出计数。
  \item \textbf{顺序与冗余窗口：}跨 VL 的发送顺序由各 VL 的 BAG 与候选选择共同决定；接收分片必须在同一 VL 内按递增片偏移顺序到达，乱序分片将被拒绝。接收副本的到达间隔超过 \code{skew_max} 时，冗余记录按过期处理。
  \item \textbf{取值范围：}\code{bag} $\le 2^{30}$ 拍；\code{qmax} $\le 63$；\code{max_payload} 为 1--1471（启用分片时 $\ge 8$）；应用消息总长 $1\le M\le \MaxMessageBytes$~B；\code{skew_max} 小于过期清理门限 \code{0x4000} 个计时单位（12 位字段自然满足）。
\end{itemize}


\chapter{时序与存储设计}\label{ch:timing}

\section{时序设计规则}

为使 4096 条 VL 的查表通路和逐字节数据通路都能在 8~ns 内完成，全部 RTL 按表~\ref{tab:rules} 的十条规则组织寄存器、RAM 访问和组合运算。

\begin{table}\centering\small
\caption{时序设计规则}\label{tab:rules}
\begin{tabularx}{\linewidth}{>{\centering\arraybackslash}p{0.8cm}L}
\toprule
编号 & 规则 \\
\midrule
1 & 全设计使用单一时钟。 \\
2 & 逐 VL 数据只存放于 RAM，不得出现按 VL 数展开的触发器阵列或并行比较。 \\
3 & 所有 RAM 仅通过 \code{afdx_ram_sdp} 例化，地址、写数据和使能均直接来自寄存器。小寄存器组须注明用途（载荷 FIFO 8 项、归还 FIFO 4 项、接收描述 32 项）。 \\
4 & RAM 读出先经寄存器，再做加、减、比较；中间最多允许一级不超过 4 选 1 的选择。 \\
5 & 每拍最多一级 32 位算术；不使用乘法器和除法器，地址以位拼接生成。 \\
6 & 多于 8 选 1 的多路选择分两拍完成。 \\
7 & 数据逐拍连续推进；RAM 与算术采用寄存流水，帧头期间执行预读。 \\
8 & 同一 RAM 的同拍同址读写必须互斥。 \\
9 & 模块输出以寄存器为主；组合输出仅限 ready 握手、\code{cfg_ok}，以及 \code{afdx_svc} 的数据选择与接收分发直通，后者只做选择和门控。 \\
10 & 以单一 125~MHz、8~ns 为时钟约束，内部设计目标为综合 setup 裕量不小于 2.0~ns；不设置 false path 或多周期例外。 \\
\bottomrule
\end{tabularx}
\end{table}

规则 2 和 4 是“表加单主状态机”结构的基础：逐 VL 状态全部在 RAM 中，每张表只由一个状态机按固定流水访问，每拍的逻辑深度与 VL 数量无关。规则 5 使 RTC 的时间戳换算和偏差计算也拆成每拍一级加减，服务模块与收发核心均在寄存器之间限制组合运算的级数。

\section{存储预算}

全部存储都由 \code{afdx_ram_sdp} 例化，表~\ref{tab:ram} 按 RTL 数组声明的深度和位宽计算逻辑存储位数；器件存储块的映射用量见测试报告。

\begin{table}\centering\small
\caption{存储预算（逻辑位数）}\label{tab:ram}
\begin{tabularx}{\linewidth}{>{\raggedright\arraybackslash}p{3.6cm}>{\raggedleft\arraybackslash}p{2.8cm}>{\raggedleft\arraybackslash}p{2.2cm}L}
\toprule
表 & 深度$\times$位宽 & 位数 & 说明 \\
\midrule
\code{tx_payload} & 131072$\times$8 & 1048576 & 256 页 $\times$ 512 字节 \\
\code{tx_cfg} & 1024$\times$32 & 32768 & 128 条 VL $\times$ 8 字 \\
\code{tx_desc} & 256$\times$44 & 11264 & 每页一项 \\
\code{tx_state} & 128$\times$63 & 8064 & 每条 VL 一项 \\
\code{tx_mnext/pnext/free} & 3$\times$256$\times$8 & 6144 & 链表与空闲页 FIFO \\
\midrule
\code{rx_map} & 65536$\times$13 & 851968 & 任意 VLID 直接映射 \\
\code{rx_state} & 4096$\times$44 & 180224 & 每条 VL 一项 \\
\code{rx_cfg} & 4096$\times$32 & 131072 & 每条 VL 一字 \\
\code{rx_ring_a/b} & 2$\times$4096$\times$8 & 65536 & 每网 4~KiB \\
\midrule
SNMP 请求、回应 & 2$\times$256$\times$8 & 4096 & \\
615A 五块 & --- & 102400 & 收发报文各 1~KiB，LUR 512~B，LUH 2~KiB，加载区 8~KiB \\
\midrule
合计 & --- & 2442112 & 发送 1106816，接收 1228800，服务 106496 \\
\bottomrule
\end{tabularx}
\end{table}

\code{tx_payload} 存放发送载荷，\code{rx_map} 完成 16 位 VLID 到接收索引的映射。这两张深表均由块 RAM 实现。


\section{流水划分}

收发通路把查表、判决和数据传送分开：查表地址由寄存器驱动，RAM 读出先寄存，下一拍再做算术或比较；组帧与接收解析沿字节流逐拍推进。各模块的寄存边界见表~\ref{tab:pipeline-rule}。

\begin{table}\centering\small
\caption{主要通路的流水划分}\label{tab:pipeline-rule}
\begin{tabularx}{\linewidth}{>{\raggedright\arraybackslash}p{2.6cm}LL}
\toprule
通路 & 组合处理 & 寄存边界与推进方式 \\
\midrule
发送调度 & 到期比较、链表指针更新、候选选择 & 配置和状态读出先寄存；入队、授权和扫描按操作阶段推进（第~\ref{sec:txeng}~节） \\
组帧与 MAC & 帧头选择、校验和、CRC 更新 & 帧头期间预读载荷；每拍输出一个字节，双网共用发送状态机（第~\ref{sec:txbld}、\ref{sec:mactx}~节） \\
接收映射与判决 & VLID 查表、配置检查、序号与冗余窗口比较 & 映射表读出寄存后进入判决，各阶段逐拍完成，结果用于提交或回退（第~\ref{sec:rxtables}~节、表~\ref{tab:rxphase}） \\
服务协议处理 & 字节解析、字段比较、时间戳换算 & 同步读出字节后再解析；时间换算和偏差计算拆成逐拍加减（第~\ref{ch:svc}~章） \\
\bottomrule
\end{tabularx}
\end{table}

\code{cfg_ok} 将配置写与引擎访问错开，避免同拍读写同一表项；具体许可和写入时刻见第~\ref{sec:cfgok}~节。综合资源、建立与保持时间检查结果集中列于测试报告。

\chapter{可验证性设计}\label{ch:verif}

本章说明设计为验证提供的观察手段。验证平台、用例、核对项、测量口径和结果见测试报告。

\section{设计原则}

设计让每条消息和每一帧的去向都能只从顶层端口确定，从而满足三点：
\begin{itemize}
  \item \textbf{只看端口。}系统级验证不需要探测内部信号，同一套激励可以直接用于 RTL 和综合后网表。
  \item \textbf{期望值可独立生成。}帧格式、CRC、校验和、BAG 间隔和 IC/RM 判决都按协议规则定义，参考模型可以不读 RTL 独立算出期望结果。
  \item \textbf{逐条核对。}每条消息、每一帧恰好对应一个结果脉冲，验证可以逐条比对去向，而不只是比较总数。
\end{itemize}

\section{可观察性}

\begin{itemize}
  \item \textbf{事件端口。}\code{ev_tx}、\code{ev_a}、\code{ev_b}、\code{ev_rm_drop}、\code{ev_deliver} 对每条发送消息给出 accept 或四种丢弃之一，对线上每一帧在其所在网络给出一个类别脉冲，完好帧再给出一个结果脉冲（第~\ref{sec:rxclass}~节）。脉冲出现的时间有确定上界：发送消息的脉冲在末字节后 96 拍内，接收类别脉冲在帧尾后 32 拍内，结果脉冲在帧尾后 150 拍内。
  \item \textbf{计数器。}27 个统计计数器经配置口读回，读值的取样时刻有精确定义（$t$ 拍发出读请求，$t+2$ 拍返回 $t-1$ 拍及以前的脉冲数，第~\ref{ch:cfg}~章），可以与事件记录逐拍对齐。
  \item \textbf{线路字节。}GMII 输出直接来自 MAC 的输出寄存器，帧的每个字节和起始拍都在端口上可见；两网同拍起始，便于对照两网的输出。
  \item \textbf{服务状态。}RTC 的偏差和校时次数、SNMP 的请求和 community 错误计数、615A 的状态、原因码和加载区内容都可经配置口读回（附录~\ref{app:regs}）。
  \item \textbf{RAM 碰撞监视。}\code{afdx_ram_sdp} 在仿真中检查同拍读写同一地址，出现即报错；该监视用 \code{`ifndef SYNTHESIS} 包住，不影响综合网表。它直接检验第~\ref{sec:cfgok}~节“避让”设计是否成立。
\end{itemize}

\section{可控性}

链路上的各类异常都可以从端口构造：在 GMII 输入上改动字节、丢弃或重放整帧，即可产生 FCS 错、帧头错、序号乱序、双网时间差等情况；配置口可以随时改写 VL 配置以构造配置不允许、未知 VL 和队列满。设计不需要专门的测试模式。

\chapter{设计总结}\label{ch:limit}

\section{总体架构}

\code{AFDX_top} 采用单一 125~MHz 时钟，以发送核心、接收核心、服务通道、配置模块和时间基准组成端系统。发送和接收两侧都按每拍一字节组织数据通路，逐 VL 配置与状态存放在块 RAM 中，状态机按固定阶段完成读出、判决和写回。全设计通过 19 个 \code{afdx_ram_sdp} 实例承载载荷、配置、状态和协议文件，模块层次见表~\ref{tab:hierarchy}。

发送核心把消息存储与线路发送分开。分页缓存接纳应用消息，发送引擎按 VL 的 BAG 和队列状态调度，组帧器预读载荷并生成协议字段，双网 MAC 共用状态机和发送时刻。提前候选和帧头预备使允许时刻到来前完成控制准备；线路争用及其配置条件由第~\ref{sec:budget}~节的预算给出。

接收核心为 A、B 两网分别设置 MAC、解析器和环形缓冲，共享 VL 判决与应用交付。载荷先试写，经过帧格式、配置、完整性和冗余判决后提交，拒绝时回退。每网独立保存序号状态，冗余管理决定有效副本的交付，异常分类与统计事件记录各帧的处理结果。

\section{服务与主机控制}

RTC、SNMP 和 615A 通过服务通道共享收发核心，服务报文与应用报文按各自 VL 配置参与发送调度和接收检查。RTC 在报文握手处采样时间戳，分拍完成时间格式换算与偏差计算；SNMP 通过请求与回应 RAM 完成 BER 解析、对象访问和回应组装；615A 以 TFTP 传输协议文件及加载区数据，并通过状态文件报告进度与结果。各项协议的处理范围和字段检查见第~\ref{ch:svc}~章。

主机通过配置口管理 VL、服务参数和统计计数器。配置写在引擎许可的周期写入，配置读与计数器取样有固定时序；带宽、排队、缓存和交付约束由主机配置共同满足，编程步骤见第~\ref{ch:cfg}~章。顶层 GMII 字节流、事件端口和状态寄存器提供运行观察接口，其定义见第~\ref{ch:if}、\ref{ch:verif}~章。

按上述结构，端系统以 2\,442\,112 位逻辑存储和逐拍一级组合深度的数据通路达到设计指标：每网每方向载荷容量 956.44~Mbps；线路空闲时端到端时延上界 26.360~µs、发送抖动为 0；争用时抖动由式~(\ref{eq:jitter}) 给出，满足配置约束时不超过 50~µs；128 条发送 VL 与 4096 条接收 VL 可全部配置并同时启用。综合资源、时序检查、功能测试和性能测量结果集中列于测试报告。

\appendix
\chapter{源文件与默认参数}

\section{源文件}\label{app:files}

RTL 共 \RtlFileCount{} 个 Verilog 文件（合计 \countsep{\RtlLineCount} 行，\DesignUnitCount{} 个设计单元），均位于仓库 \code{rtl/} 目录。

\begin{longtable}{>{\raggedright\arraybackslash}p{4.6cm}>{\raggedright\arraybackslash}p{9.4cm}}
\caption{RTL 源文件清单（\RtlFileCount{} 个文件）}\label{tab:files}\\
\longtablehead{2}{文件 & 作用}
\code{AFDX_top.v} & 赛题模板顶层封装：2B VLID 流式包装器、配置窥探与 \code{afdx_es} 核心 \\
\code{afdx_es.v} & 端系统核心顶层：双网收发核心、服务通道、三项服务、配置与时间基准 \\
\code{afdx_app_tx.v} & 顶层发送流包装器：2B VLID 前缀解析、3 级 128 路比较映射与 32~KiB 缓存 \\
\code{afdx_app_rx.v} & 顶层接收流包装器：消息缓存、VLID 反查、2B VLID 前缀插入与溢出保护 \\
\code{afdx_app_map.v} & 接收包装层 $4096\times 16$~bit $\text{idx}\rightarrow\text{VLID}$ 专用同步映射 RAM \\
\code{afdx_app_payload.v} & 包装层 $8\times 4$~KiB 分体同步载荷存储（含 \code{_bank}/\code{_slice} 及写隔离） \\
\code{afdx_app_read.v} & 包装层 8 信用流式同步预取队列（屏蔽 4 拍块 RAM 读延迟） \\
\code{afdx_timebase.v} & 周期计数 \code{now}、粗粒度计时 \code{now_c} 与字节使能 \code{byte_en} \\
\code{afdx_cfg.v} & 18 位字地址译码、寄存器、27 个计数器两级流水维护、配置表写握手 \\
\code{afdx_svc.v} & 应用与服务的整消息轮询发送（授权寄存）、按 UDP 目的端口分发接收 \\
\code{afdx_tx_core.v} & 发送核心集成模块，例化发送各子模块、输出流缓冲及 \code{afdx_tx_payload_store} \\
\code{afdx_tx_buf.v} & $1\text{--}\MaxMessageBytes$~B 分页写入、两槽入口弹性级与 $16\times 8$ 回滚 RAM \\
\code{afdx_tx_free.v} & 256 页空闲页 FIFO，支持双页面回收端口并发回收 \\
\code{afdx_tx_engine.v} & BAG 与分片调度：扫描、入队、单帧/分片授权、分片进度表维护与起始确认 \\
\code{afdx_tx_builder.v} & 42B/34B 组帧、12 拍 IPv4 校验和准备、载荷预读与 $-8$ 拍提前补偿 \\
\code{afdx_mac_tx.v} & 双网发送 MAC：前导码、FCS、帧间隔、发送下溢（TX Underrun）检测与异常终止 \\
\code{afdx_crc32_d8.v} & 逐字节 8 位并行 CRC-32 组合逻辑 \\
\code{afdx_rx_core.v} & 接收核心集成模块，例化双 MAC、解析器、32-Bank 双写分片映射表、重组与交付 \\
\code{afdx_mac_rx.v} & 接收 MAC：SFD 检测、CRC 校验、剥离 FCS \\
\code{afdx_rx_parse.v} & 单帧/首片与后续分片帧头检查、载荷预写入环形缓冲、帧尾分类 \\
\code{afdx_rx_engine.v} & 30 拍日程：32-Bank 查表、配置检查、IC、RM、重组槽位查询预约与过期清理 \\
\code{afdx_rx_deliver.v} & 读取 $32\times 73$~bit \code{rx_dq}，从双网环形缓冲逐字节流式送出 \\
\code{afdx_rx_reassembly.v} & 双 Bank $\times 8$ 槽位（\ReasmContexts{} 上下文）顺序 IPv4 分片重组顶层 \\
\code{afdx_rx_reassembly_bank.v} & 单 Bank 8 槽位重组状态机、连续偏移检查、整包 UDP 校验与 2~s 超时回收 \\
\code{afdx_rx_reassembly_payload.v} & $16\times 8$~KiB 重组载荷存储池（64 条 $8192\times 2$ 通道与 128 个写隔离 LUT） \\
\code{afdx_rx_reassembly_ram.v} & 分片重组元数据与接收窄分片表复用的专用同步双端口 RAM 封装 \\
\code{afdx_ram_sdp.v} & 通用简单双端口（SDP）同步 RAM 原语封装（例化 \RamSdpCount{} 处） \\
\code{afdx_rtc.v} & 实时时钟与 SNTPv4 纳秒级时钟同步（含 40~ns 流水补偿） \\
\code{afdx_ntp_fraction.v} & 纳秒与 NTP 32 位二进制小数的 32 轮逐位换算 \\
\code{afdx_snmp.v} & SNMP 顶层包装：请求捕获、协议 RAM 接口与 \code{CE_DIV=2} 使能树 \\
\code{afdx_snmp_engine.v} & SNMP 协议核心：ASN.1/BER 编解码、MIB 读写与响应生成（16~ns 域） \\
\code{afdx_615a.v} & 615A 顶层包装：TFTP 捕获、6 块协议 RAM 接口与 \code{CE_DIV=2} 使能树 \\
\code{afdx_615a_engine.v} & 615A 协议核心：数据加载/卸载状态机与文件 CRC-16 硬件校验（16~ns 域） \\
\code{afdx_mgmt_tx.v} & 管理面流式预取发送适配器：8 字节预取队列，供 SNMP 与 615A 复用解耦 \\
\end{longtable}

\section{默认参数}\label{app:param}

\begin{longtable}{>{\raggedright\arraybackslash}p{4.0cm}>{\raggedleft\arraybackslash}p{2.5cm}>{\raggedright\arraybackslash}p{\dimexpr\linewidth-6.5cm-6\tabcolsep\relax}}
\caption{顶层与核心参数及默认容量}\label{tab:param}\\
\longtablehead{3}{参数 & 默认值 & 含义}
\code{TX_PAGE_AW} / \code{TX_OFF_AW} & 8 / 9 & 发送缓存 256 页 $\times 512$~B，共 128~KiB \\
\code{LOOKAHEAD} & 512 & 提前选择窗口，512 拍（4.096~µs） \\
\code{RX_RING_AW} / \code{RX_DQ_AW} & 12 / 5 & 每网接收环形缓冲 4~KiB，交付队列 32 项（73 位宽） \\
\code{RX_CONFIG_BANKS} & 16 & 接收配置总线与描述符局部复制扇出分支数 \\
\code{CE_DIV} & 2 & SNMP / 615A 协议引擎多周期时钟使能分频比（16~ns 域） \\
\code{REASSEMBLY_CONTEXTS} & \ReasmContexts & 接收 IPv4 顺序分片重组并发槽位数（2 Bank $\times 8$ 槽 $\times 8$~KiB） \\
\code{REASSEMBLY_TIMEOUT} & $2.5\times 10^8$ 拍 & 分片重组超时老化阈值（2.0~s） \\
\midrule
单条应用消息长度 & $1\text{--}\MaxMessageBytes$~B & $>1471$~B（或 $>\text{max\_payload}$）自动按 IPv4 分片传输 \\
顶层应用包装缓存 & $2\times 32$~KiB & \code{afdx_app_tx} 与 \code{afdx_app_rx} 各 32~KiB + 32 项描述符 \\
发送 / 接收 VL 容量 & 128 / 4096 & 接收映射表 65536 项；声明满载规划为 63 条并发收发 VL \\
帧间隔（IFG） & 12 拍 & 96~ns \\
上电表清零初始化 & 65,539 拍 & 核心 \code{init_done}（524.312~µs）；顶层 \code{tx_ready} 约 65,550 拍 \\
SNMP 请求缓存 & 256~B & 单 VarBind，支持 36 个 MIB 标量/表格节点 \\
615A 加载区 & 8~KiB & 头文件区 2~KiB，LUR 区 512~B，超时 2~s / 重传 1 次 \\
RTC 同步周期 & 125000 拍 & 1~ms，可经寄存器修改 \\
\end{longtable}

\chapter{服务寄存器}\label{app:regs}

基本寄存器、计数器和配置表地址见表~\ref{tab:addrmap}，服务寄存器见表~\ref{tab:regs}。RW 为可读写，RO 为只读。

\begingroup
\renewcommand{\tablestyle}{\small\linespread{1.04}\selectfont
  \renewcommand{\arraystretch}{1.10}}
\begin{longtable}{>{\raggedright\arraybackslash}p{2.4cm}>{\raggedright\arraybackslash}p{3.6cm}>{\centering\arraybackslash}p{0.9cm}>{\raggedright\arraybackslash}p{7.0cm}}
\caption{服务寄存器}\label{tab:regs}\\
\longtablehead{4}{地址 & 名称 & 读写 & 含义}
\multicolumn{4}{l}{\textit{RTC}} \\*
\code{0x80} & \code{CTRL} & RW & 位 0 启用，位 1 Server 模式；写位 2 装入暂存时间 \\
\code{0x81} & \code{PERIOD} & RW & Client 端请求间隔（拍），复位 125000 \\
\code{0x82} & \code{TX_VL} & RW & 发送 VL，低 7 位 \\
\code{0x83} & \code{UDP_PORT} & RW & UDP 端口，复位 123 \\
\code{0x84} & \code{SEC} & RW & 写秒暂存值；读当前秒并锁存纳秒 \\
\code{0x85} & \code{NS} & RW & 写纳秒暂存值；读锁存的纳秒值 \\
\code{0x86} & \code{OFFSET_SEC} & RO & 最近偏差的有符号秒部分 \\
\code{0x87} & \code{OFFSET_NS} & RO & 最近偏差的纳秒余数 \\
\code{0x88} & \code{SYNC_COUNT} & RO & Client 时钟同步次数 \\
\midrule
\multicolumn{4}{l}{\textit{SNMP}} \\*
\code{0x90} & \code{ENABLE} & RW & 位 0 启用代理 \\
\code{0x91} & \code{TX_VL} & RW & 响应与 Trap 的发送 VL \\
\code{0x92} & \code{COMMUNITY_LENGTH} & RW & 1--16，复位 6 \\
\code{0x93}--\code{0x96} & \code{COMMUNITY_0}--\code{3} & RW & community 内容，复位为 \code{public} \\
\code{0x97} & \code{SAMPLE_CYCLE} & RO & 上次 MIB 读取时的 \code{now} \\
\code{0x98} & \code{UPTIME} & RO & sysUpTime，单位 10~ms \\
\code{0x99} & \code{REQUESTS} & RO & 端口匹配的请求数 \\
\code{0x9A} & \code{BAD_COMMUNITY} & RO & community 不符次数 \\
\code{0x9B} & \code{TRAP_ENABLE} & RW & 位 0 启用 Trap \\
\code{0x9C} & \code{TRAPS_SENT} & RO & 已发送 Trap 数 \\
\code{0xA0}--\code{0xA3} & \code{DESCRIPTION_0}--\code{3} & RO & sysDescr 字符串 \\
\midrule
\multicolumn{4}{l}{\textit{615A}} \\*
\code{0xB0} & \code{ENABLE} & RW & 位 0 启用 \\
\code{0xB1} & \code{TX_VL} & RW & 发送 VL \\
\code{0xB2} & \code{STATUS} & RO & 位 1:0 空闲/进行中/成功/失败；位 15:8 原因码 \\
\code{0xB3} & \code{LOADED_BYTES} & RO & 本次已写字节数 \\
\code{0xB4} & \code{MEMORY_ADDRESS} & RW & 加载区读回地址，低 13 位 \\
\code{0xB5} & \code{MEMORY_DATA} & RO & 加载区读回字节 \\
\code{0xB6} & \code{LUS_COUNTER} & RO & 本次状态文件数 \\
\code{0xB7} & \code{RETRY_TOTAL} & RO & 累计超时重传次数 \\
\code{0xB8} & \code{REQUEST_PORT} & RW & TFTP 请求端口，复位 59 \\
\code{0xB9} & \code{TIMEOUT_CYCLES} & RW & 超时拍数，复位 250000000 \\
\code{0xBA} & \code{RETRY_LIMIT} & RW & 重传上限，复位 1 \\
\code{0xBB} & \code{PROTOCOL_VERSION} & RW & \code{A4} 或 \code{A3} \\
\end{longtable}
\endgroup

\chapter{内部接口信号}\label{app:intif}

本附录按信号说明第~\ref{sec:intif}~节汇总的内部接口。

\section{消息接口}

服务通道 \code{afdx_svc} 把选中的消息以字节流交给发送核心，见表~\ref{tab:ifm}。没有服务报文时，这组信号与应用口组合直通。

\begin{longtable}{>{\raggedright\arraybackslash}p{4.4cm}>{\raggedleft\arraybackslash}p{1.8cm}>{\raggedright\arraybackslash}p{\dimexpr\linewidth-6.2cm-6\tabcolsep\relax}}
\caption{消息接口（\texttt{afdx\_svc} $\rightarrow$ \texttt{afdx\_tx\_buf}）}\label{tab:ifm}\\
\longtablehead{3}{信号 & 位宽 & 含义}
\code{m_valid}、\code{m_ready} & 1、1 & 字节握手；\code{m_ready} 只在 \code{tx_buf} 处于 IDLE、RECV、DRAIN 时为 1 \\
\code{m_data}、\code{m_last} & 8、1 & 消息字节、末字节标记 \\
\code{m_vl} & 7 & 发送 VL 索引，首字节锁存 \\
\code{m_port_en}、\code{m_sport}、\code{m_dport} & 1、16、16 & \code{m_port_en} 为 1 时使用消息自带的 UDP 端口；首字节锁存 \\
\end{longtable}

\section{入队与授权}

\code{afdx_tx_buf} 写完消息后，将 47 位描述符写入 \code{tx_desc}，再向引擎发出入队请求。引擎选中一条消息（或其中一片 IPv4 分片）后，将单帧/分片字段一次性交给组帧器；组帧器在帧开始发送时报告起始时刻，并在分片载荷读毕时回传下一片游标，见表~\ref{tab:ifgrant}。

\begin{longtable}{>{\raggedright\arraybackslash}p{4.4cm}>{\raggedleft\arraybackslash}p{1.8cm}>{\raggedright\arraybackslash}p{\dimexpr\linewidth-6.2cm-6\tabcolsep\relax}}
\caption{入队、分片授权和发送起始指示接口}\label{tab:ifgrant}\\
\longtablehead{3}{信号 & 位宽 & 含义}
\multicolumn{3}{l}{\textit{入队（\code{tx_buf} $\rightarrow$ \code{tx_engine}，请求应答）}} \\*
\code{enq_valid} & 1 & 保持有效，直至收到 \code{enq_ack} \\
\code{enq_vl}、\code{enq_head}、\code{enq_len} & 7、8、14 & 发送 VL 索引、消息首页号、整包字节数（$1\text{--}\MaxMessageBytes$） \\
\code{enq_ack}、\code{enq_result} & 1、3 & 单拍应答；0 接纳，1 VL 未启用，2 队列满，3 超长（4 无空闲页由 \code{tx_buf} 自行判定） \\
\code{tx_desc}[46:0] & 47 & \{\code{len}[13:0], \code{port_en}, \code{sport}[15:0], \code{dport}[15:0]\}，以首页号为地址 \\
\midrule
\multicolumn{3}{l}{\textit{授权（\code{tx_engine} $\rightarrow$ \code{tx_builder}，单深缓冲握手）}} \\*
\code{grant_valid}、\code{grant_ready} & 1、1 & \code{grant_ready} $=\neg$\code{pending}，即单帧待发槽位空闲 \\*
\code{grant_a}、\code{grant_b} & 1、1 & 两网启用位（配置字 0 的 [1]、[2]） \\*
\code{grant_timed}、\code{grant_time} & 1、32 & \code{grant_timed} 为 1 时，帧不得早于 \code{grant_time} 起始；为 0 时可立即发送 \\*
\code{grant_vlid} & 16 & 配置字 2 \\*
\code{grant_src}、\code{grant_dst} & 32、32 & 配置字 3、4 \\*
\code{grant_sport}、\code{grant_dport} & 16、16 & \code{port_en} 为 1 取消息自带端口，否则取配置字 5 \\*
\code{grant_len}、\code{grant_sn}、\code{grant_page} & 11、8、8 & 本帧载荷字节数、AFDX 帧序号、本片起始页号 \\*
\code{grant_head}、\code{grant_offset} & 8、9 & 消息首页号（首片保留至末片释放）、本片页内起始字节偏移 \\*
\code{grant_ip_offset}、\code{grant_message_len} & 13、14 & IPv4 8B 单位片偏移、整条应用消息总字节数 \\*
\code{grant_ident}、\code{grant_first}、\code{grant_more} & 16、1、1 & IPv4 包标识、首片标志（含 8B UDP 头）、\code{MF} 更多分片标志 \\*
\midrule
\multicolumn{3}{l}{\textit{发送起始与分片游标回传（\code{tx_builder} $\rightarrow$ \code{tx_engine}）}} \\*
\code{start_valid}、\code{t_start} & 1、32 & 单拍脉冲；\code{t_start} 为输出首个前导字节时的 \code{now} \\
\code{fragment_done}、\code{_vl}、\code{_page}、\code{_offset} & 1、7、8、9 & 非末片载荷读毕时回传下一待发分片的起始页号与页内偏移 \\
\end{longtable}

\section{组帧器与 MAC}

组帧器将目的 MAC 地址至 SN 字段的帧内容以字节流送至 MAC，由 MAC 控制字节接收时刻，见表~\ref{tab:ifs}。

\begin{longtable}{>{\raggedright\arraybackslash}p{4.4cm}>{\raggedleft\arraybackslash}p{1.8cm}>{\raggedright\arraybackslash}p{\dimexpr\linewidth-6.2cm-6\tabcolsep\relax}}
\caption{帧字节接口（\texttt{tx\_builder} $\leftrightarrow$ \texttt{mac\_tx}）}\label{tab:ifs}\\
\longtablehead{3}{信号 & 位宽 & 含义}
\code{s_valid}、\code{s_ready} & 1、1 & 字节握手；\code{s_ready} 在 \code{byte_en}$=1$ 且满足以下任一条件时为 1：前导码末拍；DATA 且 \code{last_byte}$=0$；DRAIN \\
\code{s_data_a}、\code{s_data_b} & 8、8 & 两网字节，仅第 11 字节不同 \\
\code{s_last} & 1 & 与 SN 字节同拍 \\
\code{s_en_a}、\code{s_en_b} & 1、1 & 起始拍锁存，选择本帧发往哪一网 \\
\code{tx_sof}、\code{tx_eof} & 1、1 & 首个前导字节、帧尾（FCS 发完）的单拍脉冲 \\
\code{tx_underrun} & 1 & 数据阶段上游未备好字节的下溢（Underrun）指示 \\
\end{longtable}

\section{帧描述与交付队列}

每网解析器为正确帧生成一份 221 位帧描述（兼容不分片单帧与 IPv4 分片帧）。引擎读取描述并完成 30 拍判定与重组槽位查询后，向解析器发出提交或回退指示；对于通过判定的帧，引擎向 $32\times 73$~bit 交付队列写入一项，见表~\ref{tab:ifdesc}。

\begin{longtable}{>{\raggedright\arraybackslash}p{4.4cm}>{\raggedleft\arraybackslash}p{1.8cm}>{\raggedright\arraybackslash}p{\dimexpr\linewidth-6.2cm-6\tabcolsep\relax}}
\caption{帧描述与交付队列项}\label{tab:ifdesc}\\
\longtablehead{3}{字段 & 位宽 & 含义}
\multicolumn{3}{l}{\textit{帧描述（\code{rx_parse} $\rightarrow$ \code{rx_engine}，单深缓冲握手，共 221 位，高位在前）}} \\*
\code{ip_src}、\code{ip_dst}、\code{ip_id} & 32、32、16 & IPv4 源/目的地址与 16 位分片标识（用于分片重组五元组与伪头） \\*
\code{frag_offset}、\code{mf} & 13、1 & IPv4 8B 单位片偏移与更多分片标志（均为 0 表示不分片单帧） \\*
\code{udp_len}、\code{udp_check} & 16、16 & 首片/单帧的 UDP 总长与校验和字段 \\*
\code{vlid}、\code{sn} & 16、8 & 目的 MAC 第 4--5 字节 VLID 与帧末 AFDX SN \\*
\code{sport}、\code{dport} & 16、16 & 首片/单帧提取的 UDP 源端口与目的端口 \\*
\code{n}、\code{start}、\code{arrival} & 11、12、16 & 本帧载荷字节数、环形缓冲起点与帧尾时刻的 \code{now_c} \\*
\code{desc_valid}、\code{desc_take} & 1、1 & 交接握手（按配置扇出分支复制） \\*
\code{resolve_valid}、\code{resolve_commit} & 1、1 & 引擎判定：提交则 \code{committed} 前进到本帧末，否则 \code{trial} 回退 \\*
\midrule
\multicolumn{3}{l}{\textit{交付队列项（\code{rx_engine} $\rightarrow$ \code{rx_deliver}，32 项 $\times$ 73~bit 块 RAM \code{rx_dq}，高位在前）}} \\*
\code{frag[72]}、\code{slot[71:68]} & 1、4 & 分片标志（0 单帧直通，1 分片重组）与预约的重组槽位号（\texttt{\{bank,row\}}） \\
\code{net[67]} & 1 & 0 为 A 网，1 为 B 网 \\
\code{start[66:55]}、\code{n[54:44]} & 12、11 & 环形缓冲起点和本帧载荷字节数 \\
\code{idx[43:32]} & 12 & 接收 VL 索引 \\
\code{sport[31:16]}、\code{dport[15:0]} & 16、16 & UDP 源端口与目的端口（共 73 位） \\
\end{longtable}

核心交付口（重组模块输出 \code{d_*}）的字段与 \code{afdx_es} 应用接收口对应（目的端口在内部记为 \code{d_dport}，在核心记为 \code{app_rx_port}）。服务通道将这些信号同时广播给应用和三个服务。各服务在首字节到达当拍，根据 UDP 目的端口产生组合匹配信号 \code{svc_rx_claim}；匹配结果寄存并保持至末字节。端口匹配的消息由对应服务接收；其余应用消息在顶层 \code{AFDX_top} 中经 \code{afdx_app_rx} 加上 2 字节大端 VLID 前缀后自 \code{rx_data}/\code{rx_valid} 输出。

\section{服务发送总线}

三个服务的发送口打包成一组总线接到 \code{afdx_svc}，按 RTC、SNMP、615A 从低位到高位拼接，见表~\ref{tab:ifsvc}。各服务按相同的字节流握手约定驱动对应位段。

\begin{longtable}{>{\raggedright\arraybackslash}p{4.4cm}>{\raggedleft\arraybackslash}p{1.8cm}>{\raggedright\arraybackslash}p{\dimexpr\linewidth-6.2cm-6\tabcolsep\relax}}
\caption{服务发送总线 \texttt{svc\_tx\_*}}\label{tab:ifsvc}\\
\longtablehead{3}{信号 & 位宽 & 拼接}
\code{svc_tx_valid}、\code{_ready}、\code{_last}、\code{_port_en} & 各 $3\times1$ & [0] RTC，[1] SNMP，[2] 615A \\
\code{svc_tx_data} & $3\times8$ & [7:0] RTC，[15:8] SNMP，[23:16] 615A \\
\code{svc_tx_vl} & $3\times7$ & [6:0]、[13:7]、[20:14] \\
\code{svc_tx_sport}、\code{svc_tx_dport} & 各 $3\times16$ & [15:0]、[31:16]、[47:32] \\
\end{longtable}

\chapter{术语与参考文献}

\section{术语}\label{app:terms}

\begin{table}\centering\small
\caption{本文专用术语}\label{tab:terms}
\begin{tabularx}{\linewidth}{>{\raggedright\arraybackslash}p{3.4cm}L}
\toprule
术语 & 含义 \\
\midrule
拍 & 一个 125~MHz 时钟周期，8~ns \\
PSN & 同一网上同一 VL 上一帧的序列号 \\
预写入 / 提交 / 回退 & 载荷先写入环形缓冲；通过检查后更新提交写指针，拒绝时回退预写入指针 \\
候选 & 调度引擎选出、等待授权给组帧器的 VL \\
\code{pend_vl} & 已授权但尚未开始发送的 VL \\
\code{cfg_ok} & 引擎在本拍不发出配置表读的指示 \\
有效副本 & 帧格式、FCS 与协议头检查通过，且满足接收配置和已启用的完整性检查的单网帧副本 \\
页面释放 & 将页号送回空闲页池（经页面回收口 \code{ret1_*} 或 \code{ret2_*}）\\
端口匹配 & 服务根据 UDP 目的端口确定是否接收当前消息（信号 \code{svc_rx_claim}）\\
表初始化 & 复位后逐项清零收发配置表与状态表（\code{rx_map}、\code{rx_cfg}、\code{rx_state}、\code{tx_cfg}、\code{tx_state}）的过程 \\
线路占用 & 一帧在线路上的占用时长 $C(n)$，单位为拍 \\
Server / Client & SNTP 角色：Server 端应答请求并回显时间戳；Client 端周期发送请求并执行时钟同步校准 \\
应答 & 请求应答式握手的确认脉冲（\code{ack}），用于确认接口请求，与协议响应报文不同 \\
计时单位 & \code{now_c} 的单位，32.768~µs \\
发送抖动 & 帧起始相对 BAG 允许时刻的偏差 \\
起始指示 / 起始确认 & 组帧器在 \code{tx_sof} 有效拍经 \code{start_valid}/\code{t_start} 发送起始指示（TX Start Indication）；发送引擎据此完成发送起始确认（TX Start Confirm）并写回更新后的 \code{next_time} \\
\bottomrule
\end{tabularx}
\end{table}

\section{参考文献}\label{app:ref}

\begingroup
\small\linespread{1.08}\selectfont
\begin{enumerate}[label={[\arabic*]},leftmargin=2.5em]
  \item ARINC Specification 664 Part 7: Avionics Full Duplex Switched Ethernet Network, 2005-06-27.
  \item ARINC Report 615A-2: Software Data Loader Using Ethernet Interface, 2002-05-10.
  \item ARINC Report 665-1: Loadable Software Standards, 2001-01-12.
  \item RFC 1350: The TFTP Protocol (Revision 2), 1992.
  \item RFC 2347: TFTP Option Extension, 1998.
  \item RFC 1157: A Simple Network Management Protocol (SNMP), 1990.
  \item RFC 3416: Version 2 of the Protocol Operations for the Simple Network Management Protocol (SNMP), 2002.
  \item RFC 3418: Management Information Base (MIB) for the Simple Network Management Protocol (SNMP), 2002.
  \item RFC 4330: Simple Network Time Protocol (SNTP) Version 4 for IPv4, IPv6 and OSI, 2006.
  \item IEEE Std 802.3: IEEE Standard for Ethernet, 2018.
  \item 中科亿海微. 全国大学生嵌入式芯片与系统设计竞赛 FPGA 创新设计赛道选题指南.
\end{enumerate}
\endgroup


\end{document}
